% ch3_analysis.tex — 3장 분석
% Called by: main.tex line 35 (% ch3_analysis.tex — 3장 분석
% Called by: main.tex line 35 (% ch3_analysis.tex — 3장 분석
% Called by: main.tex line 35 (% ch3_analysis.tex — 3장 분석
% Called by: main.tex line 35 (\input{sections/ch3_analysis.tex})
\section{분석}
\label{sec:analysis}

\noindent{\sffamily\small\color{muted}\textbf{CHAPTER 3}}

\begin{quoteblock}
측정 $\rightarrow$ HLM $\rightarrow$ 베이지안 통합의 3단계 분석으로 서울 디지털 격차의 공간 분포, 개인-지역 기여 분리, 그리고 텍스트 기반 잠재 격차 신호까지 통합 추정한다.
\end{quoteblock}

\figfull{image6.png}{핵심 발견 요약}{fig:3-keyfindings}
\figfull{image7.png}{디지털 사막 --- 종합지수 하위 4개 자치구}{fig:3-desert}

본 장의 세 분석은 순차적으로 연결된다. 먼저 서울 자치구별 UMC 지수를 구축하여 지역 격차를 측정하고(\ref{subsec:3-1}), 이 결과를 지역 수준 변수로 활용하여 개인 및 지역 특성이 개인의 디지털 연결성에 미치는 영향을 분석한다(\ref{subsec:3-2}). 마지막으로 디지털 플랫폼 텍스트를 통해 양적 분석이 포착하지 못하는 구체적 사례를 도출한다. 이 흐름은 측정 $\rightarrow$ 인과 규명 $\rightarrow$ 현상 보완으로 구성되며, \ref{subsec:3-4}절에서 세 분석의 결과를 종합한다.

\subsection{서울의 UMC 지역 격차}
\label{subsec:3-1}

\subsubsection{측정 지표 및 데이터}
\label{subsubsec:3-1-1}

본 프로젝트에서 구성한 측정 지표는 UMC의 6개 차원에 대한 개념적 정의를 바탕으로 한다. 각 차원은 Connectivity, Available for Use, Affordability, Devices, Skills, Safety로 구분되며, 개별 측정 지표는 의미성과 보편성을 동시에 고려한 두 가지 핵심 구성 요소에 따라 설정하였다~\citep{ITU2023}. 본 분석에 활용된 서울의 UMC를 구성하는 6개 차원별 측정 지표는 표~\ref{tab:1-umc-dims}과 같다.

측정 지표 선정 과정에서는 개념적 정의와의 부합성, 서울의 디지털 기술 환경, 데이터 가용성을 종합적으로 검토하였다. 분석에는 데이터를 확보할 수 있는 가장 최근 시점인 2023년과 2024년, 총 2개 시점의 자료를 활용하였다. 2개 시점의 국가 행정 자료 등을 활용하였다. 단, Devices, Skills, Security 차원은 데이터 제공 시점의 한계로 인해 2023년 조사 자료를 기준으로 두 시점 모두 동일한 값을 적용하였다.

Connectivity의 경우, 한국전파누리에서 제공하는 이동통신 무선국 데이터와 서울시 행정데이터인 생활인구 데이터를 활용하여 측정 지표를 산출하였다. \citet{ITU2023}에 따르면, Connectivity는 네트워크 커버리지 내에 거주하는 인구를 기준으로 측정될 수 있다. 그러나 한국은 직주 분리가 뚜렷한 경우가 많아, 거주지역과 실질적으로 디지털 환경에 접속하는 지역이 상이할 수 있다. 이에 본 연구에서는 단순 거주인구 대신 주간 유동인구 데이터를 활용하여 4G 및 5G 네트워크 커버리지를 산출함으로써, 한국의 맥락에 보다 적합한 지표를 고안하였다. 구체적으로, 4G 기지국 수와 5G 기지국 수를 각각 주간 생활인구 천 명당 기지국 수로 환산하여 집계구별로 산출한 뒤, 이를 자치구별로 평균한 값을 최종 지표로 사용하였다.

Available for Use는 디지털 서비스의 실질적 이용 가능성을 측정하는 차원이다. SKT 통신정보에서 제공하는 월별 지역 내 가구당 모바일 데이터 사용량과 금융, 쇼핑, 동영상, 배달 4개 범주의 온라인 서비스 이용일수를 인구가중 평균으로 자치구별 월별로 집계하고, 12개월 평균값을 사용하였다. 여기에 자치구 지방정부의 공공WiFi 등록 현황 데이터를 활용하여 전체 인구 1,000명당 공공WiFi AP 수를 산출하여 포함하였다. 공공WiFi 밀도는 가정 내 인터넷 접속 환경이 제한적인 고령층에게 특히 중요한 디지털 접근 경로로 기능한다~\citep{vanDijk2020}.

Affordability는 ITU ICT Development Index 2023의 두 가지 축 구조를 따라 구성하였다~\citep{ITU2024}. 구매력(capacity) 측면에서는 서울시 상권분석서비스에서 제공하는 자치구별 월평균 가구소득을 활용하고, 가격부담(burden) 측면에서는 SKT 통신정보의 최근 3개월 이내 요금 연체 비율을 활용하였다. 연체율은 행정동, 성별, 연령대별 인구가중 평균으로 산출한 뒤 자치구 단위로 집계하였으며, 값이 높을수록 경제적 접근성이 낮음을 의미하는 역방향 지표이다.

Devices는 디지털 기기 보유 및 이용의 다양성을 측정한다. 서울시민 디지털역량실태조사의 가구 내 핵심 디지털 기기 보유 유형 수와 개인이 실제로 사용하는 기기 유형 수를 자치구별 가중평균으로 산출하였다. \citet{ITU2023}의 보고서에 따르면, 휴대폰과 컴퓨터 등 기기의 특성에 따라서 기기의 보유와 이용은 서로 다르게 해석된다. 이를 한국 국내 맥락에선, 컴퓨터의 보유 여부도 디지털 연결성을 드러내는 지표가 될 수 있다. 스마트폰 보급률이 사실상 포화 상태인 서울에서는 데스크톱, 노트북, 태블릿 등 보조 기기의 보유 여부가 디지털 활용 역량의 차별화 요인으로 작용한다~\citep{vanDeursenHelsper2015}.

Digital Skills는 기초, 응용, 고급, 사이버활동의 4개 하위 영역으로 구성된다. 기초 역량은 스마트 기기 기본 조작 능력을, 응용 역량은 디지털 서비스 활용 능력을, 고급 역량은 스마트 기기 활동 숙련도를, 사이버활동 역량은 사이버 공간에서의 활동 능력을 각각 측정한다. 각 하위 영역의 원점수를 만점으로 나누어 0\,$\sim$\,1 비율로 변환한 뒤 자치구별 가중평균을 산출하였다.

Safety는 디지털 보안 행동과 보안 인식의 2개 하위 지표로 구성된다. 보안 행동은 비밀번호 관리, 소프트웨어 업데이트, 2단계 인증 사용 등 실제 보안 실천 행위를 측정하고, 보안 인식은 개인정보 보호의 중요성에 대한 인지 수준을 측정한다. 각 하위 지표를 만점 기준 비율로 변환하여 자치구별 가중평균을 산출하였다.

\begin{table*}[htbp]
  \centering\small
  \caption{UMC 차원별 측정 지표}
  \label{tab:1-umc-dims}
  \begin{tabularx}{\textwidth}{l Y Y Y}
    \toprule
    \textbf{Dimension} & \textbf{Indicators} & \textbf{Data Source} & \textbf{Note} \\
    \midrule
    Connectivity & station\_density\_4g, station\_density\_5g, download\_speed & Spectrum Resource Mgmt System; NIA Communication Quality & Per 1K daytime living pop; CQ uses 2022--2024 avg \\
    \addlinespace[2pt]
    Available for Use & mobile\_data\_usage, online\_service\_days, wifi\_density & SKT Telecom Data; Local Admin WiFi License Data & Pop-weighted district avg; WiFi per 1K total pop \\
    \addlinespace[2pt]
    Affordability & avg\_income, delinq\_rate (reverse) & Seoul Commercial Analysis Service; SKT Telecom Data & ITU IDI capacity + burden; delinq = reverse-coded \\
    \addlinespace[2pt]
    Devices & device\_home\_core, device\_use\_core & Seoul Digital Competency Survey 2023 & Home: excl. smartphone; Use: incl. smartphone \\
    \addlinespace[2pt]
    Digital Skills & Basic / Applied / Advanced / Cyber activity & Seoul Digital Competency Survey 2023 & --- \\
    \addlinespace[2pt]
    Safety & safety\_behavior, safety\_awareness & Seoul Digital Competency Survey 2023 & Security practice (7 items) + Awareness (3 items) \\
    \bottomrule
  \end{tabularx}
\end{table*}

\subsubsection{측정 지표 지수화}
\label{subsubsec:3-1-2}

모든 원시 지표는 25개 자치구 내에서 Min-Max 스케일링을 통해 $[0,1]$ 구간으로 정규화한다. 정규화된 값 $x^{*}$는 식~\eqref{eq:minmax}과 같이 산출한다.

\begin{equation}\label{eq:minmax}
  x^{*} = \frac{x_{i} - x_{\min}}{x_{\max} - x_{\min}}
\end{equation}

이때, 값이 클수록 나쁜 결과를 의미하는 역방향 지표는 변환하여 모든 지표의 방향성을 통일하였다. 그 다음, 각 차원 내에서 정규화된 지표들을 동일 가중치 산술 평균으로 결합하여 차원별 점수를 산출하였다. 동일 가중치로 분석을 진행한 이유는, 6개 차원 간 상대적 중요도에 대한 합의된 기준이 부재한 상태에서 연구자의 자의적 판단을 최소화하기 위함이다~\citep{ITU2024}. 산출된 종합지수는 0(최저)에서 1(최고) 범위를 가지며, 서울 25개 자치구 전체에 대해 2023년과 2024년 각각 독립적으로 산출한다. 해석에 있어 중요한 점은 이 종합 지수가 지역의 절대적 UMC가 아닌, 서울 내 다른 자치구 간의 상대적 위치를 보여준다는 것이다.

\subsubsection{지역별 UMC 지수 특성}
\label{subsubsec:3-1-3}

25개 자치구의 2024년 UMC 종합지수는 0.279(중랑구)에서 0.695(서초구)까지 분포하였으며, 최상위와 최하위 간 약 2.5배의 격차가 존재하였다. 서울이 높은 인터넷 보급률을 가진 도시라는 일반적 인식에도 불구하고, 하위 행정단위에서는 상당한 상대적 수준의 디지털 연결성 격차가 관찰되었다. 상위 5개 자치구는 서초구(0.695), 영등포구(0.686), 마포구(0.665), 서대문구(0.615), 용산구(0.603)이다. 하위 5개 자치구는 중랑구(0.279), 도봉구(0.303), 강북구(0.380), 구로구(0.389), 노원구(0.389)이다. 상위 자치구는 대체로 도심부 또는 강남권에 위치하여 밀집된 상업 활동, 높은 소득 수준, 풍부한 민간 디지털 인프라의 혜택을 받는 반면, 하위 자치구는 북부 외곽 및 남서부에 집중되어 있었다.

종합지수의 공간적 분포를 살펴보면(그림~\ref{fig:3-umc-comp}), UMC 지수가 높게 나온 자치구는 두 구역에 집중되어 있었다. 이는 각각 서초-영등포-마포로 이어지는 한강과 도심권 지역과, 용산-종로-중구 중심의 전통적 도심부이다. 반면, 북부 외곽(강북구, 도봉구, 노원구)과 남서부(구로구, 금천구)는 도시 평균 이하의 점수를 보였다. 이러한 공간 분포는 서울의 사회경제적인 자원 분포와 연관되어 있었다. 즉, 디지털 격차가 도시의 구조적 공간 불평등과 밀접히 연관되어 있을 가능성을 시사하였다.

차원별 분석에서 한편으로, 어떤 자치구도 6개 차원 모두에서 균일하게 우수하지는 않다는 점이 발견되었다(그림~\ref{fig:3-umc-heatmap}). 예컨대 서울 북부에 위치한 은평구의 경우, Connectivity 차원에서 0.781로 상위권에 위치하지만, Affordability 차원에서는 0.328로 하위권에 그쳤다. 유사하게, 관악구는 Skills 차원에서 0.795로 1위를 기록하였지만 Affordability에서는 0.152로 최하위권이다. 이러한 다차원적 지역간 불균등은, 종합 순위만으로는 정책적으로 의미 있는 차원별 취약성을 포착할 수 없음을 의미하였다.

\begin{keymsg}[title=핵심 발견 \textperiodcentered{} 디지털 사막 4개 자치구]
종합지수 하위 자치구(중랑\textperiodcentered{}도봉\textperiodcentered{}강북\textperiodcentered{}노원)는 Devices\textperiodcentered{}Safety 차원에서 동반 결핍을 보이며, 공간적으로도 서울 북부에 집중되어 있다. 이는 권역 단위 협력 거버넌스 구축의 정당화 근거를 제공한다.
\end{keymsg}

결론적으로, 그림~\ref{fig:3-umc-dim}에 제시된 바와 같이 공간적으로 밀집한 종합지수 하위 4개 자치구(중랑구, 도봉구, 강북구, 노원구)를 본 프로젝트에서는 잠정적으로 디지털 사막(Digital Desert)으로 분류한다. 이들 자치구는 공통적으로 Devices와 Safety 차원에서 뚜렷한 결핍을 보이며, Affordability 역시 서울시 평균에 미치지 못했다. 한편, Connectivity 차원에서는 일부 자치구(중랑구 0.630, 강북구 0.562)가 서울시 평균을 상회하여, 물리적 네트워크 인프라 자체의 부재보다는 이를 활용하기 위한 인적\textperiodcentered{}경제적 역량의 부족이 핵심 격차 요인일 가능성을 보여준다. 이러한 발견은 \ref{subsec:3-2}절의 다층 분석에서 개인 수준 변수와 자치구 수준 변수의 상대적 기여를 분리 및 검증함으로써 보다 체계적으로 검토된다.

또한, UMC의 일부 차원에서 지리적으로 인접한 지역이 서로 유사한 패턴을 보였다. 그림~\ref{fig:3-umc-dim}은 차원별 지역의 공간적 자기상관성을 나타낸 것이다. 공간적으로 군집해 있는 자치구가 서로 유사한 패턴을 보이는 점은 디지털 사막 해소를 위한 광역적 거버넌스 체계 구축의 가능성을 시사한다.

결론적으로, 그림~\ref{fig:3-umc-dim}에 제시된 바와 같이 종합지수 하위 5개 자치구(중랑구, 도봉구, 강북구, 구로구, 노원구)를 본 프로젝트에서는 잠정적으로 디지털 사막(Digital Desert)으로 분류한다. 이들 자치구는 공통적으로 Devices와 Safety 차원에서 뚜렷한 결핍을 보이며, Affordability 역시 서울시 평균에 미치지 못했다. 한편, Connectivity 차원에서는 일부 자치구(중랑구 0.630, 강북구 0.562)가 서울시 평균을 상회하여, 물리적 네트워크 인프라 자체의 부재보다는 이를 활용하기 위한 인적\textperiodcentered{}경제적 역량의 부족이 핵심 격차 요인일 가능성을 보여준다. 이러한 발견은 \ref{subsec:3-2}절의 다층 분석에서 개인 수준 변수와 자치구 수준 변수의 상대적 기여를 분리 및 검증함으로써 보다 체계적으로 검토된다. 또한, UMC의 일부 차원에서 지리적으로 인접한 지역이 서로 유사한 패턴을 보였다. 그림~\ref{fig:3-lisa-grid}은 차원별 지역의 공간적 자기상관성을 나타낸 것이다. 공간적으로 군집해 있는 자치구가 서로 유사한 패턴을 보이는 점은 디지털 사막 해소를 위한 광역적 거버넌스 체계 구축의 가능성을 시사한다.

\begin{figure}[htbp]
  \centering
  \begin{subfigure}[t]{0.55\linewidth}
    \centering
    \includegraphics[width=\linewidth]{image9.png}
    \caption{자치구별 UMC 종합점수}
    \label{fig:3-umc-comp}
  \end{subfigure}\hfill
  \begin{subfigure}[t]{0.43\linewidth}
    \centering
    \includegraphics[width=\linewidth]{image10.png}
    \caption{UMC 차원별 점수 히트맵}
    \label{fig:3-umc-heatmap}
  \end{subfigure}
  \caption{자치구별 UMC 종합점수 및 차원별 히트맵}
  \label{fig:3-umc-two}
\end{figure}

\figfull{image11.png}{UMC 차원별 점수}{fig:3-umc-dim}

\begin{figure*}[htbp]
  \centering
  \begin{subfigure}[t]{0.32\textwidth}
    \centering
    \includegraphics[width=\linewidth]{image12.jpeg}
    \caption{Connectivity}
  \end{subfigure}\hfill
  \begin{subfigure}[t]{0.32\textwidth}
    \centering
    \includegraphics[width=\linewidth]{image13.jpeg}
    \caption{Available for Use}
  \end{subfigure}\hfill
  \begin{subfigure}[t]{0.32\textwidth}
    \centering
    \includegraphics[width=\linewidth]{image14.jpeg}
    \caption{Affordability}
  \end{subfigure}

  \vspace{4pt}
  \begin{subfigure}[t]{0.32\textwidth}
    \centering
    \includegraphics[width=\linewidth]{image15.jpeg}
    \caption{Devices}
  \end{subfigure}\hfill
  \begin{subfigure}[t]{0.32\textwidth}
    \centering
    \includegraphics[width=\linewidth]{image16.jpeg}
    \caption{Skills}
  \end{subfigure}\hfill
  \begin{subfigure}[t]{0.32\textwidth}
    \centering
    \includegraphics[width=\linewidth]{image17.jpeg}
    \caption{Safety}
  \end{subfigure}
  \caption{UMC 수준의 공간적 자기상관 (서울 25개 자치구)}
  \label{fig:3-lisa-grid}
\end{figure*}

\subsection{개인 및 지역 특성과 디지털 연결성}
\label{subsec:3-2}

본 연구의 데이터는 개인(Level 1)이 자치구(Level 2) 내에 군집되어 있는 위계적 구조를 갖는다. 여기서 Level 2는 개인의 거주지에 해당한다. 이러한 구조에서 통상적인 OLS 회귀분석은 관측치의 비독립성 가정을 위반하여, 자치구 수준 예측변수의 표준오차를 과소추정하고 1종 오류를 팽창시킬 수 있다~\citep{RaudenbushBryk2002}. 따라서 분산을 25개 자치구를 기준으로 구 내와 구 간 성분으로 분해하고, 개인 수준과 자치구 수준 변수를 동시에 투입할 수 있는 2수준 위계적 선형 모형(HLM)을 채택하였다.

두 번째 분석에서 활용한 자료인 서울서베이는 2003년부터 서울시에서 2만 개 가구를 대상으로 년마다 실시해온 행정 자료이다. 자료의 수집은 매년 8\,$\sim$\,9월에 비대면과 방문면접을 병행하여 이뤄진다.

\subsubsection{변수 구성}
\label{subsubsec:3-2-1}

먼저 분석의 종속변수인 디지털 활용 점수는 서울서베이의 응답항목 중 디지털 서비스 이용빈도 항목의 8개 항목을 합산한 점수이다. 8개 항목은 생활정보 서비스, 전자상거래 서비스, 금융거래 서비스, 공공 서비스, 커뮤니케이션 서비스, 여가 서비스, 교육학습 서비스, 모바일 결제로 구성된다. 이들은 전화와 문자와 같은 기초적 연결 행위가 아니라, 디지털 기기를 매개로 한 실질적인 추가 활용에 해당하여, 디지털 환경이 발달한 서울에 적합한 지표로 판단하였다. 각 항목은 원래 4점 서열 척도를 0--10 연속 척도로 변환하여 총 0--80점 범위의 연속변수를 구성한다. 이 점수가 높을수록 개인의 디지털 서비스 활용이 다양하고 원하는 서비스를 자유롭게 이용할 수 있음을 의미한다.

본 분석의 독립변수는 개인 특성의 변수와 지역 특성의 변수로 구분된다(표~\ref{tab:2-hlm-vars}). 개인 특성의 변수는 개인의 디지털 활용에 영향을 미칠 수 있는 변수인 성별, 연령, 직업, 학력, 소득수준, 장애여부로 구성된다. 교육수준은 중졸 이하, 고졸, 대졸(기준범주), 대학원 이상의 4범주로, 직업은 관리\textperiodcentered{}전문직, 사무직(기준범주), 생산\textperiodcentered{}기능직, 기타(학생\textperiodcentered{}주부\textperiodcentered{}무직)의 4범주로 구분하였다. 주거유형은 아파트(기준범주), 단독주택, 다세대 및 다가구, 기타의 4범주로 분류하였다. 연령과 소득은 총평균 중심화를 적용하여 절편의 해석 가능성을 높였다. 추가로, 1인가구 더미, 장애 여부 더미, 연도 더미(기준: 2023년)를 포함하였다.

지역 수준(Level 2) 변수는 \ref{subsec:3-1}절에서 구축한 UMC 차원 점수 중 지역의 디지털 환경을 결정하는 외생적 요인의 세 차원 Connectivity, Available for Use, Affordability을 활용하였다. 이는, Devices, Digital Skills, Safety가 서울디지털역량실태조사의 개인 응답을 자치구 단위로 집계한 것을 고려한 것이다. 그로 인한 내생성 문제를 피하기 위해 주 분석에서 세 가지 차원을 제외하여 지역 차원의 보편적인 의미있는 연결성에 더 적합하다고 판단되는 차원을 선택하였다. 통제변수로 자치구 평균소득과 고령화율을 포함하였으며, 모든 Level-2 변수는 앞선 \ref{subsec:3-1}절에서와 같이 Min-Max 정규화를 적용하여 $[0,1]$ 범위로 표준화하였다.

분석은 4단계로 진행된다. 무조건 모형에서 급내상관계수를 산출하여 다층 분석의 정당성을 확인한 뒤, 개인 수준 변수를 투입한 M1, 자치구 수준 변수를 추가한 M2, 교차수준 상호작용을 포함한 M3를 순차적으로 추정한다. 교차수준 상호작용은 UMC 인프라 변수(Connectivity, Available for Use, Affordability)와 취약계층 지표의 곱으로 구성하여, 자치구의 디지털 인프라 수준이 취약계층의 디지털 활용에 차별적으로 작용하는지를 검증한다.

\subsubsection{분석 결과}
\label{subsubsec:3-2-2}

분석에 활용된 표본은 서울서베이 2023\,$\sim$\,2024년 응답자 총 10,000명이며, 25개 자치구에서 평균적으로 400명이 표집되었다. 주요 변수에 대한 기술통계량을 나타낸 것은 표~\ref{tab:3-descstat}과 같다. 개인의 디지털 연결성의 평균이 80점 만점에 60점으로 나타났는데, 서울시민의 높은 디지털 연결성을 잘 보여준다. 종속변수가 음의 왜도가 존재할 수 있으나 기관 표준 척도와 10000명에 근접한 표본을 고려해 척도 선형화한 원점수대로 분석을 진행하였다.

표~\ref{tab:4-hlm-results}은 순차적으로 다층 모형의 추정 결과를 나타낸 것이다. 무조건 모형의 급내상관계수는 0.49\%로 자치구 간 분산 비율 자체는 크지 않았으나, 자치구당 평균 표본 크기 400명을 감안한 설계효과가 2.94로 임계값 2.0을 초과하여 다수준 모형의 사용에 대한 타당성을 확인하였다~\citep{HoxEtAl2018}. 그러나, 25개의 자치구의 소규모 표본 / 자치구가 설명하는 분산이 매우 적은 점을 종합하면 개인의 디지털 연결성에 지역의 연결 관련 특성이 미치는 영향이 극히 적을 것으로 추정된다. 이러한 결과는 서울의 우수한 디지털 연결성과 높은 수준의 디지털 역량으로 지역 수준의 연결성 인프라가 영향을 미치지 못한 것일 수 있다. 또한, 서울서베이는 개인의 거주 지역을 조사하기 때문에, 실제로 연결성에 직접적으로 영향을 받는 지역과 일치하지 않을 수 있다.

개인 특성 중 교육 수준의 관계가 가장 두드러졌다. 중졸 이하 학력은 대졸 대비 디지털 활용 점수를 약 22.6점 낮게 나타났다. 이는 80점 만점 기준 28\%에 해당하는 상당한 격차다. 고졸 역시 대졸보다 7.3점 낮게 나타난 반면, 대학원 이상은 오히려 2.9점 높아 교육 수준과 디지털 활용 능력 간의 뚜렷한 정적 관계를 확인할 수 있다. 연령 또한 유의한 영향을 미쳤다. 평균 연령으로부터 1세 증가할수록 개인의 디지털 연결성에 관한 점수가 0.35점 감소하였으며, 65세 이상 고령층의 경우 연령 효과와는 별개로 추가적인 6.0점 하락이 확인되었다. 65세 이상의 연령층에서 디지털 연결성의 두드러진 저하가 나타났다. 직업 측면에서는 관리직\textperiodcentered{}전문직 종사자가 사무직 대비 2.3점 높게 나타나, 고소득 직종일수록 디지털 활용 수준도 높은 경향을 보였다. 주거환경의 경우 단독주택 거주자는 아파트 대비 2.5점, 다세대주택 거주자는 1.8점 낮았는데, 이는

\begin{table*}[htbp]
  \centering\small
  \caption{분석 변수 설명}
  \label{tab:2-hlm-vars}
  \begin{tabularx}{\textwidth}{l Y c l l}
    \toprule
    \textbf{Variable} & \textbf{Description} & \textbf{Level} & \textbf{Scale} & \textbf{Source} \\
    \midrule
    digital\_use\_score & Digital service utilization score (8 items summed) & L1 & 0--80 (continuous) & Seoul Survey 2023--24 \\
    \addlinespace[2pt]
    age\_c & Age (grand-mean centered, M = 49.3) & L1 & years & Seoul Survey \\
    \addlinespace[2pt]
    female & Female (ref: male) & L1 & 0/1 & Seoul Survey \\
    \addlinespace[2pt]
    Edu & Education: $\leq$middle / high / grad+ (ref: college) & L1 & 0/1 each & Seoul Survey \\
    \addlinespace[2pt]
    Income & Monthly income (grand-mean centered) & L1 & 10K KRW & Seoul Survey \\
    \addlinespace[2pt]
    job & Job: professional / blue-collar / other (ref: white-collar) & L1 & 0/1 each & Seoul Survey \\
    \addlinespace[2pt]
    disabled\_bin & Disability status & L1 & 0/1 & Seoul Survey \\
    \addlinespace[2pt]
    hh\_size\_1 & Single-person household & L1 & 0/1 & Seoul Survey \\
    \addlinespace[2pt]
    housing\_detached / \_multi / \_other & Housing: detached / multi-family / other (ref: apt) & L1 & 0/1 each & Seoul Survey \\
    \addlinespace[2pt]
    year & Year 2024 dummy (ref: 2023) & L1 & 0/1 & Seoul Survey \\
    \addlinespace[2pt]
    elderly & Elderly ($\geq$65) & L1 & 0/1 & Seoul Survey \\
    \addlinespace[2pt]
    low\_edu & Low education ($\leq$middle school) & L1 & 0/1 & Seoul Survey \\
    \addlinespace[2pt]
    connectivity & Connectivity dimension score (normalized) & L2 & 0--1 & UMC Part 1 \\
    \addlinespace[2pt]
    wifi\_total\_n & Public WiFi density per 1K pop (normalized) & L2 & 0--1 & Local Admin License \\
    \bottomrule
  \end{tabularx}
\end{table*}

주택 유형에 따른 소득 수준의 차이뿐 아니라 건물 유형별 인터넷 연결 여건의 차이도 반영된 결과로 해석된다. 1인 가구 역시 다인 가구 대비 1.8점 낮은 점수를 보였다. 그 외 개인 특성으로는, 여성이 남성 대비 2.6점 높게 나타났다. 여성이 남성에 비해 더 높은 디지털 연결성을 갖는 것은, 국내 다양한 디지털 관련 서비스에 여성이 더 민감하게 반응한 결과일 수 있다. \textbf{[인용 필요]} 이러한 결과는 국가에 따른 사회문화적인 맥락을 고려할 필요가 있다. 장애인의 경우 3.1점 높게 나타났는데, 표본의 특성에 기인한 것으로 보인다.

자치구 수준의 UMC 인프라 차원은 주효과만으로는 대체로 유의하지 않았다. Model 2에서 세 차원 모두 통계적 유의성에 도달하지 못하였고, 우도비 검정 결과도 Model 1 대비 유의한 개선이 없었다. 그러나 교차수준 상호작용을 투입한 Model 3에서는 Infrastructure 차원이 유의한 부정적 효과를 보였다. 이는 인프라의 평균적 효과가 아닌, 집단별 차별적 효과가 존재함을 의미한다.

본 분석의 핵심 발견은 교차수준 상호작용에서 나타났다. 공공 WiFi 밀도와 고령 더미의 상호작용은 3.03점의 유의한 정적 효과를 보였다. 이는 공공 WiFi 인프라가 비고령자보다 고령자에게 상대적으로 더 큰 혜택을 제공함을 뜻한다. 디지털 인프라 수준이 높은 자치구에 거주하는 저학력 주민은 인프라 수준이 낮은 곳의 저학력 주민 대비 현저히 높은 디지털 활용 점수를 기록하였다. 이 결과는 인프라 투자가 교육 취약계층의 디지털 격차를 완화하는 데 불균형적으로 큰 효과를 가질 수 있음을 시사한다.

개인 수준 상호작용에서는 이중 취약성이 확인되었다. 독거 고령자는 비독거 고령자 대비 추가적으로 4.0점이 낮았으며, 고령일수록 저학력의 부정적인 효과가 더 커지는 것으로 나타났다. 이러한 결과는 단일 취약 요인보다 복합 취약 요인이 디지털 격차를 심화시킴을 보여준다.

\begin{table*}[htbp]
  \centering\small
  \caption{주요 변수의 기술통계량}
  \label{tab:3-descstat}
  \begin{tabularx}{\textwidth}{Y r r r r r}
    \toprule
    \textbf{변수(Variable)} & \textbf{N} & \textbf{평균(Mean)} & \textbf{표준편차(SD)} & \textbf{최소(Min)} & \textbf{최대(Max)} \\
    \midrule
    디지털 활용 점수(Digital use score) & 10,000 & 60.88 & 16.84 & 0 & 80 \\
    연령(Age) & 10,000 & 49.32 & 16.23 & 15 & 99 \\
    여성 비율(Female, proportion) & 10,000 & 0.51 & --- & 0 & 1 \\
    학력: 중졸 이하(Education: $\leq$Middle school) & 10,000 & 0.07 & --- & 0 & 1 \\
    학력: 고졸(Education: High school) & 10,000 & 0.25 & --- & 0 & 1 \\
    학력: 대학원 이상(Education: Graduate+) & 10,000 & 0.13 & --- & 0 & 1 \\
    장애 여부(Disability) & 10,000 & 0.03 & --- & 0 & 1 \\
    월평균 소득(만원 단위, Monthly income [10K KRW]) & 10,000 & 237.55 & 209.73 & 0 & 1,050 \\
    1인 가구(Single-person household) & 10,000 & 0.17 & --- & 0 & 1 \\
    고령(65세 이상, Elderly 65+) & 10,000 & 0.22 & --- & 0 & 1 \\
    (군집) 연결성 지수(Connectivity, norm.) & 50 & 0.5 & 0.18 & 0.18 & 0.98 \\
    이용 가능성(Available for Use, norm.) & 50 & 0.5 & 0.19 & 0.12 & 0.96 \\
    비용 접근성(Affordability, norm.) & 50 & 0.5 & 0.22 & 0.07 & 1 \\
    \bottomrule
    \multicolumn{6}{p{0.97\textwidth}}{\footnotesize\textit{Note.} N = 10,000명(2023\,$\sim$\,2024년, 25개 자치구, 표본 합산). 소득은 만원 단위. 자치구(군집) 단위 변수는 0\,$\sim$\,1(min-max) 정규화 지수임. N = 10,000 individuals nested within $j$ = 25 districts (2023--2024 pooled). Income measured in 10,000 KRW units. District-level variables are min-max normalized (0--1 scale).}
  \end{tabularx}
\end{table*}

\begin{table*}[htbp]
  \centering\small
  \caption{HLM 추정 결과}
  \label{tab:4-hlm-results}
  \begin{tabularx}{\textwidth}{l Y Y Y}
    \toprule
    \textbf{변수} & \textbf{Model 1} & \textbf{Model 2} & \textbf{Model 3} \\
    \midrule
    \multicolumn{4}{l}{\textit{Level 1: 개인 특성}} \\
    \addlinespace[2pt]
    연령(중심화) & $-0.347^{***}$ (0.012) & $-0.347^{***}$ (0.012) & $-0.278^{***}$ (0.014) \\
    여성 & $2.636^{***}$ (0.305) & $2.634^{***}$ (0.305) & $2.601^{***}$ (0.305) \\
    학력: 중학교 이하 & $-22.618^{***}$ (0.770) & $-22.610^{***}$ (0.770) & $-19.019^{***}$ (2.379) \\
    학력: 고등학교 & $-7.289^{***}$ (0.380) & $-7.282^{***}$ (0.380) & $-7.258^{***}$ (0.380) \\
    학력: 대학원 이상 & $2.906^{***}$ (0.483) & $2.909^{***}$ (0.483) & $2.912^{***}$ (0.481) \\
    소득(중심화) & $0.003^{**}$ (0.001) & $0.003^{**}$ (0.001) & $0.002^{**}$ (0.001) \\
    직업: 전문직 & $2.329^{***}$ (0.451) & $2.326^{***}$ (0.451) & $2.28^{***}$ (0.449) \\
    직업: 생산직 & $-0.373$ (0.424) & $-0.379$ (0.424) & $-0.385$ (0.423) \\
    직업: 기타 & $-1.233^{**}$ (0.456) & $-1.226^{**}$ (0.456) & $-1.105^{*}$ (0.454) \\
    장애 & $3.084^{***}$ (0.888) & $3.081^{***}$ (0.888) & $3.006^{***}$ (0.884) \\
    1인 가구 & $-1.775^{***}$ (0.338) & $-1.776^{***}$ (0.338) & $-0.877^{*}$ (0.367) \\
    단독주택 & $-2.478^{***}$ (0.377) & $-2.480^{***}$ (0.377) & $-2.574^{***}$ (0.379) \\
    다세대주택 & $-1.798^{***}$ (0.408) & $-1.795^{***}$ (0.408) & $-1.933^{***}$ (0.408) \\
    연도(2024) & $0.294$ (0.283) & $0.292$ (0.283) & $0.292$ (0.283) \\
    고령자(65세 이상) & $-6.046^{***}$ (1.696) & & \\
    \addlinespace[4pt]
    \multicolumn{4}{l}{\textit{Level 2: 자치구 특성}} \\
    \addlinespace[2pt]
    연결성 & $-1.127$ (0.747) & $-1.688^{*}$ (0.804) & \\
    이용 가능성(WiFi) & $0.085$ (0.780) & $-0.200$ (0.829) & \\
    비용 접근성 & $-1.578$ (1.135) & & \\
    자치구 소득 & $-2.035^{\dagger}$ (1.181) & & \\
    고령화율 & $-0.298$ (1.095) & $-0.293$ (1.111) & \\
    \addlinespace[4pt]
    \multicolumn{4}{l}{\textit{상호작용}} \\
    \addlinespace[2pt]
    고령자 $\times$ 1인 가구 & $-3.958^{***}$ (0.936) & & \\
    연령 $\times$ 저학력 & $-0.277^{***}$ (0.080) & & \\
    \addlinespace[4pt]
    \multicolumn{4}{l}{\textit{교차수준 상호작용}} \\
    \addlinespace[2pt]
    연결성 $\times$ 고령자 & $-1.920$ (1.901) & & \\
    WiFi $\times$ 고령자 & $3.026^{*}$ (1.515) & & \\
    자치구 소득 $\times$ 고령자 & $2.358$ (1.532) & & \\
    연결성 $\times$ 저학력 & $10.936^{***}$ (3.370) & & \\
    \bottomrule
    \multicolumn{4}{p{0.97\textwidth}}{\footnotesize\textit{Note.} Unstandardized coefficients with standard errors in parentheses. Model 0 = null model (intercept only). Reference categories: male, college education, white-collar job, apartment, multi-person household. $^{\dagger}p<.10$, $^{*}p<.05$, $^{**}p<.01$, $^{***}p<.001$.}
  \end{tabularx}
\end{table*}

\subsection{플랫폼 텍스트 분석: LLM 기반 디지털 연결성 사례분석}
\label{subsec:3-3}

앞선 \ref{subsec:3-1}절과 \ref{subsec:3-2}절에서는 행정 데이터와 설문 조사를 이용해 서울의 디지털 격차를 정량적으로 측정하고, 개인 특성과 지역 인프라의 역할을 분리하였다. 그러나 정량적 지표만으로는 주민들이 일상에서 언제, 어떤 맥락에서 디지털 연결의 결여를 경험하는지까지 포착하기 어렵다. 디지털 연결성에 대한 생활 세계의 체감과 구조적 지표 사이의 거리를 좁히기 위해 정성적 증거를 통합할 필요가 있다.

본 절의 목적은 지역 기반 플랫폼의 비정형 텍스트를 이용해 UMC의 6개 차원(Connectivity, Availability for Use, Affordability, Devices, Skills, Safety)에 잠재된 디지털 결여 사례를 체계적으로 추출하고, 이러한 사례를 \ref{subsec:3-1}절과 \ref{subsec:3-2}절의 정량 분석과 결합하는 방법론적 틀을 제시하는 데 있다. 플랫폼 자체를 분석의 대상으로 삼기보다는, 플랫폼이 제공하는 지역 단위의 생활 경험 데이터를 통해 구조와 경험 사이의 연결을 복원하고자 한다.

\subsubsection{당근 플랫폼 특성}
\label{subsubsec:3-3-1}

비정형 텍스트의 출처는 한국의 지역 기반 커뮤니티 플랫폼 '당근'의 동네생활 게시글이다. 당근을 선정한 이유는 세 가지다. 첫째, 각 게시글에 행정동과 자치구 수준의 지리적 메타데이터가 자동으로 부여되어 \ref{subsec:3-1}절에서 구축한 UMC 지수와 직접적으로 연동할 수 있다. 둘째, 플랫폼 이용자가 서울의 25개 자치구 전역에 골고루 분포해 있어 특정 지역이나 연령층에 치우치지 않는다. 셋째, 게시글은 일상생활 속 디지털 문제를 자발적으로 기록하기 때문에 조사자가 설문에서 미처 포착하지 못한 생활상의 결여를 드러낼 수 있다. 연구의 초점은 플랫폼 자체에 대한 평가가 아니라, 이와 같은 구조적 장점을 활용해 생활 세계의 경험을 수집하는 데 있다.

\subsubsection{왜 LLM인가: 기존 방법론의 한계와 LLM의 구조적 이점}
\label{subsubsec:3-3-2}

UMC의 여섯 차원에 따라 비정형 텍스트를 분류하려면 두 가지 요구를 충족해야 한다. 첫째, 게시글이 디지털 연결성 문제와 관련된 지 여부와 관련된 경우 해당 차원을 정확히 판별해야 한다. 둘째, 일상 언어로 쓰인 구어체 한국어는 조사나 어미 변화가 복잡한 교착어적 특성으로 인해 단순한 형태소 분석기만으로는 정확히 처리하기 어렵다. 기존의 토픽 모델링이나 키워드 기반 분류는 잠재 주제를 파악하는 데 유용하지만 사전 정의된 차원과 일치하지 않거나, 대화체 텍스트를 제대로 분석하지 못하는 한계가 있다. 이러한 한계는 한국어 형태소 분석기가 구어체 데이터에서 정확도가 낮고, 데이터 훈련이 필요한 머신러닝 모델이 충분한 라벨을 요구한다는 점에서 더욱 두드러진다.

대규모 언어 모델(LLM)은 이런 제약을 상당 부분 극복한다. 다국어로 사전 학습되어 한국어의 형태소 분석이 없어도 맥락을 파악할 수 있으며, 104개 이상의 언어를 동시에 다루는 모델은 언어간 의미론적 표현을 공유해 적은 양의 예시만으로도 새로운 태스크에 적응하는 능력을 갖추고 있다. LLM을 도구로 사용할 때는 출력의 확률성과 편차를 고려해야 한다. 동일한 입력에 대해서도 결과가 달라질 수 있으므로, 본 연구는 분석 절차를 고정하고 LLM은 분류와 추론을 수행하는 계산 모듈로만 활용했다. 이를 통해 연구자가 정의한 분류 규칙과 추론 프로토콜을 일관되게 적용하면서도 LLM의 문맥 이해력을 활용할 수 있도록 설계하였다.

\subsubsection{분석 구조}
\label{subsubsec:3-3-3}

텍스트 분석은 네 단계로 구성된다. 먼저, 크롤링한 원시 CSV 데이터에서 삭제되어 접근할 수 없는 게시글과 광고성\textperiodcentered{}중복 게시글을 제거한 뒤, 여섯 차원별 키워드 사전을 이용해 UMC 후보 게시글을 추출한다. 키워드 사전은 디지털 연결성과 직간접적으로 관련된 단어를 정리한 것으로 해당 단어를 포함하는 게시글만 1차로 추출하였다. 이후, 게시글이 디지털 연결성 문제와 관련이 있는지 여부를 판정하는 1차 분류 (umc\_related = Y/N/?)를 수행하며, 판정의 범위를 넓히고자 LLM에게 예시, 제한을 상세하게 밝히지 않고 분석하였다. 그 다음, 관련성이 확인된 게시글만을 대상으로 사용자가 제공한 상세한 예시를 포함한 few-shot 프롬프트를 LLM에게 제시해 게시글을 6개 개념적 차원의 사례로 분류한다. 이 단계에서는 모델의 자율성을 제한하여 프롬프트에 예시를 포함한 구체적인 분류 기준과 명세만을 제공했고, 모델이 판단에 참조하는 정보 공간인 컨텍스트를 엄격하게 관리하였다.

표~\ref{tab:5-text-sample}는 분석 구조와 표본 추출 과정을 요약한 것이다. 원시 데이터 1,287,761건 중 약 51.5\%는 게시자 탈퇴나 차단으로 본문을 볼 수 없어 분석에서 제외되었다. 중복\textperiodcentered{}광고성 게시글 제거 후에는 약 48만 건이 남았으며, 여섯 차원별 키워드 사전을 적용한 결과 26,742건이 UMC 관련 후보로 추출되었다. few-shot 프롬프트를 이용한 1차 판정에서는 약 57.7\%(15,438건)가 관련(Y)으로, 29.6\%(7,921건)가 비관련(N)으로, 12.6\%(3,383건)가 판단 유보(?)로 분류되었다. 이 가운데 Y로 판정된 게시글만을 대상으로 차원 분류를 진행했다.

차원 분류 판정에서 LLM에게 각 차원에서 대표적인 사례를 담은 3\,$\sim$\,4개의 예시를 프롬프트로 제공했다. 이러한 few-shot 프롬프트는 모델이 사람의 판단 기준을 학습하는 데 도움이 되며, 복잡한 작업에서 zero-shot보다 나은 성능을 보인다~\textbf{[인용 필요]}. 예시에는 어떤 유형의 게시글이 디지털 연결성 문제와 무관한지를 보여주는 음성 예시도 포함해, 관련성을 최대한 좁혀 높은 정밀도를 확보하도록 했다. 2차 차원 분류에서는 구체적인 분류 기준과 정의만을 프롬프트로 제공해 모델이 보다 자율적으로 사례를 분류하도록 설계했다. 이처럼 앞 단계는 예시 기반 few-shot 방식이고, 뒤 단계는 규칙 기반 프롬프트를 사용한다는 점에서 추론과 few-shot 사이의 균형을 보여준다(그림~\ref{fig:3-pipeline-roles} 참조).

\begin{table}[htbp]
  \centering\small
  \caption{텍스트 분석 표본 추출 및 판정 절차}
  \label{tab:5-text-sample}
  \begin{tabularx}{\linewidth}{l Y r}
    \toprule
    \textbf{단계} & \textbf{설명} & \textbf{남은 게시글 수} \\
    \midrule
    원시 데이터 & 2024년 5월\,$\sim$\,2026년 3월에 작성된 동네생활 게시글 전수 & 1,287,761 \\
    접근가능 게시글 & 탈퇴\textperiodcentered{}차단으로 본문 확인이 불가능한 게시글 제거 & 624,564 \\
    중복\textperiodcentered{}광고 제거 & 중복된 내용, 광고\textperiodcentered{}홍보 게시글 제거 & 480,972 \\
    UMC 키워드 필터 & 6개 차원별 키워드 사전을 적용해 후보 추출 & 26,742 \\
    1차 판정 & few-shot 프롬프트로 UMC 관련성 분류 & --- \\
    \bottomrule
  \end{tabularx}
\end{table}

\fig{image18.png}{분류와 추론의 역할 구분}{fig:3-pipeline-roles}

이후, 분류 결과를 자치구 단위로 집계하여 표출 프로파일을 만들고, \ref{subsec:3-1}절의 정량 지표와 결합하기 위해 베이지안 업데이트를 적용한다. 정량 지표와 텍스트 분류 결과를 결합하기 위해 베이지안 업데이트를 적용하였다. 사전 분포는 UMC의 6가지 개념의 잠재적 공간이다. 차원 $d$에서 자치구 $j$의 UMC 실측값을 $s_{jd}$로 할 때, 사전 분포는 $\mathrm{Beta}(\alpha_{jd}, \beta_{jd})$이며, 평균이 $s_{jd}$가 되도록 사전 분포를 가정한다. 강도 모수 $\kappa$를 사용해 $\alpha_{jd} = s_{jd}\cdot\kappa$, $\beta_{jd} = (1-s_{jd})\cdot\kappa$로 표현한다. 텍스트 분류에서 해당 차원으로 판정된 게시글 수를 $y_{jd}$, 총 분류된 게시글 수를 $n_{jd}$라 할 때, 우도는 $\mathrm{Binomial}(n_{jd}, \theta_{jd})$로 표현된다. 사후 분포는 $\theta_{jd}\mid y_{jd} \sim \mathrm{Beta}(\alpha_{jd}+y_{jd}, \beta_{jd}+n_{jd}-y_{jd})$로 계산된다. 사전 강도 $\kappa$는 기본값 20을 사용하고 10과 30에 대한 민감도 분석도 수행하였다. 게시글 수가 적은 자치구에서 사후 분포가 사전 분포에 의해 지나치게 지배되는 것을 막기 위해, 게시글 수 $n_{jd}$를 생활인구로 보정한 $n_{jd}^{*} = n_{jd}\times(\text{생활인구}_{j}/\overline{\text{생활인구}})$로 사용하였다.

\begin{align}
\alpha_{jd} &= \kappa \times s_{jd} \label{eq:alpha} \\
\beta_{jd}  &= \kappa \times (1 - s_{jd}) \label{eq:beta} \\
\theta_{jd}\mid y_{jd} &\sim \mathrm{Beta}\bigl(\alpha_{jd}+y_{jd},\, \beta_{jd}+n_{jd}-y_{jd}\bigr) \label{eq:posterior} \\
n_{jd}^{*} &= n_{jd}\times\bigl(\text{생활인구}_{j}/\overline{\text{생활인구}}\bigr) \label{eq:popadj}
\end{align}

베이지안 집계에서 게시글로 추정된 사후 분포가 사전 분포로 유의미하게 이동한 자치구\textperiodcentered{}차원 조합을 대상으로 인과기제에 관한 정성적 추론을 수행하였다. 구체적으로는 (1) 사후 이동이 통계적으로 유의하며, (2) 해당 자치구\textperiodcentered{}차원에서 Y 판정 게시글이 최소 5건 이상이고, (3) \ref{subsec:3-2}절의 다층모형 결과가 허용하는 범위 내에서만 해석을 실시하였다. 사전 분포에 의해 도출된 조합에 대해 대표 게시글 3\,$\sim$\,5건을 선정하고, 1\,$\sim$\,2문장으로 구조적 지표와 생활 경험 사이의 연결 경로를 서술하였다. 분류와 추론 작업은 세 명의 독립된 에이전트가 서로 다른 추론 경로를 통해 수행했으며, 최종 판정은 네 번째 에이전트가 세 결과를 비교\textperiodcentered{}검증하는 방식으로 이루어졌다. 이는 한 가지 프로토콜에 의존하지 않고 여러 관점의 상호 검증을 통해 편향을 통제하는 삼각검증(triangulation) 접근을 따른 것이다.

모든 분류와 추론 작업은 세 명의 독립된 분석 에이전트가 서로 다른 추론 방식 프롬프트와 컨텍스트를 사용해 수행하고, 네 번째 분석자가 세 결과를 비교\textperiodcentered{}검증하는 방식으로 이루어졌다. 이는 질적 연구에서 여러 자료원과 분석가를 동원해 편향을 줄이는 삼각검증(triangulation) 접근에 기반한다. 세 관점이 일치할 때만 해당 게시글을 최종 분석에 포함했으며, 해석상의 불확실성은 결과 표기에 반영하였다. 이러한 과정을 정리한 결과는 표~\ref{tab:6-units}, 그림~\ref{fig:3-pipeline-full}과 같다.

\begin{table}[htbp]
  \centering\small
  \caption{\ref{subsec:3-3}절의 분석 단위와 해석 단위}
  \label{tab:6-units}
  \begin{tabularx}{\linewidth}{l l Y Y}
    \toprule
    \textbf{구분} & \textbf{단위} & \textbf{무엇을 보여주는가} & \textbf{해석상 주의점} \\
    \midrule
    관찰 단위 & 개별 게시글 & 디지털 연결성 경험의 단서 & 단일 게시글만으로 지역 구조를 주장할 수 없음 \\
    패턴 단위 & 행정동 게시글 & 차원별 표출 비율과 분포 & 집계 기준과 표본 편향을 함께 검토해야 함 \\
    구조 단위 & UMC 지역 지표 & 지역 수준의 연결성 조건 & 텍스트 집계 결과와의 정합을 통해 제한적으로 해석 \\
    인과기제 단위 & 추론 메커니즘 & 사례와 지역 패턴의 연결 설명 & \ref{subsec:3-2}의 정량 결과가 허용하는 범위 안에서만 해석 \\
    \bottomrule
  \end{tabularx}
\end{table}

\subsubsection{데이터와 적용 범위}
\label{subsubsec:3-3-4}

원시 데이터는 서울 25개 자치구를 대상으로 2024년 5월부터 2026년 3월까지 작성된 동네생활 게시글 전수를 크롤링한 것으로, 중복 제거 후 총 1,287,761건이다. 이 중 삭제(DELETED) 663,163건(51.5\%)과 차단(BLOCKED) 103,728건(8.1\%)을 합한 766,891건(59.6\%)은 본문 확인이 불가능하여 분석에서 제외하였다. 분석 대상은 최종 약 52만 건이다.

표본 탈락의 체계적 편향 가능성을 검토한 결과, 자치구별 제거율의 표준편차는 4.6\%로 자치구 간 편차가 크지 않았으며, 디지털 결여 추론의 주요 대상인 생활/편의(47.0\%), 동네사건사고(48.4\%), 병원/약국(35.5\%) 카테고리는 사교\textperiodcentered{}홍보 목적 카테고리(동네친구 86.6\%, 교육 69.0\%)에 비해 제거율이 낮았다. 삭제의 상당 부분은 사용자 탈퇴에 수반된 일괄 삭제에 기인하는 것으로 판단된다. 다만, 삭제 게시글 중 일부가 디지털 결여 관련 사례를 포함할 가능성을 완전히 배제할 수는 없고, 삭제 행위 자체가 디지털 결여와 체계적으로 연관되었을 수 있다. 이는 본 분석의 한계에 해당한다.

당근 데이터의 구조적 한계도 인식할 필요가 있다. 플랫폼 이용자 자체가 디지털 접근이 가능한 집단으로 한정되므로, 디지털 소외 계층의 경험은 구조적으로 과소 대표된다. 게시글의 시간적\textperiodcentered{}공간적 분포가 균일하지 않아 자치구 간 비교에서 표본 편향의 가능성이 존재한다. 이러한 한계는 데이터 소스의 제약이지 방법론의 제약이 아니다. 보다 직접적인 질적 자료가 존재하는 맥락에서는 동일한 파이프라인을 적용하여 보다 정밀한 분석이 가능할 것이다.

\figfull{image19.png}{LLM 기반 텍스트 분석 파이프라인}{fig:3-pipeline-full}

\subsubsection{분석 결과}
\label{subsubsec:3-3-5}

텍스트 분류 결과는 표~\ref{tab:6-units}와 같다. 연결 품질(Connection Quality) 차원은 6개 차원 가운데 전이의 폭이 가장 작다($-0.006 \sim +0.005$). 실측과 텍스트 증거가 대체로 수렴하는 차원으로 볼 수 있다. 양천구는 실측 UMC 0.278로 낮은 편이나 관련 게시글 비율이 상대적으로 높아 $+0.005$의 전이를 보였고, 용산구는 실측 0.884로 높은 반면 관련 게시글이 드물어 $-0.006$의 전이를 보였다.

이용 가능성(Availability for Use) 차원에서는 도봉구가 실측 최하위(0.213)에 해당하면서도 관련 게시글 비율이 높아 $+0.006$의 전이를 보였다. 광진구는 실측 0.832로 높고 관련 게시글이 드물어 $-0.010$의 전이를 기록하였다.

경제적 접근성(Affordability) 차원은 서울특별시 자치구의 소득 수준과 일관된 방향의 경향을 보였다. 중랑구와 강서구는 각각 $+0.005$의 양의 전이를 보인 반면, 서초구는 실측이 0.993으로 최고 수준이면서 관련 게시글이 드물어 $-0.007$의 전이를 나타냈다. 이 차원은 통신비와 기기 구입 비용에 관한 고민이 소득 수준에 따라 차등적으로 표출되는 경향과 정합적이다.

기기 접근성(Devices) 차원은 자치구 간 편차가 가장 큰 차원이다(gap = 0.090). 중랑구는 실측 0.052에 해당하면서도 기기 관련 게시글 비율이 4.95\%로 전 자치구 가운데 가장 높아 $+0.048$의 전이를 보였다. 송파구 또한 $+0.026$의 유사한 패턴을 나타냈다. 반면 용산구($-0.042$)와 동작구($-0.033$)는 실측이 높으면서 관련 게시글이 드물게 나타났다. 다만 이 차원에서 실측 지표는 기기 보유율을 포착하는 반면, 텍스트 증거는 중고 거래, 수리 요청 등 기기와 관련한 일상적 거래와 고충을 포착하므로, 두 지표가 지시하는 대상이 부분적으로 상이할 가능성을 배제하기 어렵다.

디지털 역량(Digital Skills) 차원은 기기 차원에 이어 두 번째로 편차가 큰 차원이다(gap = 0.092). 중랑구는 실측 0.242에 관련 게시글 비율 4.8\%로 $+0.047$의 전이를 보였다. 강남구 또한 실측이 0.276으로 낮은 편이면서 관련 게시글 비율이 상대적으로 높아 $+0.013$의 전이를 나타냈다. 반면 관악구는 실측 0.795로 최고 수준이면서 관련 게시글이 드물어 $-0.045$의 전이를 보였다.

디지털 역량(Digital Skills) 차원에서는 기기 차원에 이어 두 번째로 큰 편차가 나타났다. 중랑구는 실측 0.242에 비해 관련 게시글 비율 4.8\%로 $+0.047$의 이동을 보였으며, 강서구는 실측값이 높음에도 관련 게시글이 적어 $-0.031$의 이동을 보였다. 안전\textperiodcentered{}보안(Safety) 차원은 피싱 및 보이스피싱 사례가 게시글의 주요 내용을 이루었고, 중랑구($+0.018$)와 도봉구($+0.015$)에서 양의 이동이 확인되었다. 종로구는 실측값이 매우 높고 관련 게시글이 적어 사후 분포가 낮아졌다.

안전\textperiodcentered{}보안(Safety \& Security) 차원에서는 중랑구($+0.018$)와 도봉구($+0.015$)에서 피싱 사기 및 보이스피싱 관련 게시글의 비율이 높았다. 종로구는 실측 0.984로 최고 수준이면서 관련 게시글이 드물어 $-0.019$의 전이를 나타냈는데, 관광\textperiodcentered{}상업 중심지의 주거 구성과 당근 이용자 집단의 특성이 상이할 가능성을 함께 고려할 필요가 있다.

자치구 단위로 집계했을 때, 지역의 실측값과 게시글 분포 간 차이는 크게 세 가지 유형. 첫째, 평균 전이가 양의 방향으로 나타나는 자치구이다. 중랑구는 6개 차원 중 4개 차원에서 양의 전이가 가장 크게 나타나 평균 $+0.020$으로 가장 높은 값을 보였다. 도봉구와 강남구 또한 평균적으로 양의 방향을 보였으나 그 크기는 0.001 수준으로 작았다. 특히 강남구의 경우 평균 소득이 높은 자치구임에도 이용 가능성과 디지털 역량 차원에서 양의 전이가 관찰되었는데, 자치구 내부의 이질적 집단이 관련 이슈를 표출하고 있을 가능성을 시사한다. 둘째, 전이가 0에 가까운 자치구이다. 강북구($-0.000$), 강동구($-0.001$), 동대문구($-0.001$)가 이에 해당한다. 이들 자치구에서는 차원별로 양\textperiodcentered{}음의 전이가 섞여 있어 평균적으로 상쇄되는 양상을 보인다. 셋째, 평균 전이가 음의 방향으로 나타나는 자치구이다. 마포구($-0.014$), 영등포구($-0.013$), 관악구($-0.013$), 서초구($-0.009$)가 이에 해당한다. 관악구는 대학가 중심의 인구 구성과 디지털 역량 실측값이 맞물려 도움 요청 성격의 게시글이 드물게 나타난 것으로 볼 수 있으며, 서초구는 경제적 접근성과 기기 차원에서 특히 큰 음의 전이를 보였다. 사전 강도 $\kappa$를 10, 20, 30으로 달리한 민감도 분석에서도 사후 평균의 차이는 $0.00001\sim0.00010$ 수준으로 실질적 변화가 없었다. 이는 생활인구 보정 후 유효 관측 수가 충분해 사전의 가중치 변화가 결과에 큰 영향을 미치지 않음을 시사한다.

해석에서 주의할 점은 양의 전이가 곧바로 실측 지표가 과소평가되었음을 의미하지 않는다는 것이다. 양의 이동은 실측 지표가 포착하지 못한 생활상의 불편을 반영할 수 있지만, 특정 자치구에서 특정 문제에 대한 논의가 활발한 문화적\textperiodcentered{}인구학적 요인을 동시에 포함할 수 있다~\textbf{[인용 필요]}. 또한 플랫폼 이용자의 표집 편향과 게시글의 시공간적 불균형이 결과 해석에 영향을 미칠 수 있다. 따라서 사후 분포의 방향과 크기를 정책적 함의로 연결할 때는 정량\textperiodcentered{}정성적 지표를 종합적으로 고려해야 한다. 본 절의 결과에서 유의미한 것은 전이의 부호와 크기가 차원과 자치구에 따라 체계적으로 달라진다는 것이다. 단순히 당근마켓 게시글의 빈도만을 집계할 경우 플랫폼 활성도 자체가 UMC 수준에 대한 추정과 혼동될 우려가 있다. 사전 분포로 실측값을 투입하고 텍스트 분류 결과를 우도로 결합하는 절차는, 관측된 텍스트 빈도에서 플랫폼 활성도에 기인하는 부분과 UMC 수준에 기인하는 부분을 부분적으로 분리하는 효과를 갖는다. 이 점에서 사전-사후 분포의 전이는 텍스트 증거가 실측을 얼마나 보정하는지를 보여주는 지표로 기능한다.

다만 shift의 방향을 해석할 때 두 가지 제약을 함께 고려해야 한다. 첫째, 양의 shift가 관찰된다고 해서 해당 자치구의 UMC 수준이 실측보다 낮다고 곧바로 결론지을 수 없다. 양의 shift는 실측 지표가 포착하지 못하는 생활 수준의 체감이 존재할 가능성을 시사하지만, 관련 주제가 당근마켓이라는 특정 플랫폼에서 표출되는 방식 자체가 자치구별 이용자 구성에 의존하기 때문이다. 둘째, 음의 shift 또한 해당 자치구의 UMC 수준이 실측보다 높다는 의미로 해석하기보다, 관련 이슈가 플랫폼 상에서 표출될 유인이 상대적으로 적은 상태로 읽는 것이

\begin{table*}[htbp]
  \centering\small
  \caption{자치구별 게시글 판정 결과}
  \label{tab:7-district-cls}
  \begin{tabularx}{\textwidth}{l r r r r r}
    \toprule
    \textbf{자치구} & \textbf{총 게시글(n)} & \textbf{Y 판정} & \textbf{? 판정} & \textbf{N 판정} & \textbf{Y 비율(\%)} \\
    \midrule
    \textbf{중랑구}   & 5,476  & 794 & 28  & 4,654  & 14.5 \\
    \textbf{송파구}   & 9,367  & 860 & 99  & 8,408  & 9.18 \\
    \textbf{서초구}   & 6,170  & 434 & 216 & 5,520  & 7.03 \\
    \textbf{은평구}   & 6,140  & 429 & 50  & 5,661  & 6.99 \\
    \textbf{도봉구}   & 3,491  & 233 & 5   & 3,253  & 6.67 \\
    \textbf{서대문구} & 4,500  & 289 & 5   & 4,206  & 6.42 \\
    \textbf{강북구}   & 3,885  & 247 & 31  & 3,607  & 6.36 \\
    \textbf{용산구}   & 3,012  & 181 & 12  & 2,819  & 6.01 \\
    \textbf{양천구}   & 4,975  & 296 & 7   & 4,672  & 5.95 \\
    \textbf{강동구}   & 6,254  & 342 & 231 & 5,681  & 5.47 \\
    \textbf{영등포구} & 5,964  & 306 & 5   & 5,653  & 5.13 \\
    \textbf{동대문구} & 4,754  & 238 & 29  & 4,487  & 5.01 \\
    \textbf{성북구}   & 5,437  & 266 & 141 & 5,030  & 4.89 \\
    \textbf{중구}     & 1,351  & 63  & 16  & 1,272  & 4.66 \\
    \textbf{강서구}   & 8,174  & 371 & 236 & 7,567  & 4.54 \\
    \textbf{금천구}   & 2,298  & 101 & 7   & 2,190  & 4.40 \\
    \textbf{종로구}   & 2,267  & 97  & 27  & 2,143  & 4.28 \\
    \textbf{동작구}   & 3,675  & 150 & 22  & 3,503  & 4.08 \\
    \textbf{강남구}   & 10,047 & 404 & 459 & 9,184  & 4.02 \\
    \textbf{성동구}   & 404    & 16  & 9   & 379    & 3.96 \\
    \textbf{관악구}   & 11,665 & 407 & 75  & 11,183 & 3.49 \\
    \textbf{마포구}   & 7,974  & 242 & 254 & 7,478  & 3.03 \\
    \textbf{구로구}   & 4,082  & 121 & 96  & 3,865  & 2.96 \\
    \textbf{광진구}   & 5,440  & 147 & 42  & 5,251  & 2.70 \\
    \textbf{노원구}   & 5,040  & 103 & 46  & 4,891  & 2.04 \\
    \bottomrule
  \end{tabularx}
\end{table*}

\figfull{image20.png}{지역별 사전-사후 분포 차이}{fig:3-prior-posterior}

안전하다. 중랑구와 마포구의 대조에서 드러나듯, shift의 해석은 자치구의 인구 구성과 플랫폼 이용 특성을 비롯한 다양한 맥락 정보와 결합될 때 의미를 띈다. 둘째, 음의 전이 또한 해당 자치구의 UMC 수준이 실측보다 높다는 의미로 해석하기보다, 관련 이슈가 플랫폼 상에서 표출될 유인이 상대적으로 적은 상태로 읽는 것이 타당할 수 있다. 이런 측면에서 지역 자원이 실제로 결여되어 있음에도 플랫폼 상에서 논의되지 않는 이유를 고려하는 것 또한 중요하다.

\begin{table*}[htbp]
  \centering\small
  \caption{결여 유형별 가설 생성 결과 분포}
  \label{tab:8-deficit-types}
  \begin{tabularx}{\textwidth}{l Y r r}
    \toprule
    \textbf{결여 유형} & \textbf{관련 UMC 차원} & \textbf{관련 게시글 수 (n=100)} & \textbf{비율 (\%)} \\
    \midrule
    정보부재         & Digital Skills      & 63 & 63.0 \\
    제도적 경로 결여 & Available for Use   & 61 & 61.0 \\
    기기 접근 결여   & Devices             & 50 & 50.0 \\
    디지털 역량 부재 & Digital Skills      & 45 & 45.0 \\
    연결 인프라 결여 & Connectivity        & 40 & 40.0 \\
    공공 접근점 결여 & Available for Use   & 32 & 32.0 \\
    경제적 접근 결여 & Affordability       & 27 & 27.0 \\
    안전\textperiodcentered{}보안 결여 & Safety & 16 & 16.0 \\
    복합 결여 (2개 이상) & ---             & 88 & 88.0 \\
    \bottomrule
    \multicolumn{4}{p{0.97\textwidth}}{\footnotesize\textit{Note.} 하나의 게시글에서 복수의 결여 유형이 추론될 수 있으므로 비율의 합은 100\%를 초과한다. Phase 1은 가설 생성 단계이며, Phase 2(판정)에서 삼각검증을 거쳐 최종 분류가 확정된다.}
  \end{tabularx}
\end{table*}

\subsection{분석 결과 종합}
\label{subsec:3-4}

본 절에서는 \ref{subsec:3-1}절(UMC 지수 구축), \ref{subsec:3-2}절(HLM 다수준 분석), \ref{subsec:3-3}절(LLM 기반 텍스트 분석)의 결과를 종합하여, 서울의 디지털 격차가 어떤 구조로 작동하는지를 통합적으로 해석한다. 세 분석은 각각 '어디에 격차가 있는가'(측정), '왜 격차가 발생하는가'(인과 규명), '격차가 어떻게 경험되는가'(질적 맥락)에 답하며, 이들의 교차 지점에서 정책적으로 유의미한 함의가 도출된다.

세 분석이 공통으로 지지하는 첫 번째 발견이자 본 프로젝트의 핵심 실증 결과는, 서울의 디지털 격차가 자치구 간 차이가 아닌 개인 간 차이에 의해 구조화되어 있다는 것이다. \ref{subsec:3-2}절의 HLM 분석에서 ICC는 0.49\%에 불과하여, 디지털 활용 점수 분산의 99.5\%가 동일 자치구 내 개인 간 차이에 기인하였다. 이 결과는 \ref{subsec:3-1}절에서 구축한 UMC 종합지수의 자치구 간 편차가 실재함에도 불구하고, 해당 편차가 개인의 디지털 활용 수준에 미치는 직접적 영향은 제한적임을 의미한다. 그러나 이러한 낮은 ICC는 자치구 수준 변수의 무의미함을 뜻하는 것이 아니라, 그 효과가 집단별로 차별적으로 발현됨을 시사한다. Model 3의 교차수준 상호작용 결과가 이를 뒷받침한다.

\ref{subsec:3-3}절의 텍스트 분석은 이 양적 발견이 일상에서 어떻게 발현되는지를 구체화한다. '정보부재'가 가장 빈번한 결여 유형(63\%)으로 나타난 것은, 교육수준이 낮을수록 디지털 서비스의 존재 자체를 인지하지 못하는 정보 격차가 더 심화될 수 있음을 시사한다. 프린터 이용을 위해 편의점 무인 인쇄를 모르는 사례, ISP 고객센터 번호를 모르는 사례, 정부24의 공문서 출력 기능을 모르는 사례 등은 모두 '서비스가 존재하지만 이용자에게 보이지 않는' 구조적 문제를 가리킨다. 이는 HLM에서 교육 기울기가 가파른 이유를 경험적 수준에서 설명하는 보완적 증거로 해석할 수 있다.

두 번째 발견은, 자치구 인프라의 효과가 주효과 수준에서는 유의하지 않으나 특정 취약집단에 대해서는 조절 효과로서 유의하게 작용한다는 것이다. Model 2에서 UMC 차원 중 Connectivity와 관련된 인프라 변수의 주효과는 통계적 유의성에 도달하지 못하였는데($p = 0.237$), 이는 ICC가 0.49\%인 상황에서 자치구 수준 변수의 주효과가 제한적인 것과 일관된 결과이다. 그러나 Model 3에서 교차수준 상호작용이 유의하게 나타남으로써, 인프라의 효과가 평균적 효과가 아닌 집단별 조절 효과로서 작동함이 확인되었다.

\ref{subsec:3-3}절의 텍스트 분석은 이 차별적 효과의 질적 맥락을 제공한다. 공공 접근점(프린트 카페, 도서관 복합기, 무인 인쇄)에 대한 정보부재가 32\%의 게시글에서 추론된 것은, 인프라가 물리적으로 존재하더라도 이용자의 인지와 접근 역량이 수반되지 않으면 실질적 혜택으로 전환되지 않음을 보여준다. 은평구 신사동의 사례('정녕 없는 건지요')는 온라인 검색을 이미 시도했으나 결과가 불충분했을 가능성을 시사하며, 이는 인프라의 존재와 인프라의 '발견가능성(discoverability)' 사이의 간극을 조명한다. HLM에서 인프라 주효과가 유의하지 않았던 결과와 결합하면, 인프라 투자의 효과는 취약계층의 정보 접근과 역량 지원이 병행될 때 비로소 발현된다는 해석이 가능하다.

세 번째 발견은, 디지털 격차가 단일 축이 아닌 복합 취약성의 교차로 심화된다는 것이다. HLM에서 독거 고령자는 비독거 고령자 대비 추가적으로 4.0점이 낮았고, 연령 증가에 따른 저학력의 부정적 효과가 가속화되는 패턴이 확인되었다. \ref{subsec:3-3}절에서도 결여의 복합성이 일관되게 나타났다. 게시글당 평균 2.5개 이상의 결여 유형이 동시에 식별되었으며, 정보부재\textperiodcentered{}역량부재\textperiodcentered{}제도적 경로 차단이 중첩될 때 이용자가 공식 경로에서 이탈하여 커뮤니티에 의존하는 구조가 확인되었다. 디지털 지원망의 부재('자녀분들이 없으시거나 떨어져 살아서 묻기 힘들거나')는 가구 구성이 디지털 역량 획득의 비공식적 경로로 기능함을 보여주며, 독거 고령자의 추가적 격차를 설명하는 기제로 해석된다.

네 번째 발견은, 디지털 사막의 공간적 군집과 차원별 이질성이다. \ref{subsec:3-1}절에서 하위 5개 지역 중 특히 서울 북부(도봉구, 강북구, 노원구, 중랑구)가 UMC 종합지수의 Low-Low 클러스터를 형성하였으며, 이들은 Connectivity 차원에서는 도시 평균 이상을 기록하면서도 Devices와 Safety 차원에서 극단적 결핍을 보였다. 이러한 차원별 프로파일의 이질성은 일률적 인프라 투자보다 차원별 맞춤 개입의 필요성을 시사한다.

방법론적 관점에서, 세 분석의 결합은 각각의 한계를 상호 보완한다. \ref{subsec:3-1}절의 종합지수는 자치구 간 상대적 위치를 파악하되 개인 내 변이를 포착하지 못하며, \ref{subsec:3-2}절의 HLM은 개인과 지역의 기여를 분리하되 디지털 격차의 질적 발현 경로를 설명하지 못한다. \ref{subsec:3-3}절의 텍스트 분석은 결여의 경험적 맥락을 가시화하되 대표성에 한계가 있다. 세 분석을 병렬이 아닌 순차적 삼각검증(측정 $\rightarrow$ 인과 규명 $\rightarrow$ 질적 맥락)으로 결합함으로써, '어디에'(자치구), '왜'(개인 특성과 인프라의 상호작용), '어떻게'(결여의 발현 경로)라는 세 질문에 대한 통합적 답을 구성할 수 있었다.

이상의 종합 결과는 4장의 정책 제언으로 직결된다. 개인 특성이 지배적인 설명변수라는 발견은 people-based 정책의 우선성을, 인프라의 차별적 효과는 취약계층 타겟형 place-based 정책의 필요성을, 복합 취약성의 확인은 단일 프로그램이 아닌 통합 지원 체계의 필요성을, 공간적 군집의 확인은 권역 단위 거버넌스의 근거를 각각 제공한다.
)
\section{분석}
\label{sec:analysis}

\noindent{\sffamily\small\color{muted}\textbf{CHAPTER 3}}

\begin{quoteblock}
측정 $\rightarrow$ HLM $\rightarrow$ 베이지안 통합의 3단계 분석으로 서울 디지털 격차의 공간 분포, 개인-지역 기여 분리, 그리고 텍스트 기반 잠재 격차 신호까지 통합 추정한다.
\end{quoteblock}

\figfull{image6.png}{핵심 발견 요약}{fig:3-keyfindings}
\figfull{image7.png}{디지털 사막 --- 종합지수 하위 4개 자치구}{fig:3-desert}

본 장의 세 분석은 순차적으로 연결된다. 먼저 서울 자치구별 UMC 지수를 구축하여 지역 격차를 측정하고(\ref{subsec:3-1}), 이 결과를 지역 수준 변수로 활용하여 개인 및 지역 특성이 개인의 디지털 연결성에 미치는 영향을 분석한다(\ref{subsec:3-2}). 마지막으로 디지털 플랫폼 텍스트를 통해 양적 분석이 포착하지 못하는 구체적 사례를 도출한다. 이 흐름은 측정 $\rightarrow$ 인과 규명 $\rightarrow$ 현상 보완으로 구성되며, \ref{subsec:3-4}절에서 세 분석의 결과를 종합한다.

\subsection{서울의 UMC 지역 격차}
\label{subsec:3-1}

\subsubsection{측정 지표 및 데이터}
\label{subsubsec:3-1-1}

본 프로젝트에서 구성한 측정 지표는 UMC의 6개 차원에 대한 개념적 정의를 바탕으로 한다. 각 차원은 Connectivity, Available for Use, Affordability, Devices, Skills, Safety로 구분되며, 개별 측정 지표는 의미성과 보편성을 동시에 고려한 두 가지 핵심 구성 요소에 따라 설정하였다~\citep{ITU2023}. 본 분석에 활용된 서울의 UMC를 구성하는 6개 차원별 측정 지표는 표~\ref{tab:1-umc-dims}과 같다.

측정 지표 선정 과정에서는 개념적 정의와의 부합성, 서울의 디지털 기술 환경, 데이터 가용성을 종합적으로 검토하였다. 분석에는 데이터를 확보할 수 있는 가장 최근 시점인 2023년과 2024년, 총 2개 시점의 자료를 활용하였다. 2개 시점의 국가 행정 자료 등을 활용하였다. 단, Devices, Skills, Security 차원은 데이터 제공 시점의 한계로 인해 2023년 조사 자료를 기준으로 두 시점 모두 동일한 값을 적용하였다.

Connectivity의 경우, 한국전파누리에서 제공하는 이동통신 무선국 데이터와 서울시 행정데이터인 생활인구 데이터를 활용하여 측정 지표를 산출하였다. \citet{ITU2023}에 따르면, Connectivity는 네트워크 커버리지 내에 거주하는 인구를 기준으로 측정될 수 있다. 그러나 한국은 직주 분리가 뚜렷한 경우가 많아, 거주지역과 실질적으로 디지털 환경에 접속하는 지역이 상이할 수 있다. 이에 본 연구에서는 단순 거주인구 대신 주간 유동인구 데이터를 활용하여 4G 및 5G 네트워크 커버리지를 산출함으로써, 한국의 맥락에 보다 적합한 지표를 고안하였다. 구체적으로, 4G 기지국 수와 5G 기지국 수를 각각 주간 생활인구 천 명당 기지국 수로 환산하여 집계구별로 산출한 뒤, 이를 자치구별로 평균한 값을 최종 지표로 사용하였다.

Available for Use는 디지털 서비스의 실질적 이용 가능성을 측정하는 차원이다. SKT 통신정보에서 제공하는 월별 지역 내 가구당 모바일 데이터 사용량과 금융, 쇼핑, 동영상, 배달 4개 범주의 온라인 서비스 이용일수를 인구가중 평균으로 자치구별 월별로 집계하고, 12개월 평균값을 사용하였다. 여기에 자치구 지방정부의 공공WiFi 등록 현황 데이터를 활용하여 전체 인구 1,000명당 공공WiFi AP 수를 산출하여 포함하였다. 공공WiFi 밀도는 가정 내 인터넷 접속 환경이 제한적인 고령층에게 특히 중요한 디지털 접근 경로로 기능한다~\citep{vanDijk2020}.

Affordability는 ITU ICT Development Index 2023의 두 가지 축 구조를 따라 구성하였다~\citep{ITU2024}. 구매력(capacity) 측면에서는 서울시 상권분석서비스에서 제공하는 자치구별 월평균 가구소득을 활용하고, 가격부담(burden) 측면에서는 SKT 통신정보의 최근 3개월 이내 요금 연체 비율을 활용하였다. 연체율은 행정동, 성별, 연령대별 인구가중 평균으로 산출한 뒤 자치구 단위로 집계하였으며, 값이 높을수록 경제적 접근성이 낮음을 의미하는 역방향 지표이다.

Devices는 디지털 기기 보유 및 이용의 다양성을 측정한다. 서울시민 디지털역량실태조사의 가구 내 핵심 디지털 기기 보유 유형 수와 개인이 실제로 사용하는 기기 유형 수를 자치구별 가중평균으로 산출하였다. \citet{ITU2023}의 보고서에 따르면, 휴대폰과 컴퓨터 등 기기의 특성에 따라서 기기의 보유와 이용은 서로 다르게 해석된다. 이를 한국 국내 맥락에선, 컴퓨터의 보유 여부도 디지털 연결성을 드러내는 지표가 될 수 있다. 스마트폰 보급률이 사실상 포화 상태인 서울에서는 데스크톱, 노트북, 태블릿 등 보조 기기의 보유 여부가 디지털 활용 역량의 차별화 요인으로 작용한다~\citep{vanDeursenHelsper2015}.

Digital Skills는 기초, 응용, 고급, 사이버활동의 4개 하위 영역으로 구성된다. 기초 역량은 스마트 기기 기본 조작 능력을, 응용 역량은 디지털 서비스 활용 능력을, 고급 역량은 스마트 기기 활동 숙련도를, 사이버활동 역량은 사이버 공간에서의 활동 능력을 각각 측정한다. 각 하위 영역의 원점수를 만점으로 나누어 0\,$\sim$\,1 비율로 변환한 뒤 자치구별 가중평균을 산출하였다.

Safety는 디지털 보안 행동과 보안 인식의 2개 하위 지표로 구성된다. 보안 행동은 비밀번호 관리, 소프트웨어 업데이트, 2단계 인증 사용 등 실제 보안 실천 행위를 측정하고, 보안 인식은 개인정보 보호의 중요성에 대한 인지 수준을 측정한다. 각 하위 지표를 만점 기준 비율로 변환하여 자치구별 가중평균을 산출하였다.

\begin{table*}[htbp]
  \centering\small
  \caption{UMC 차원별 측정 지표}
  \label{tab:1-umc-dims}
  \begin{tabularx}{\textwidth}{l Y Y Y}
    \toprule
    \textbf{Dimension} & \textbf{Indicators} & \textbf{Data Source} & \textbf{Note} \\
    \midrule
    Connectivity & station\_density\_4g, station\_density\_5g, download\_speed & Spectrum Resource Mgmt System; NIA Communication Quality & Per 1K daytime living pop; CQ uses 2022--2024 avg \\
    \addlinespace[2pt]
    Available for Use & mobile\_data\_usage, online\_service\_days, wifi\_density & SKT Telecom Data; Local Admin WiFi License Data & Pop-weighted district avg; WiFi per 1K total pop \\
    \addlinespace[2pt]
    Affordability & avg\_income, delinq\_rate (reverse) & Seoul Commercial Analysis Service; SKT Telecom Data & ITU IDI capacity + burden; delinq = reverse-coded \\
    \addlinespace[2pt]
    Devices & device\_home\_core, device\_use\_core & Seoul Digital Competency Survey 2023 & Home: excl. smartphone; Use: incl. smartphone \\
    \addlinespace[2pt]
    Digital Skills & Basic / Applied / Advanced / Cyber activity & Seoul Digital Competency Survey 2023 & --- \\
    \addlinespace[2pt]
    Safety & safety\_behavior, safety\_awareness & Seoul Digital Competency Survey 2023 & Security practice (7 items) + Awareness (3 items) \\
    \bottomrule
  \end{tabularx}
\end{table*}

\subsubsection{측정 지표 지수화}
\label{subsubsec:3-1-2}

모든 원시 지표는 25개 자치구 내에서 Min-Max 스케일링을 통해 $[0,1]$ 구간으로 정규화한다. 정규화된 값 $x^{*}$는 식~\eqref{eq:minmax}과 같이 산출한다.

\begin{equation}\label{eq:minmax}
  x^{*} = \frac{x_{i} - x_{\min}}{x_{\max} - x_{\min}}
\end{equation}

이때, 값이 클수록 나쁜 결과를 의미하는 역방향 지표는 변환하여 모든 지표의 방향성을 통일하였다. 그 다음, 각 차원 내에서 정규화된 지표들을 동일 가중치 산술 평균으로 결합하여 차원별 점수를 산출하였다. 동일 가중치로 분석을 진행한 이유는, 6개 차원 간 상대적 중요도에 대한 합의된 기준이 부재한 상태에서 연구자의 자의적 판단을 최소화하기 위함이다~\citep{ITU2024}. 산출된 종합지수는 0(최저)에서 1(최고) 범위를 가지며, 서울 25개 자치구 전체에 대해 2023년과 2024년 각각 독립적으로 산출한다. 해석에 있어 중요한 점은 이 종합 지수가 지역의 절대적 UMC가 아닌, 서울 내 다른 자치구 간의 상대적 위치를 보여준다는 것이다.

\subsubsection{지역별 UMC 지수 특성}
\label{subsubsec:3-1-3}

25개 자치구의 2024년 UMC 종합지수는 0.279(중랑구)에서 0.695(서초구)까지 분포하였으며, 최상위와 최하위 간 약 2.5배의 격차가 존재하였다. 서울이 높은 인터넷 보급률을 가진 도시라는 일반적 인식에도 불구하고, 하위 행정단위에서는 상당한 상대적 수준의 디지털 연결성 격차가 관찰되었다. 상위 5개 자치구는 서초구(0.695), 영등포구(0.686), 마포구(0.665), 서대문구(0.615), 용산구(0.603)이다. 하위 5개 자치구는 중랑구(0.279), 도봉구(0.303), 강북구(0.380), 구로구(0.389), 노원구(0.389)이다. 상위 자치구는 대체로 도심부 또는 강남권에 위치하여 밀집된 상업 활동, 높은 소득 수준, 풍부한 민간 디지털 인프라의 혜택을 받는 반면, 하위 자치구는 북부 외곽 및 남서부에 집중되어 있었다.

종합지수의 공간적 분포를 살펴보면(그림~\ref{fig:3-umc-comp}), UMC 지수가 높게 나온 자치구는 두 구역에 집중되어 있었다. 이는 각각 서초-영등포-마포로 이어지는 한강과 도심권 지역과, 용산-종로-중구 중심의 전통적 도심부이다. 반면, 북부 외곽(강북구, 도봉구, 노원구)과 남서부(구로구, 금천구)는 도시 평균 이하의 점수를 보였다. 이러한 공간 분포는 서울의 사회경제적인 자원 분포와 연관되어 있었다. 즉, 디지털 격차가 도시의 구조적 공간 불평등과 밀접히 연관되어 있을 가능성을 시사하였다.

차원별 분석에서 한편으로, 어떤 자치구도 6개 차원 모두에서 균일하게 우수하지는 않다는 점이 발견되었다(그림~\ref{fig:3-umc-heatmap}). 예컨대 서울 북부에 위치한 은평구의 경우, Connectivity 차원에서 0.781로 상위권에 위치하지만, Affordability 차원에서는 0.328로 하위권에 그쳤다. 유사하게, 관악구는 Skills 차원에서 0.795로 1위를 기록하였지만 Affordability에서는 0.152로 최하위권이다. 이러한 다차원적 지역간 불균등은, 종합 순위만으로는 정책적으로 의미 있는 차원별 취약성을 포착할 수 없음을 의미하였다.

\begin{keymsg}[title=핵심 발견 \textperiodcentered{} 디지털 사막 4개 자치구]
종합지수 하위 자치구(중랑\textperiodcentered{}도봉\textperiodcentered{}강북\textperiodcentered{}노원)는 Devices\textperiodcentered{}Safety 차원에서 동반 결핍을 보이며, 공간적으로도 서울 북부에 집중되어 있다. 이는 권역 단위 협력 거버넌스 구축의 정당화 근거를 제공한다.
\end{keymsg}

결론적으로, 그림~\ref{fig:3-umc-dim}에 제시된 바와 같이 공간적으로 밀집한 종합지수 하위 4개 자치구(중랑구, 도봉구, 강북구, 노원구)를 본 프로젝트에서는 잠정적으로 디지털 사막(Digital Desert)으로 분류한다. 이들 자치구는 공통적으로 Devices와 Safety 차원에서 뚜렷한 결핍을 보이며, Affordability 역시 서울시 평균에 미치지 못했다. 한편, Connectivity 차원에서는 일부 자치구(중랑구 0.630, 강북구 0.562)가 서울시 평균을 상회하여, 물리적 네트워크 인프라 자체의 부재보다는 이를 활용하기 위한 인적\textperiodcentered{}경제적 역량의 부족이 핵심 격차 요인일 가능성을 보여준다. 이러한 발견은 \ref{subsec:3-2}절의 다층 분석에서 개인 수준 변수와 자치구 수준 변수의 상대적 기여를 분리 및 검증함으로써 보다 체계적으로 검토된다.

또한, UMC의 일부 차원에서 지리적으로 인접한 지역이 서로 유사한 패턴을 보였다. 그림~\ref{fig:3-umc-dim}은 차원별 지역의 공간적 자기상관성을 나타낸 것이다. 공간적으로 군집해 있는 자치구가 서로 유사한 패턴을 보이는 점은 디지털 사막 해소를 위한 광역적 거버넌스 체계 구축의 가능성을 시사한다.

결론적으로, 그림~\ref{fig:3-umc-dim}에 제시된 바와 같이 종합지수 하위 5개 자치구(중랑구, 도봉구, 강북구, 구로구, 노원구)를 본 프로젝트에서는 잠정적으로 디지털 사막(Digital Desert)으로 분류한다. 이들 자치구는 공통적으로 Devices와 Safety 차원에서 뚜렷한 결핍을 보이며, Affordability 역시 서울시 평균에 미치지 못했다. 한편, Connectivity 차원에서는 일부 자치구(중랑구 0.630, 강북구 0.562)가 서울시 평균을 상회하여, 물리적 네트워크 인프라 자체의 부재보다는 이를 활용하기 위한 인적\textperiodcentered{}경제적 역량의 부족이 핵심 격차 요인일 가능성을 보여준다. 이러한 발견은 \ref{subsec:3-2}절의 다층 분석에서 개인 수준 변수와 자치구 수준 변수의 상대적 기여를 분리 및 검증함으로써 보다 체계적으로 검토된다. 또한, UMC의 일부 차원에서 지리적으로 인접한 지역이 서로 유사한 패턴을 보였다. 그림~\ref{fig:3-lisa-grid}은 차원별 지역의 공간적 자기상관성을 나타낸 것이다. 공간적으로 군집해 있는 자치구가 서로 유사한 패턴을 보이는 점은 디지털 사막 해소를 위한 광역적 거버넌스 체계 구축의 가능성을 시사한다.

\begin{figure}[htbp]
  \centering
  \begin{subfigure}[t]{0.55\linewidth}
    \centering
    \includegraphics[width=\linewidth]{image9.png}
    \caption{자치구별 UMC 종합점수}
    \label{fig:3-umc-comp}
  \end{subfigure}\hfill
  \begin{subfigure}[t]{0.43\linewidth}
    \centering
    \includegraphics[width=\linewidth]{image10.png}
    \caption{UMC 차원별 점수 히트맵}
    \label{fig:3-umc-heatmap}
  \end{subfigure}
  \caption{자치구별 UMC 종합점수 및 차원별 히트맵}
  \label{fig:3-umc-two}
\end{figure}

\figfull{image11.png}{UMC 차원별 점수}{fig:3-umc-dim}

\begin{figure*}[htbp]
  \centering
  \begin{subfigure}[t]{0.32\textwidth}
    \centering
    \includegraphics[width=\linewidth]{image12.jpeg}
    \caption{Connectivity}
  \end{subfigure}\hfill
  \begin{subfigure}[t]{0.32\textwidth}
    \centering
    \includegraphics[width=\linewidth]{image13.jpeg}
    \caption{Available for Use}
  \end{subfigure}\hfill
  \begin{subfigure}[t]{0.32\textwidth}
    \centering
    \includegraphics[width=\linewidth]{image14.jpeg}
    \caption{Affordability}
  \end{subfigure}

  \vspace{4pt}
  \begin{subfigure}[t]{0.32\textwidth}
    \centering
    \includegraphics[width=\linewidth]{image15.jpeg}
    \caption{Devices}
  \end{subfigure}\hfill
  \begin{subfigure}[t]{0.32\textwidth}
    \centering
    \includegraphics[width=\linewidth]{image16.jpeg}
    \caption{Skills}
  \end{subfigure}\hfill
  \begin{subfigure}[t]{0.32\textwidth}
    \centering
    \includegraphics[width=\linewidth]{image17.jpeg}
    \caption{Safety}
  \end{subfigure}
  \caption{UMC 수준의 공간적 자기상관 (서울 25개 자치구)}
  \label{fig:3-lisa-grid}
\end{figure*}

\subsection{개인 및 지역 특성과 디지털 연결성}
\label{subsec:3-2}

본 연구의 데이터는 개인(Level 1)이 자치구(Level 2) 내에 군집되어 있는 위계적 구조를 갖는다. 여기서 Level 2는 개인의 거주지에 해당한다. 이러한 구조에서 통상적인 OLS 회귀분석은 관측치의 비독립성 가정을 위반하여, 자치구 수준 예측변수의 표준오차를 과소추정하고 1종 오류를 팽창시킬 수 있다~\citep{RaudenbushBryk2002}. 따라서 분산을 25개 자치구를 기준으로 구 내와 구 간 성분으로 분해하고, 개인 수준과 자치구 수준 변수를 동시에 투입할 수 있는 2수준 위계적 선형 모형(HLM)을 채택하였다.

두 번째 분석에서 활용한 자료인 서울서베이는 2003년부터 서울시에서 2만 개 가구를 대상으로 년마다 실시해온 행정 자료이다. 자료의 수집은 매년 8\,$\sim$\,9월에 비대면과 방문면접을 병행하여 이뤄진다.

\subsubsection{변수 구성}
\label{subsubsec:3-2-1}

먼저 분석의 종속변수인 디지털 활용 점수는 서울서베이의 응답항목 중 디지털 서비스 이용빈도 항목의 8개 항목을 합산한 점수이다. 8개 항목은 생활정보 서비스, 전자상거래 서비스, 금융거래 서비스, 공공 서비스, 커뮤니케이션 서비스, 여가 서비스, 교육학습 서비스, 모바일 결제로 구성된다. 이들은 전화와 문자와 같은 기초적 연결 행위가 아니라, 디지털 기기를 매개로 한 실질적인 추가 활용에 해당하여, 디지털 환경이 발달한 서울에 적합한 지표로 판단하였다. 각 항목은 원래 4점 서열 척도를 0--10 연속 척도로 변환하여 총 0--80점 범위의 연속변수를 구성한다. 이 점수가 높을수록 개인의 디지털 서비스 활용이 다양하고 원하는 서비스를 자유롭게 이용할 수 있음을 의미한다.

본 분석의 독립변수는 개인 특성의 변수와 지역 특성의 변수로 구분된다(표~\ref{tab:2-hlm-vars}). 개인 특성의 변수는 개인의 디지털 활용에 영향을 미칠 수 있는 변수인 성별, 연령, 직업, 학력, 소득수준, 장애여부로 구성된다. 교육수준은 중졸 이하, 고졸, 대졸(기준범주), 대학원 이상의 4범주로, 직업은 관리\textperiodcentered{}전문직, 사무직(기준범주), 생산\textperiodcentered{}기능직, 기타(학생\textperiodcentered{}주부\textperiodcentered{}무직)의 4범주로 구분하였다. 주거유형은 아파트(기준범주), 단독주택, 다세대 및 다가구, 기타의 4범주로 분류하였다. 연령과 소득은 총평균 중심화를 적용하여 절편의 해석 가능성을 높였다. 추가로, 1인가구 더미, 장애 여부 더미, 연도 더미(기준: 2023년)를 포함하였다.

지역 수준(Level 2) 변수는 \ref{subsec:3-1}절에서 구축한 UMC 차원 점수 중 지역의 디지털 환경을 결정하는 외생적 요인의 세 차원 Connectivity, Available for Use, Affordability을 활용하였다. 이는, Devices, Digital Skills, Safety가 서울디지털역량실태조사의 개인 응답을 자치구 단위로 집계한 것을 고려한 것이다. 그로 인한 내생성 문제를 피하기 위해 주 분석에서 세 가지 차원을 제외하여 지역 차원의 보편적인 의미있는 연결성에 더 적합하다고 판단되는 차원을 선택하였다. 통제변수로 자치구 평균소득과 고령화율을 포함하였으며, 모든 Level-2 변수는 앞선 \ref{subsec:3-1}절에서와 같이 Min-Max 정규화를 적용하여 $[0,1]$ 범위로 표준화하였다.

분석은 4단계로 진행된다. 무조건 모형에서 급내상관계수를 산출하여 다층 분석의 정당성을 확인한 뒤, 개인 수준 변수를 투입한 M1, 자치구 수준 변수를 추가한 M2, 교차수준 상호작용을 포함한 M3를 순차적으로 추정한다. 교차수준 상호작용은 UMC 인프라 변수(Connectivity, Available for Use, Affordability)와 취약계층 지표의 곱으로 구성하여, 자치구의 디지털 인프라 수준이 취약계층의 디지털 활용에 차별적으로 작용하는지를 검증한다.

\subsubsection{분석 결과}
\label{subsubsec:3-2-2}

분석에 활용된 표본은 서울서베이 2023\,$\sim$\,2024년 응답자 총 10,000명이며, 25개 자치구에서 평균적으로 400명이 표집되었다. 주요 변수에 대한 기술통계량을 나타낸 것은 표~\ref{tab:3-descstat}과 같다. 개인의 디지털 연결성의 평균이 80점 만점에 60점으로 나타났는데, 서울시민의 높은 디지털 연결성을 잘 보여준다. 종속변수가 음의 왜도가 존재할 수 있으나 기관 표준 척도와 10000명에 근접한 표본을 고려해 척도 선형화한 원점수대로 분석을 진행하였다.

표~\ref{tab:4-hlm-results}은 순차적으로 다층 모형의 추정 결과를 나타낸 것이다. 무조건 모형의 급내상관계수는 0.49\%로 자치구 간 분산 비율 자체는 크지 않았으나, 자치구당 평균 표본 크기 400명을 감안한 설계효과가 2.94로 임계값 2.0을 초과하여 다수준 모형의 사용에 대한 타당성을 확인하였다~\citep{HoxEtAl2018}. 그러나, 25개의 자치구의 소규모 표본 / 자치구가 설명하는 분산이 매우 적은 점을 종합하면 개인의 디지털 연결성에 지역의 연결 관련 특성이 미치는 영향이 극히 적을 것으로 추정된다. 이러한 결과는 서울의 우수한 디지털 연결성과 높은 수준의 디지털 역량으로 지역 수준의 연결성 인프라가 영향을 미치지 못한 것일 수 있다. 또한, 서울서베이는 개인의 거주 지역을 조사하기 때문에, 실제로 연결성에 직접적으로 영향을 받는 지역과 일치하지 않을 수 있다.

개인 특성 중 교육 수준의 관계가 가장 두드러졌다. 중졸 이하 학력은 대졸 대비 디지털 활용 점수를 약 22.6점 낮게 나타났다. 이는 80점 만점 기준 28\%에 해당하는 상당한 격차다. 고졸 역시 대졸보다 7.3점 낮게 나타난 반면, 대학원 이상은 오히려 2.9점 높아 교육 수준과 디지털 활용 능력 간의 뚜렷한 정적 관계를 확인할 수 있다. 연령 또한 유의한 영향을 미쳤다. 평균 연령으로부터 1세 증가할수록 개인의 디지털 연결성에 관한 점수가 0.35점 감소하였으며, 65세 이상 고령층의 경우 연령 효과와는 별개로 추가적인 6.0점 하락이 확인되었다. 65세 이상의 연령층에서 디지털 연결성의 두드러진 저하가 나타났다. 직업 측면에서는 관리직\textperiodcentered{}전문직 종사자가 사무직 대비 2.3점 높게 나타나, 고소득 직종일수록 디지털 활용 수준도 높은 경향을 보였다. 주거환경의 경우 단독주택 거주자는 아파트 대비 2.5점, 다세대주택 거주자는 1.8점 낮았는데, 이는

\begin{table*}[htbp]
  \centering\small
  \caption{분석 변수 설명}
  \label{tab:2-hlm-vars}
  \begin{tabularx}{\textwidth}{l Y c l l}
    \toprule
    \textbf{Variable} & \textbf{Description} & \textbf{Level} & \textbf{Scale} & \textbf{Source} \\
    \midrule
    digital\_use\_score & Digital service utilization score (8 items summed) & L1 & 0--80 (continuous) & Seoul Survey 2023--24 \\
    \addlinespace[2pt]
    age\_c & Age (grand-mean centered, M = 49.3) & L1 & years & Seoul Survey \\
    \addlinespace[2pt]
    female & Female (ref: male) & L1 & 0/1 & Seoul Survey \\
    \addlinespace[2pt]
    Edu & Education: $\leq$middle / high / grad+ (ref: college) & L1 & 0/1 each & Seoul Survey \\
    \addlinespace[2pt]
    Income & Monthly income (grand-mean centered) & L1 & 10K KRW & Seoul Survey \\
    \addlinespace[2pt]
    job & Job: professional / blue-collar / other (ref: white-collar) & L1 & 0/1 each & Seoul Survey \\
    \addlinespace[2pt]
    disabled\_bin & Disability status & L1 & 0/1 & Seoul Survey \\
    \addlinespace[2pt]
    hh\_size\_1 & Single-person household & L1 & 0/1 & Seoul Survey \\
    \addlinespace[2pt]
    housing\_detached / \_multi / \_other & Housing: detached / multi-family / other (ref: apt) & L1 & 0/1 each & Seoul Survey \\
    \addlinespace[2pt]
    year & Year 2024 dummy (ref: 2023) & L1 & 0/1 & Seoul Survey \\
    \addlinespace[2pt]
    elderly & Elderly ($\geq$65) & L1 & 0/1 & Seoul Survey \\
    \addlinespace[2pt]
    low\_edu & Low education ($\leq$middle school) & L1 & 0/1 & Seoul Survey \\
    \addlinespace[2pt]
    connectivity & Connectivity dimension score (normalized) & L2 & 0--1 & UMC Part 1 \\
    \addlinespace[2pt]
    wifi\_total\_n & Public WiFi density per 1K pop (normalized) & L2 & 0--1 & Local Admin License \\
    \bottomrule
  \end{tabularx}
\end{table*}

주택 유형에 따른 소득 수준의 차이뿐 아니라 건물 유형별 인터넷 연결 여건의 차이도 반영된 결과로 해석된다. 1인 가구 역시 다인 가구 대비 1.8점 낮은 점수를 보였다. 그 외 개인 특성으로는, 여성이 남성 대비 2.6점 높게 나타났다. 여성이 남성에 비해 더 높은 디지털 연결성을 갖는 것은, 국내 다양한 디지털 관련 서비스에 여성이 더 민감하게 반응한 결과일 수 있다. \textbf{[인용 필요]} 이러한 결과는 국가에 따른 사회문화적인 맥락을 고려할 필요가 있다. 장애인의 경우 3.1점 높게 나타났는데, 표본의 특성에 기인한 것으로 보인다.

자치구 수준의 UMC 인프라 차원은 주효과만으로는 대체로 유의하지 않았다. Model 2에서 세 차원 모두 통계적 유의성에 도달하지 못하였고, 우도비 검정 결과도 Model 1 대비 유의한 개선이 없었다. 그러나 교차수준 상호작용을 투입한 Model 3에서는 Infrastructure 차원이 유의한 부정적 효과를 보였다. 이는 인프라의 평균적 효과가 아닌, 집단별 차별적 효과가 존재함을 의미한다.

본 분석의 핵심 발견은 교차수준 상호작용에서 나타났다. 공공 WiFi 밀도와 고령 더미의 상호작용은 3.03점의 유의한 정적 효과를 보였다. 이는 공공 WiFi 인프라가 비고령자보다 고령자에게 상대적으로 더 큰 혜택을 제공함을 뜻한다. 디지털 인프라 수준이 높은 자치구에 거주하는 저학력 주민은 인프라 수준이 낮은 곳의 저학력 주민 대비 현저히 높은 디지털 활용 점수를 기록하였다. 이 결과는 인프라 투자가 교육 취약계층의 디지털 격차를 완화하는 데 불균형적으로 큰 효과를 가질 수 있음을 시사한다.

개인 수준 상호작용에서는 이중 취약성이 확인되었다. 독거 고령자는 비독거 고령자 대비 추가적으로 4.0점이 낮았으며, 고령일수록 저학력의 부정적인 효과가 더 커지는 것으로 나타났다. 이러한 결과는 단일 취약 요인보다 복합 취약 요인이 디지털 격차를 심화시킴을 보여준다.

\begin{table*}[htbp]
  \centering\small
  \caption{주요 변수의 기술통계량}
  \label{tab:3-descstat}
  \begin{tabularx}{\textwidth}{Y r r r r r}
    \toprule
    \textbf{변수(Variable)} & \textbf{N} & \textbf{평균(Mean)} & \textbf{표준편차(SD)} & \textbf{최소(Min)} & \textbf{최대(Max)} \\
    \midrule
    디지털 활용 점수(Digital use score) & 10,000 & 60.88 & 16.84 & 0 & 80 \\
    연령(Age) & 10,000 & 49.32 & 16.23 & 15 & 99 \\
    여성 비율(Female, proportion) & 10,000 & 0.51 & --- & 0 & 1 \\
    학력: 중졸 이하(Education: $\leq$Middle school) & 10,000 & 0.07 & --- & 0 & 1 \\
    학력: 고졸(Education: High school) & 10,000 & 0.25 & --- & 0 & 1 \\
    학력: 대학원 이상(Education: Graduate+) & 10,000 & 0.13 & --- & 0 & 1 \\
    장애 여부(Disability) & 10,000 & 0.03 & --- & 0 & 1 \\
    월평균 소득(만원 단위, Monthly income [10K KRW]) & 10,000 & 237.55 & 209.73 & 0 & 1,050 \\
    1인 가구(Single-person household) & 10,000 & 0.17 & --- & 0 & 1 \\
    고령(65세 이상, Elderly 65+) & 10,000 & 0.22 & --- & 0 & 1 \\
    (군집) 연결성 지수(Connectivity, norm.) & 50 & 0.5 & 0.18 & 0.18 & 0.98 \\
    이용 가능성(Available for Use, norm.) & 50 & 0.5 & 0.19 & 0.12 & 0.96 \\
    비용 접근성(Affordability, norm.) & 50 & 0.5 & 0.22 & 0.07 & 1 \\
    \bottomrule
    \multicolumn{6}{p{0.97\textwidth}}{\footnotesize\textit{Note.} N = 10,000명(2023\,$\sim$\,2024년, 25개 자치구, 표본 합산). 소득은 만원 단위. 자치구(군집) 단위 변수는 0\,$\sim$\,1(min-max) 정규화 지수임. N = 10,000 individuals nested within $j$ = 25 districts (2023--2024 pooled). Income measured in 10,000 KRW units. District-level variables are min-max normalized (0--1 scale).}
  \end{tabularx}
\end{table*}

\begin{table*}[htbp]
  \centering\small
  \caption{HLM 추정 결과}
  \label{tab:4-hlm-results}
  \begin{tabularx}{\textwidth}{l Y Y Y}
    \toprule
    \textbf{변수} & \textbf{Model 1} & \textbf{Model 2} & \textbf{Model 3} \\
    \midrule
    \multicolumn{4}{l}{\textit{Level 1: 개인 특성}} \\
    \addlinespace[2pt]
    연령(중심화) & $-0.347^{***}$ (0.012) & $-0.347^{***}$ (0.012) & $-0.278^{***}$ (0.014) \\
    여성 & $2.636^{***}$ (0.305) & $2.634^{***}$ (0.305) & $2.601^{***}$ (0.305) \\
    학력: 중학교 이하 & $-22.618^{***}$ (0.770) & $-22.610^{***}$ (0.770) & $-19.019^{***}$ (2.379) \\
    학력: 고등학교 & $-7.289^{***}$ (0.380) & $-7.282^{***}$ (0.380) & $-7.258^{***}$ (0.380) \\
    학력: 대학원 이상 & $2.906^{***}$ (0.483) & $2.909^{***}$ (0.483) & $2.912^{***}$ (0.481) \\
    소득(중심화) & $0.003^{**}$ (0.001) & $0.003^{**}$ (0.001) & $0.002^{**}$ (0.001) \\
    직업: 전문직 & $2.329^{***}$ (0.451) & $2.326^{***}$ (0.451) & $2.28^{***}$ (0.449) \\
    직업: 생산직 & $-0.373$ (0.424) & $-0.379$ (0.424) & $-0.385$ (0.423) \\
    직업: 기타 & $-1.233^{**}$ (0.456) & $-1.226^{**}$ (0.456) & $-1.105^{*}$ (0.454) \\
    장애 & $3.084^{***}$ (0.888) & $3.081^{***}$ (0.888) & $3.006^{***}$ (0.884) \\
    1인 가구 & $-1.775^{***}$ (0.338) & $-1.776^{***}$ (0.338) & $-0.877^{*}$ (0.367) \\
    단독주택 & $-2.478^{***}$ (0.377) & $-2.480^{***}$ (0.377) & $-2.574^{***}$ (0.379) \\
    다세대주택 & $-1.798^{***}$ (0.408) & $-1.795^{***}$ (0.408) & $-1.933^{***}$ (0.408) \\
    연도(2024) & $0.294$ (0.283) & $0.292$ (0.283) & $0.292$ (0.283) \\
    고령자(65세 이상) & $-6.046^{***}$ (1.696) & & \\
    \addlinespace[4pt]
    \multicolumn{4}{l}{\textit{Level 2: 자치구 특성}} \\
    \addlinespace[2pt]
    연결성 & $-1.127$ (0.747) & $-1.688^{*}$ (0.804) & \\
    이용 가능성(WiFi) & $0.085$ (0.780) & $-0.200$ (0.829) & \\
    비용 접근성 & $-1.578$ (1.135) & & \\
    자치구 소득 & $-2.035^{\dagger}$ (1.181) & & \\
    고령화율 & $-0.298$ (1.095) & $-0.293$ (1.111) & \\
    \addlinespace[4pt]
    \multicolumn{4}{l}{\textit{상호작용}} \\
    \addlinespace[2pt]
    고령자 $\times$ 1인 가구 & $-3.958^{***}$ (0.936) & & \\
    연령 $\times$ 저학력 & $-0.277^{***}$ (0.080) & & \\
    \addlinespace[4pt]
    \multicolumn{4}{l}{\textit{교차수준 상호작용}} \\
    \addlinespace[2pt]
    연결성 $\times$ 고령자 & $-1.920$ (1.901) & & \\
    WiFi $\times$ 고령자 & $3.026^{*}$ (1.515) & & \\
    자치구 소득 $\times$ 고령자 & $2.358$ (1.532) & & \\
    연결성 $\times$ 저학력 & $10.936^{***}$ (3.370) & & \\
    \bottomrule
    \multicolumn{4}{p{0.97\textwidth}}{\footnotesize\textit{Note.} Unstandardized coefficients with standard errors in parentheses. Model 0 = null model (intercept only). Reference categories: male, college education, white-collar job, apartment, multi-person household. $^{\dagger}p<.10$, $^{*}p<.05$, $^{**}p<.01$, $^{***}p<.001$.}
  \end{tabularx}
\end{table*}

\subsection{플랫폼 텍스트 분석: LLM 기반 디지털 연결성 사례분석}
\label{subsec:3-3}

앞선 \ref{subsec:3-1}절과 \ref{subsec:3-2}절에서는 행정 데이터와 설문 조사를 이용해 서울의 디지털 격차를 정량적으로 측정하고, 개인 특성과 지역 인프라의 역할을 분리하였다. 그러나 정량적 지표만으로는 주민들이 일상에서 언제, 어떤 맥락에서 디지털 연결의 결여를 경험하는지까지 포착하기 어렵다. 디지털 연결성에 대한 생활 세계의 체감과 구조적 지표 사이의 거리를 좁히기 위해 정성적 증거를 통합할 필요가 있다.

본 절의 목적은 지역 기반 플랫폼의 비정형 텍스트를 이용해 UMC의 6개 차원(Connectivity, Availability for Use, Affordability, Devices, Skills, Safety)에 잠재된 디지털 결여 사례를 체계적으로 추출하고, 이러한 사례를 \ref{subsec:3-1}절과 \ref{subsec:3-2}절의 정량 분석과 결합하는 방법론적 틀을 제시하는 데 있다. 플랫폼 자체를 분석의 대상으로 삼기보다는, 플랫폼이 제공하는 지역 단위의 생활 경험 데이터를 통해 구조와 경험 사이의 연결을 복원하고자 한다.

\subsubsection{당근 플랫폼 특성}
\label{subsubsec:3-3-1}

비정형 텍스트의 출처는 한국의 지역 기반 커뮤니티 플랫폼 '당근'의 동네생활 게시글이다. 당근을 선정한 이유는 세 가지다. 첫째, 각 게시글에 행정동과 자치구 수준의 지리적 메타데이터가 자동으로 부여되어 \ref{subsec:3-1}절에서 구축한 UMC 지수와 직접적으로 연동할 수 있다. 둘째, 플랫폼 이용자가 서울의 25개 자치구 전역에 골고루 분포해 있어 특정 지역이나 연령층에 치우치지 않는다. 셋째, 게시글은 일상생활 속 디지털 문제를 자발적으로 기록하기 때문에 조사자가 설문에서 미처 포착하지 못한 생활상의 결여를 드러낼 수 있다. 연구의 초점은 플랫폼 자체에 대한 평가가 아니라, 이와 같은 구조적 장점을 활용해 생활 세계의 경험을 수집하는 데 있다.

\subsubsection{왜 LLM인가: 기존 방법론의 한계와 LLM의 구조적 이점}
\label{subsubsec:3-3-2}

UMC의 여섯 차원에 따라 비정형 텍스트를 분류하려면 두 가지 요구를 충족해야 한다. 첫째, 게시글이 디지털 연결성 문제와 관련된 지 여부와 관련된 경우 해당 차원을 정확히 판별해야 한다. 둘째, 일상 언어로 쓰인 구어체 한국어는 조사나 어미 변화가 복잡한 교착어적 특성으로 인해 단순한 형태소 분석기만으로는 정확히 처리하기 어렵다. 기존의 토픽 모델링이나 키워드 기반 분류는 잠재 주제를 파악하는 데 유용하지만 사전 정의된 차원과 일치하지 않거나, 대화체 텍스트를 제대로 분석하지 못하는 한계가 있다. 이러한 한계는 한국어 형태소 분석기가 구어체 데이터에서 정확도가 낮고, 데이터 훈련이 필요한 머신러닝 모델이 충분한 라벨을 요구한다는 점에서 더욱 두드러진다.

대규모 언어 모델(LLM)은 이런 제약을 상당 부분 극복한다. 다국어로 사전 학습되어 한국어의 형태소 분석이 없어도 맥락을 파악할 수 있으며, 104개 이상의 언어를 동시에 다루는 모델은 언어간 의미론적 표현을 공유해 적은 양의 예시만으로도 새로운 태스크에 적응하는 능력을 갖추고 있다. LLM을 도구로 사용할 때는 출력의 확률성과 편차를 고려해야 한다. 동일한 입력에 대해서도 결과가 달라질 수 있으므로, 본 연구는 분석 절차를 고정하고 LLM은 분류와 추론을 수행하는 계산 모듈로만 활용했다. 이를 통해 연구자가 정의한 분류 규칙과 추론 프로토콜을 일관되게 적용하면서도 LLM의 문맥 이해력을 활용할 수 있도록 설계하였다.

\subsubsection{분석 구조}
\label{subsubsec:3-3-3}

텍스트 분석은 네 단계로 구성된다. 먼저, 크롤링한 원시 CSV 데이터에서 삭제되어 접근할 수 없는 게시글과 광고성\textperiodcentered{}중복 게시글을 제거한 뒤, 여섯 차원별 키워드 사전을 이용해 UMC 후보 게시글을 추출한다. 키워드 사전은 디지털 연결성과 직간접적으로 관련된 단어를 정리한 것으로 해당 단어를 포함하는 게시글만 1차로 추출하였다. 이후, 게시글이 디지털 연결성 문제와 관련이 있는지 여부를 판정하는 1차 분류 (umc\_related = Y/N/?)를 수행하며, 판정의 범위를 넓히고자 LLM에게 예시, 제한을 상세하게 밝히지 않고 분석하였다. 그 다음, 관련성이 확인된 게시글만을 대상으로 사용자가 제공한 상세한 예시를 포함한 few-shot 프롬프트를 LLM에게 제시해 게시글을 6개 개념적 차원의 사례로 분류한다. 이 단계에서는 모델의 자율성을 제한하여 프롬프트에 예시를 포함한 구체적인 분류 기준과 명세만을 제공했고, 모델이 판단에 참조하는 정보 공간인 컨텍스트를 엄격하게 관리하였다.

표~\ref{tab:5-text-sample}는 분석 구조와 표본 추출 과정을 요약한 것이다. 원시 데이터 1,287,761건 중 약 51.5\%는 게시자 탈퇴나 차단으로 본문을 볼 수 없어 분석에서 제외되었다. 중복\textperiodcentered{}광고성 게시글 제거 후에는 약 48만 건이 남았으며, 여섯 차원별 키워드 사전을 적용한 결과 26,742건이 UMC 관련 후보로 추출되었다. few-shot 프롬프트를 이용한 1차 판정에서는 약 57.7\%(15,438건)가 관련(Y)으로, 29.6\%(7,921건)가 비관련(N)으로, 12.6\%(3,383건)가 판단 유보(?)로 분류되었다. 이 가운데 Y로 판정된 게시글만을 대상으로 차원 분류를 진행했다.

차원 분류 판정에서 LLM에게 각 차원에서 대표적인 사례를 담은 3\,$\sim$\,4개의 예시를 프롬프트로 제공했다. 이러한 few-shot 프롬프트는 모델이 사람의 판단 기준을 학습하는 데 도움이 되며, 복잡한 작업에서 zero-shot보다 나은 성능을 보인다~\textbf{[인용 필요]}. 예시에는 어떤 유형의 게시글이 디지털 연결성 문제와 무관한지를 보여주는 음성 예시도 포함해, 관련성을 최대한 좁혀 높은 정밀도를 확보하도록 했다. 2차 차원 분류에서는 구체적인 분류 기준과 정의만을 프롬프트로 제공해 모델이 보다 자율적으로 사례를 분류하도록 설계했다. 이처럼 앞 단계는 예시 기반 few-shot 방식이고, 뒤 단계는 규칙 기반 프롬프트를 사용한다는 점에서 추론과 few-shot 사이의 균형을 보여준다(그림~\ref{fig:3-pipeline-roles} 참조).

\begin{table}[htbp]
  \centering\small
  \caption{텍스트 분석 표본 추출 및 판정 절차}
  \label{tab:5-text-sample}
  \begin{tabularx}{\linewidth}{l Y r}
    \toprule
    \textbf{단계} & \textbf{설명} & \textbf{남은 게시글 수} \\
    \midrule
    원시 데이터 & 2024년 5월\,$\sim$\,2026년 3월에 작성된 동네생활 게시글 전수 & 1,287,761 \\
    접근가능 게시글 & 탈퇴\textperiodcentered{}차단으로 본문 확인이 불가능한 게시글 제거 & 624,564 \\
    중복\textperiodcentered{}광고 제거 & 중복된 내용, 광고\textperiodcentered{}홍보 게시글 제거 & 480,972 \\
    UMC 키워드 필터 & 6개 차원별 키워드 사전을 적용해 후보 추출 & 26,742 \\
    1차 판정 & few-shot 프롬프트로 UMC 관련성 분류 & --- \\
    \bottomrule
  \end{tabularx}
\end{table}

\fig{image18.png}{분류와 추론의 역할 구분}{fig:3-pipeline-roles}

이후, 분류 결과를 자치구 단위로 집계하여 표출 프로파일을 만들고, \ref{subsec:3-1}절의 정량 지표와 결합하기 위해 베이지안 업데이트를 적용한다. 정량 지표와 텍스트 분류 결과를 결합하기 위해 베이지안 업데이트를 적용하였다. 사전 분포는 UMC의 6가지 개념의 잠재적 공간이다. 차원 $d$에서 자치구 $j$의 UMC 실측값을 $s_{jd}$로 할 때, 사전 분포는 $\mathrm{Beta}(\alpha_{jd}, \beta_{jd})$이며, 평균이 $s_{jd}$가 되도록 사전 분포를 가정한다. 강도 모수 $\kappa$를 사용해 $\alpha_{jd} = s_{jd}\cdot\kappa$, $\beta_{jd} = (1-s_{jd})\cdot\kappa$로 표현한다. 텍스트 분류에서 해당 차원으로 판정된 게시글 수를 $y_{jd}$, 총 분류된 게시글 수를 $n_{jd}$라 할 때, 우도는 $\mathrm{Binomial}(n_{jd}, \theta_{jd})$로 표현된다. 사후 분포는 $\theta_{jd}\mid y_{jd} \sim \mathrm{Beta}(\alpha_{jd}+y_{jd}, \beta_{jd}+n_{jd}-y_{jd})$로 계산된다. 사전 강도 $\kappa$는 기본값 20을 사용하고 10과 30에 대한 민감도 분석도 수행하였다. 게시글 수가 적은 자치구에서 사후 분포가 사전 분포에 의해 지나치게 지배되는 것을 막기 위해, 게시글 수 $n_{jd}$를 생활인구로 보정한 $n_{jd}^{*} = n_{jd}\times(\text{생활인구}_{j}/\overline{\text{생활인구}})$로 사용하였다.

\begin{align}
\alpha_{jd} &= \kappa \times s_{jd} \label{eq:alpha} \\
\beta_{jd}  &= \kappa \times (1 - s_{jd}) \label{eq:beta} \\
\theta_{jd}\mid y_{jd} &\sim \mathrm{Beta}\bigl(\alpha_{jd}+y_{jd},\, \beta_{jd}+n_{jd}-y_{jd}\bigr) \label{eq:posterior} \\
n_{jd}^{*} &= n_{jd}\times\bigl(\text{생활인구}_{j}/\overline{\text{생활인구}}\bigr) \label{eq:popadj}
\end{align}

베이지안 집계에서 게시글로 추정된 사후 분포가 사전 분포로 유의미하게 이동한 자치구\textperiodcentered{}차원 조합을 대상으로 인과기제에 관한 정성적 추론을 수행하였다. 구체적으로는 (1) 사후 이동이 통계적으로 유의하며, (2) 해당 자치구\textperiodcentered{}차원에서 Y 판정 게시글이 최소 5건 이상이고, (3) \ref{subsec:3-2}절의 다층모형 결과가 허용하는 범위 내에서만 해석을 실시하였다. 사전 분포에 의해 도출된 조합에 대해 대표 게시글 3\,$\sim$\,5건을 선정하고, 1\,$\sim$\,2문장으로 구조적 지표와 생활 경험 사이의 연결 경로를 서술하였다. 분류와 추론 작업은 세 명의 독립된 에이전트가 서로 다른 추론 경로를 통해 수행했으며, 최종 판정은 네 번째 에이전트가 세 결과를 비교\textperiodcentered{}검증하는 방식으로 이루어졌다. 이는 한 가지 프로토콜에 의존하지 않고 여러 관점의 상호 검증을 통해 편향을 통제하는 삼각검증(triangulation) 접근을 따른 것이다.

모든 분류와 추론 작업은 세 명의 독립된 분석 에이전트가 서로 다른 추론 방식 프롬프트와 컨텍스트를 사용해 수행하고, 네 번째 분석자가 세 결과를 비교\textperiodcentered{}검증하는 방식으로 이루어졌다. 이는 질적 연구에서 여러 자료원과 분석가를 동원해 편향을 줄이는 삼각검증(triangulation) 접근에 기반한다. 세 관점이 일치할 때만 해당 게시글을 최종 분석에 포함했으며, 해석상의 불확실성은 결과 표기에 반영하였다. 이러한 과정을 정리한 결과는 표~\ref{tab:6-units}, 그림~\ref{fig:3-pipeline-full}과 같다.

\begin{table}[htbp]
  \centering\small
  \caption{\ref{subsec:3-3}절의 분석 단위와 해석 단위}
  \label{tab:6-units}
  \begin{tabularx}{\linewidth}{l l Y Y}
    \toprule
    \textbf{구분} & \textbf{단위} & \textbf{무엇을 보여주는가} & \textbf{해석상 주의점} \\
    \midrule
    관찰 단위 & 개별 게시글 & 디지털 연결성 경험의 단서 & 단일 게시글만으로 지역 구조를 주장할 수 없음 \\
    패턴 단위 & 행정동 게시글 & 차원별 표출 비율과 분포 & 집계 기준과 표본 편향을 함께 검토해야 함 \\
    구조 단위 & UMC 지역 지표 & 지역 수준의 연결성 조건 & 텍스트 집계 결과와의 정합을 통해 제한적으로 해석 \\
    인과기제 단위 & 추론 메커니즘 & 사례와 지역 패턴의 연결 설명 & \ref{subsec:3-2}의 정량 결과가 허용하는 범위 안에서만 해석 \\
    \bottomrule
  \end{tabularx}
\end{table}

\subsubsection{데이터와 적용 범위}
\label{subsubsec:3-3-4}

원시 데이터는 서울 25개 자치구를 대상으로 2024년 5월부터 2026년 3월까지 작성된 동네생활 게시글 전수를 크롤링한 것으로, 중복 제거 후 총 1,287,761건이다. 이 중 삭제(DELETED) 663,163건(51.5\%)과 차단(BLOCKED) 103,728건(8.1\%)을 합한 766,891건(59.6\%)은 본문 확인이 불가능하여 분석에서 제외하였다. 분석 대상은 최종 약 52만 건이다.

표본 탈락의 체계적 편향 가능성을 검토한 결과, 자치구별 제거율의 표준편차는 4.6\%로 자치구 간 편차가 크지 않았으며, 디지털 결여 추론의 주요 대상인 생활/편의(47.0\%), 동네사건사고(48.4\%), 병원/약국(35.5\%) 카테고리는 사교\textperiodcentered{}홍보 목적 카테고리(동네친구 86.6\%, 교육 69.0\%)에 비해 제거율이 낮았다. 삭제의 상당 부분은 사용자 탈퇴에 수반된 일괄 삭제에 기인하는 것으로 판단된다. 다만, 삭제 게시글 중 일부가 디지털 결여 관련 사례를 포함할 가능성을 완전히 배제할 수는 없고, 삭제 행위 자체가 디지털 결여와 체계적으로 연관되었을 수 있다. 이는 본 분석의 한계에 해당한다.

당근 데이터의 구조적 한계도 인식할 필요가 있다. 플랫폼 이용자 자체가 디지털 접근이 가능한 집단으로 한정되므로, 디지털 소외 계층의 경험은 구조적으로 과소 대표된다. 게시글의 시간적\textperiodcentered{}공간적 분포가 균일하지 않아 자치구 간 비교에서 표본 편향의 가능성이 존재한다. 이러한 한계는 데이터 소스의 제약이지 방법론의 제약이 아니다. 보다 직접적인 질적 자료가 존재하는 맥락에서는 동일한 파이프라인을 적용하여 보다 정밀한 분석이 가능할 것이다.

\figfull{image19.png}{LLM 기반 텍스트 분석 파이프라인}{fig:3-pipeline-full}

\subsubsection{분석 결과}
\label{subsubsec:3-3-5}

텍스트 분류 결과는 표~\ref{tab:6-units}와 같다. 연결 품질(Connection Quality) 차원은 6개 차원 가운데 전이의 폭이 가장 작다($-0.006 \sim +0.005$). 실측과 텍스트 증거가 대체로 수렴하는 차원으로 볼 수 있다. 양천구는 실측 UMC 0.278로 낮은 편이나 관련 게시글 비율이 상대적으로 높아 $+0.005$의 전이를 보였고, 용산구는 실측 0.884로 높은 반면 관련 게시글이 드물어 $-0.006$의 전이를 보였다.

이용 가능성(Availability for Use) 차원에서는 도봉구가 실측 최하위(0.213)에 해당하면서도 관련 게시글 비율이 높아 $+0.006$의 전이를 보였다. 광진구는 실측 0.832로 높고 관련 게시글이 드물어 $-0.010$의 전이를 기록하였다.

경제적 접근성(Affordability) 차원은 서울특별시 자치구의 소득 수준과 일관된 방향의 경향을 보였다. 중랑구와 강서구는 각각 $+0.005$의 양의 전이를 보인 반면, 서초구는 실측이 0.993으로 최고 수준이면서 관련 게시글이 드물어 $-0.007$의 전이를 나타냈다. 이 차원은 통신비와 기기 구입 비용에 관한 고민이 소득 수준에 따라 차등적으로 표출되는 경향과 정합적이다.

기기 접근성(Devices) 차원은 자치구 간 편차가 가장 큰 차원이다(gap = 0.090). 중랑구는 실측 0.052에 해당하면서도 기기 관련 게시글 비율이 4.95\%로 전 자치구 가운데 가장 높아 $+0.048$의 전이를 보였다. 송파구 또한 $+0.026$의 유사한 패턴을 나타냈다. 반면 용산구($-0.042$)와 동작구($-0.033$)는 실측이 높으면서 관련 게시글이 드물게 나타났다. 다만 이 차원에서 실측 지표는 기기 보유율을 포착하는 반면, 텍스트 증거는 중고 거래, 수리 요청 등 기기와 관련한 일상적 거래와 고충을 포착하므로, 두 지표가 지시하는 대상이 부분적으로 상이할 가능성을 배제하기 어렵다.

디지털 역량(Digital Skills) 차원은 기기 차원에 이어 두 번째로 편차가 큰 차원이다(gap = 0.092). 중랑구는 실측 0.242에 관련 게시글 비율 4.8\%로 $+0.047$의 전이를 보였다. 강남구 또한 실측이 0.276으로 낮은 편이면서 관련 게시글 비율이 상대적으로 높아 $+0.013$의 전이를 나타냈다. 반면 관악구는 실측 0.795로 최고 수준이면서 관련 게시글이 드물어 $-0.045$의 전이를 보였다.

디지털 역량(Digital Skills) 차원에서는 기기 차원에 이어 두 번째로 큰 편차가 나타났다. 중랑구는 실측 0.242에 비해 관련 게시글 비율 4.8\%로 $+0.047$의 이동을 보였으며, 강서구는 실측값이 높음에도 관련 게시글이 적어 $-0.031$의 이동을 보였다. 안전\textperiodcentered{}보안(Safety) 차원은 피싱 및 보이스피싱 사례가 게시글의 주요 내용을 이루었고, 중랑구($+0.018$)와 도봉구($+0.015$)에서 양의 이동이 확인되었다. 종로구는 실측값이 매우 높고 관련 게시글이 적어 사후 분포가 낮아졌다.

안전\textperiodcentered{}보안(Safety \& Security) 차원에서는 중랑구($+0.018$)와 도봉구($+0.015$)에서 피싱 사기 및 보이스피싱 관련 게시글의 비율이 높았다. 종로구는 실측 0.984로 최고 수준이면서 관련 게시글이 드물어 $-0.019$의 전이를 나타냈는데, 관광\textperiodcentered{}상업 중심지의 주거 구성과 당근 이용자 집단의 특성이 상이할 가능성을 함께 고려할 필요가 있다.

자치구 단위로 집계했을 때, 지역의 실측값과 게시글 분포 간 차이는 크게 세 가지 유형. 첫째, 평균 전이가 양의 방향으로 나타나는 자치구이다. 중랑구는 6개 차원 중 4개 차원에서 양의 전이가 가장 크게 나타나 평균 $+0.020$으로 가장 높은 값을 보였다. 도봉구와 강남구 또한 평균적으로 양의 방향을 보였으나 그 크기는 0.001 수준으로 작았다. 특히 강남구의 경우 평균 소득이 높은 자치구임에도 이용 가능성과 디지털 역량 차원에서 양의 전이가 관찰되었는데, 자치구 내부의 이질적 집단이 관련 이슈를 표출하고 있을 가능성을 시사한다. 둘째, 전이가 0에 가까운 자치구이다. 강북구($-0.000$), 강동구($-0.001$), 동대문구($-0.001$)가 이에 해당한다. 이들 자치구에서는 차원별로 양\textperiodcentered{}음의 전이가 섞여 있어 평균적으로 상쇄되는 양상을 보인다. 셋째, 평균 전이가 음의 방향으로 나타나는 자치구이다. 마포구($-0.014$), 영등포구($-0.013$), 관악구($-0.013$), 서초구($-0.009$)가 이에 해당한다. 관악구는 대학가 중심의 인구 구성과 디지털 역량 실측값이 맞물려 도움 요청 성격의 게시글이 드물게 나타난 것으로 볼 수 있으며, 서초구는 경제적 접근성과 기기 차원에서 특히 큰 음의 전이를 보였다. 사전 강도 $\kappa$를 10, 20, 30으로 달리한 민감도 분석에서도 사후 평균의 차이는 $0.00001\sim0.00010$ 수준으로 실질적 변화가 없었다. 이는 생활인구 보정 후 유효 관측 수가 충분해 사전의 가중치 변화가 결과에 큰 영향을 미치지 않음을 시사한다.

해석에서 주의할 점은 양의 전이가 곧바로 실측 지표가 과소평가되었음을 의미하지 않는다는 것이다. 양의 이동은 실측 지표가 포착하지 못한 생활상의 불편을 반영할 수 있지만, 특정 자치구에서 특정 문제에 대한 논의가 활발한 문화적\textperiodcentered{}인구학적 요인을 동시에 포함할 수 있다~\textbf{[인용 필요]}. 또한 플랫폼 이용자의 표집 편향과 게시글의 시공간적 불균형이 결과 해석에 영향을 미칠 수 있다. 따라서 사후 분포의 방향과 크기를 정책적 함의로 연결할 때는 정량\textperiodcentered{}정성적 지표를 종합적으로 고려해야 한다. 본 절의 결과에서 유의미한 것은 전이의 부호와 크기가 차원과 자치구에 따라 체계적으로 달라진다는 것이다. 단순히 당근마켓 게시글의 빈도만을 집계할 경우 플랫폼 활성도 자체가 UMC 수준에 대한 추정과 혼동될 우려가 있다. 사전 분포로 실측값을 투입하고 텍스트 분류 결과를 우도로 결합하는 절차는, 관측된 텍스트 빈도에서 플랫폼 활성도에 기인하는 부분과 UMC 수준에 기인하는 부분을 부분적으로 분리하는 효과를 갖는다. 이 점에서 사전-사후 분포의 전이는 텍스트 증거가 실측을 얼마나 보정하는지를 보여주는 지표로 기능한다.

다만 shift의 방향을 해석할 때 두 가지 제약을 함께 고려해야 한다. 첫째, 양의 shift가 관찰된다고 해서 해당 자치구의 UMC 수준이 실측보다 낮다고 곧바로 결론지을 수 없다. 양의 shift는 실측 지표가 포착하지 못하는 생활 수준의 체감이 존재할 가능성을 시사하지만, 관련 주제가 당근마켓이라는 특정 플랫폼에서 표출되는 방식 자체가 자치구별 이용자 구성에 의존하기 때문이다. 둘째, 음의 shift 또한 해당 자치구의 UMC 수준이 실측보다 높다는 의미로 해석하기보다, 관련 이슈가 플랫폼 상에서 표출될 유인이 상대적으로 적은 상태로 읽는 것이

\begin{table*}[htbp]
  \centering\small
  \caption{자치구별 게시글 판정 결과}
  \label{tab:7-district-cls}
  \begin{tabularx}{\textwidth}{l r r r r r}
    \toprule
    \textbf{자치구} & \textbf{총 게시글(n)} & \textbf{Y 판정} & \textbf{? 판정} & \textbf{N 판정} & \textbf{Y 비율(\%)} \\
    \midrule
    \textbf{중랑구}   & 5,476  & 794 & 28  & 4,654  & 14.5 \\
    \textbf{송파구}   & 9,367  & 860 & 99  & 8,408  & 9.18 \\
    \textbf{서초구}   & 6,170  & 434 & 216 & 5,520  & 7.03 \\
    \textbf{은평구}   & 6,140  & 429 & 50  & 5,661  & 6.99 \\
    \textbf{도봉구}   & 3,491  & 233 & 5   & 3,253  & 6.67 \\
    \textbf{서대문구} & 4,500  & 289 & 5   & 4,206  & 6.42 \\
    \textbf{강북구}   & 3,885  & 247 & 31  & 3,607  & 6.36 \\
    \textbf{용산구}   & 3,012  & 181 & 12  & 2,819  & 6.01 \\
    \textbf{양천구}   & 4,975  & 296 & 7   & 4,672  & 5.95 \\
    \textbf{강동구}   & 6,254  & 342 & 231 & 5,681  & 5.47 \\
    \textbf{영등포구} & 5,964  & 306 & 5   & 5,653  & 5.13 \\
    \textbf{동대문구} & 4,754  & 238 & 29  & 4,487  & 5.01 \\
    \textbf{성북구}   & 5,437  & 266 & 141 & 5,030  & 4.89 \\
    \textbf{중구}     & 1,351  & 63  & 16  & 1,272  & 4.66 \\
    \textbf{강서구}   & 8,174  & 371 & 236 & 7,567  & 4.54 \\
    \textbf{금천구}   & 2,298  & 101 & 7   & 2,190  & 4.40 \\
    \textbf{종로구}   & 2,267  & 97  & 27  & 2,143  & 4.28 \\
    \textbf{동작구}   & 3,675  & 150 & 22  & 3,503  & 4.08 \\
    \textbf{강남구}   & 10,047 & 404 & 459 & 9,184  & 4.02 \\
    \textbf{성동구}   & 404    & 16  & 9   & 379    & 3.96 \\
    \textbf{관악구}   & 11,665 & 407 & 75  & 11,183 & 3.49 \\
    \textbf{마포구}   & 7,974  & 242 & 254 & 7,478  & 3.03 \\
    \textbf{구로구}   & 4,082  & 121 & 96  & 3,865  & 2.96 \\
    \textbf{광진구}   & 5,440  & 147 & 42  & 5,251  & 2.70 \\
    \textbf{노원구}   & 5,040  & 103 & 46  & 4,891  & 2.04 \\
    \bottomrule
  \end{tabularx}
\end{table*}

\figfull{image20.png}{지역별 사전-사후 분포 차이}{fig:3-prior-posterior}

안전하다. 중랑구와 마포구의 대조에서 드러나듯, shift의 해석은 자치구의 인구 구성과 플랫폼 이용 특성을 비롯한 다양한 맥락 정보와 결합될 때 의미를 띈다. 둘째, 음의 전이 또한 해당 자치구의 UMC 수준이 실측보다 높다는 의미로 해석하기보다, 관련 이슈가 플랫폼 상에서 표출될 유인이 상대적으로 적은 상태로 읽는 것이 타당할 수 있다. 이런 측면에서 지역 자원이 실제로 결여되어 있음에도 플랫폼 상에서 논의되지 않는 이유를 고려하는 것 또한 중요하다.

\begin{table*}[htbp]
  \centering\small
  \caption{결여 유형별 가설 생성 결과 분포}
  \label{tab:8-deficit-types}
  \begin{tabularx}{\textwidth}{l Y r r}
    \toprule
    \textbf{결여 유형} & \textbf{관련 UMC 차원} & \textbf{관련 게시글 수 (n=100)} & \textbf{비율 (\%)} \\
    \midrule
    정보부재         & Digital Skills      & 63 & 63.0 \\
    제도적 경로 결여 & Available for Use   & 61 & 61.0 \\
    기기 접근 결여   & Devices             & 50 & 50.0 \\
    디지털 역량 부재 & Digital Skills      & 45 & 45.0 \\
    연결 인프라 결여 & Connectivity        & 40 & 40.0 \\
    공공 접근점 결여 & Available for Use   & 32 & 32.0 \\
    경제적 접근 결여 & Affordability       & 27 & 27.0 \\
    안전\textperiodcentered{}보안 결여 & Safety & 16 & 16.0 \\
    복합 결여 (2개 이상) & ---             & 88 & 88.0 \\
    \bottomrule
    \multicolumn{4}{p{0.97\textwidth}}{\footnotesize\textit{Note.} 하나의 게시글에서 복수의 결여 유형이 추론될 수 있으므로 비율의 합은 100\%를 초과한다. Phase 1은 가설 생성 단계이며, Phase 2(판정)에서 삼각검증을 거쳐 최종 분류가 확정된다.}
  \end{tabularx}
\end{table*}

\subsection{분석 결과 종합}
\label{subsec:3-4}

본 절에서는 \ref{subsec:3-1}절(UMC 지수 구축), \ref{subsec:3-2}절(HLM 다수준 분석), \ref{subsec:3-3}절(LLM 기반 텍스트 분석)의 결과를 종합하여, 서울의 디지털 격차가 어떤 구조로 작동하는지를 통합적으로 해석한다. 세 분석은 각각 '어디에 격차가 있는가'(측정), '왜 격차가 발생하는가'(인과 규명), '격차가 어떻게 경험되는가'(질적 맥락)에 답하며, 이들의 교차 지점에서 정책적으로 유의미한 함의가 도출된다.

세 분석이 공통으로 지지하는 첫 번째 발견이자 본 프로젝트의 핵심 실증 결과는, 서울의 디지털 격차가 자치구 간 차이가 아닌 개인 간 차이에 의해 구조화되어 있다는 것이다. \ref{subsec:3-2}절의 HLM 분석에서 ICC는 0.49\%에 불과하여, 디지털 활용 점수 분산의 99.5\%가 동일 자치구 내 개인 간 차이에 기인하였다. 이 결과는 \ref{subsec:3-1}절에서 구축한 UMC 종합지수의 자치구 간 편차가 실재함에도 불구하고, 해당 편차가 개인의 디지털 활용 수준에 미치는 직접적 영향은 제한적임을 의미한다. 그러나 이러한 낮은 ICC는 자치구 수준 변수의 무의미함을 뜻하는 것이 아니라, 그 효과가 집단별로 차별적으로 발현됨을 시사한다. Model 3의 교차수준 상호작용 결과가 이를 뒷받침한다.

\ref{subsec:3-3}절의 텍스트 분석은 이 양적 발견이 일상에서 어떻게 발현되는지를 구체화한다. '정보부재'가 가장 빈번한 결여 유형(63\%)으로 나타난 것은, 교육수준이 낮을수록 디지털 서비스의 존재 자체를 인지하지 못하는 정보 격차가 더 심화될 수 있음을 시사한다. 프린터 이용을 위해 편의점 무인 인쇄를 모르는 사례, ISP 고객센터 번호를 모르는 사례, 정부24의 공문서 출력 기능을 모르는 사례 등은 모두 '서비스가 존재하지만 이용자에게 보이지 않는' 구조적 문제를 가리킨다. 이는 HLM에서 교육 기울기가 가파른 이유를 경험적 수준에서 설명하는 보완적 증거로 해석할 수 있다.

두 번째 발견은, 자치구 인프라의 효과가 주효과 수준에서는 유의하지 않으나 특정 취약집단에 대해서는 조절 효과로서 유의하게 작용한다는 것이다. Model 2에서 UMC 차원 중 Connectivity와 관련된 인프라 변수의 주효과는 통계적 유의성에 도달하지 못하였는데($p = 0.237$), 이는 ICC가 0.49\%인 상황에서 자치구 수준 변수의 주효과가 제한적인 것과 일관된 결과이다. 그러나 Model 3에서 교차수준 상호작용이 유의하게 나타남으로써, 인프라의 효과가 평균적 효과가 아닌 집단별 조절 효과로서 작동함이 확인되었다.

\ref{subsec:3-3}절의 텍스트 분석은 이 차별적 효과의 질적 맥락을 제공한다. 공공 접근점(프린트 카페, 도서관 복합기, 무인 인쇄)에 대한 정보부재가 32\%의 게시글에서 추론된 것은, 인프라가 물리적으로 존재하더라도 이용자의 인지와 접근 역량이 수반되지 않으면 실질적 혜택으로 전환되지 않음을 보여준다. 은평구 신사동의 사례('정녕 없는 건지요')는 온라인 검색을 이미 시도했으나 결과가 불충분했을 가능성을 시사하며, 이는 인프라의 존재와 인프라의 '발견가능성(discoverability)' 사이의 간극을 조명한다. HLM에서 인프라 주효과가 유의하지 않았던 결과와 결합하면, 인프라 투자의 효과는 취약계층의 정보 접근과 역량 지원이 병행될 때 비로소 발현된다는 해석이 가능하다.

세 번째 발견은, 디지털 격차가 단일 축이 아닌 복합 취약성의 교차로 심화된다는 것이다. HLM에서 독거 고령자는 비독거 고령자 대비 추가적으로 4.0점이 낮았고, 연령 증가에 따른 저학력의 부정적 효과가 가속화되는 패턴이 확인되었다. \ref{subsec:3-3}절에서도 결여의 복합성이 일관되게 나타났다. 게시글당 평균 2.5개 이상의 결여 유형이 동시에 식별되었으며, 정보부재\textperiodcentered{}역량부재\textperiodcentered{}제도적 경로 차단이 중첩될 때 이용자가 공식 경로에서 이탈하여 커뮤니티에 의존하는 구조가 확인되었다. 디지털 지원망의 부재('자녀분들이 없으시거나 떨어져 살아서 묻기 힘들거나')는 가구 구성이 디지털 역량 획득의 비공식적 경로로 기능함을 보여주며, 독거 고령자의 추가적 격차를 설명하는 기제로 해석된다.

네 번째 발견은, 디지털 사막의 공간적 군집과 차원별 이질성이다. \ref{subsec:3-1}절에서 하위 5개 지역 중 특히 서울 북부(도봉구, 강북구, 노원구, 중랑구)가 UMC 종합지수의 Low-Low 클러스터를 형성하였으며, 이들은 Connectivity 차원에서는 도시 평균 이상을 기록하면서도 Devices와 Safety 차원에서 극단적 결핍을 보였다. 이러한 차원별 프로파일의 이질성은 일률적 인프라 투자보다 차원별 맞춤 개입의 필요성을 시사한다.

방법론적 관점에서, 세 분석의 결합은 각각의 한계를 상호 보완한다. \ref{subsec:3-1}절의 종합지수는 자치구 간 상대적 위치를 파악하되 개인 내 변이를 포착하지 못하며, \ref{subsec:3-2}절의 HLM은 개인과 지역의 기여를 분리하되 디지털 격차의 질적 발현 경로를 설명하지 못한다. \ref{subsec:3-3}절의 텍스트 분석은 결여의 경험적 맥락을 가시화하되 대표성에 한계가 있다. 세 분석을 병렬이 아닌 순차적 삼각검증(측정 $\rightarrow$ 인과 규명 $\rightarrow$ 질적 맥락)으로 결합함으로써, '어디에'(자치구), '왜'(개인 특성과 인프라의 상호작용), '어떻게'(결여의 발현 경로)라는 세 질문에 대한 통합적 답을 구성할 수 있었다.

이상의 종합 결과는 4장의 정책 제언으로 직결된다. 개인 특성이 지배적인 설명변수라는 발견은 people-based 정책의 우선성을, 인프라의 차별적 효과는 취약계층 타겟형 place-based 정책의 필요성을, 복합 취약성의 확인은 단일 프로그램이 아닌 통합 지원 체계의 필요성을, 공간적 군집의 확인은 권역 단위 거버넌스의 근거를 각각 제공한다.
)
\section{분석}
\label{sec:analysis}

\noindent{\sffamily\small\color{muted}\textbf{CHAPTER 3}}

\begin{quoteblock}
측정 $\rightarrow$ HLM $\rightarrow$ 베이지안 통합의 3단계 분석으로 서울 디지털 격차의 공간 분포, 개인-지역 기여 분리, 그리고 텍스트 기반 잠재 격차 신호까지 통합 추정한다.
\end{quoteblock}

\figfull{image6.png}{핵심 발견 요약}{fig:3-keyfindings}
\figfull{image7.png}{디지털 사막 --- 종합지수 하위 4개 자치구}{fig:3-desert}

본 장의 세 분석은 순차적으로 연결된다. 먼저 서울 자치구별 UMC 지수를 구축하여 지역 격차를 측정하고(\ref{subsec:3-1}), 이 결과를 지역 수준 변수로 활용하여 개인 및 지역 특성이 개인의 디지털 연결성에 미치는 영향을 분석한다(\ref{subsec:3-2}). 마지막으로 디지털 플랫폼 텍스트를 통해 양적 분석이 포착하지 못하는 구체적 사례를 도출한다. 이 흐름은 측정 $\rightarrow$ 인과 규명 $\rightarrow$ 현상 보완으로 구성되며, \ref{subsec:3-4}절에서 세 분석의 결과를 종합한다.

\subsection{서울의 UMC 지역 격차}
\label{subsec:3-1}

\subsubsection{측정 지표 및 데이터}
\label{subsubsec:3-1-1}

본 프로젝트에서 구성한 측정 지표는 UMC의 6개 차원에 대한 개념적 정의를 바탕으로 한다. 각 차원은 Connectivity, Available for Use, Affordability, Devices, Skills, Safety로 구분되며, 개별 측정 지표는 의미성과 보편성을 동시에 고려한 두 가지 핵심 구성 요소에 따라 설정하였다~\citep{ITU2023}. 본 분석에 활용된 서울의 UMC를 구성하는 6개 차원별 측정 지표는 표~\ref{tab:1-umc-dims}과 같다.

측정 지표 선정 과정에서는 개념적 정의와의 부합성, 서울의 디지털 기술 환경, 데이터 가용성을 종합적으로 검토하였다. 분석에는 데이터를 확보할 수 있는 가장 최근 시점인 2023년과 2024년, 총 2개 시점의 자료를 활용하였다. 2개 시점의 국가 행정 자료 등을 활용하였다. 단, Devices, Skills, Security 차원은 데이터 제공 시점의 한계로 인해 2023년 조사 자료를 기준으로 두 시점 모두 동일한 값을 적용하였다.

Connectivity의 경우, 한국전파누리에서 제공하는 이동통신 무선국 데이터와 서울시 행정데이터인 생활인구 데이터를 활용하여 측정 지표를 산출하였다. \citet{ITU2023}에 따르면, Connectivity는 네트워크 커버리지 내에 거주하는 인구를 기준으로 측정될 수 있다. 그러나 한국은 직주 분리가 뚜렷한 경우가 많아, 거주지역과 실질적으로 디지털 환경에 접속하는 지역이 상이할 수 있다. 이에 본 연구에서는 단순 거주인구 대신 주간 유동인구 데이터를 활용하여 4G 및 5G 네트워크 커버리지를 산출함으로써, 한국의 맥락에 보다 적합한 지표를 고안하였다. 구체적으로, 4G 기지국 수와 5G 기지국 수를 각각 주간 생활인구 천 명당 기지국 수로 환산하여 집계구별로 산출한 뒤, 이를 자치구별로 평균한 값을 최종 지표로 사용하였다.

Available for Use는 디지털 서비스의 실질적 이용 가능성을 측정하는 차원이다. SKT 통신정보에서 제공하는 월별 지역 내 가구당 모바일 데이터 사용량과 금융, 쇼핑, 동영상, 배달 4개 범주의 온라인 서비스 이용일수를 인구가중 평균으로 자치구별 월별로 집계하고, 12개월 평균값을 사용하였다. 여기에 자치구 지방정부의 공공WiFi 등록 현황 데이터를 활용하여 전체 인구 1,000명당 공공WiFi AP 수를 산출하여 포함하였다. 공공WiFi 밀도는 가정 내 인터넷 접속 환경이 제한적인 고령층에게 특히 중요한 디지털 접근 경로로 기능한다~\citep{vanDijk2020}.

Affordability는 ITU ICT Development Index 2023의 두 가지 축 구조를 따라 구성하였다~\citep{ITU2024}. 구매력(capacity) 측면에서는 서울시 상권분석서비스에서 제공하는 자치구별 월평균 가구소득을 활용하고, 가격부담(burden) 측면에서는 SKT 통신정보의 최근 3개월 이내 요금 연체 비율을 활용하였다. 연체율은 행정동, 성별, 연령대별 인구가중 평균으로 산출한 뒤 자치구 단위로 집계하였으며, 값이 높을수록 경제적 접근성이 낮음을 의미하는 역방향 지표이다.

Devices는 디지털 기기 보유 및 이용의 다양성을 측정한다. 서울시민 디지털역량실태조사의 가구 내 핵심 디지털 기기 보유 유형 수와 개인이 실제로 사용하는 기기 유형 수를 자치구별 가중평균으로 산출하였다. \citet{ITU2023}의 보고서에 따르면, 휴대폰과 컴퓨터 등 기기의 특성에 따라서 기기의 보유와 이용은 서로 다르게 해석된다. 이를 한국 국내 맥락에선, 컴퓨터의 보유 여부도 디지털 연결성을 드러내는 지표가 될 수 있다. 스마트폰 보급률이 사실상 포화 상태인 서울에서는 데스크톱, 노트북, 태블릿 등 보조 기기의 보유 여부가 디지털 활용 역량의 차별화 요인으로 작용한다~\citep{vanDeursenHelsper2015}.

Digital Skills는 기초, 응용, 고급, 사이버활동의 4개 하위 영역으로 구성된다. 기초 역량은 스마트 기기 기본 조작 능력을, 응용 역량은 디지털 서비스 활용 능력을, 고급 역량은 스마트 기기 활동 숙련도를, 사이버활동 역량은 사이버 공간에서의 활동 능력을 각각 측정한다. 각 하위 영역의 원점수를 만점으로 나누어 0\,$\sim$\,1 비율로 변환한 뒤 자치구별 가중평균을 산출하였다.

Safety는 디지털 보안 행동과 보안 인식의 2개 하위 지표로 구성된다. 보안 행동은 비밀번호 관리, 소프트웨어 업데이트, 2단계 인증 사용 등 실제 보안 실천 행위를 측정하고, 보안 인식은 개인정보 보호의 중요성에 대한 인지 수준을 측정한다. 각 하위 지표를 만점 기준 비율로 변환하여 자치구별 가중평균을 산출하였다.

\begin{table*}[htbp]
  \centering\small
  \caption{UMC 차원별 측정 지표}
  \label{tab:1-umc-dims}
  \begin{tabularx}{\textwidth}{l Y Y Y}
    \toprule
    \textbf{Dimension} & \textbf{Indicators} & \textbf{Data Source} & \textbf{Note} \\
    \midrule
    Connectivity & station\_density\_4g, station\_density\_5g, download\_speed & Spectrum Resource Mgmt System; NIA Communication Quality & Per 1K daytime living pop; CQ uses 2022--2024 avg \\
    \addlinespace[2pt]
    Available for Use & mobile\_data\_usage, online\_service\_days, wifi\_density & SKT Telecom Data; Local Admin WiFi License Data & Pop-weighted district avg; WiFi per 1K total pop \\
    \addlinespace[2pt]
    Affordability & avg\_income, delinq\_rate (reverse) & Seoul Commercial Analysis Service; SKT Telecom Data & ITU IDI capacity + burden; delinq = reverse-coded \\
    \addlinespace[2pt]
    Devices & device\_home\_core, device\_use\_core & Seoul Digital Competency Survey 2023 & Home: excl. smartphone; Use: incl. smartphone \\
    \addlinespace[2pt]
    Digital Skills & Basic / Applied / Advanced / Cyber activity & Seoul Digital Competency Survey 2023 & --- \\
    \addlinespace[2pt]
    Safety & safety\_behavior, safety\_awareness & Seoul Digital Competency Survey 2023 & Security practice (7 items) + Awareness (3 items) \\
    \bottomrule
  \end{tabularx}
\end{table*}

\subsubsection{측정 지표 지수화}
\label{subsubsec:3-1-2}

모든 원시 지표는 25개 자치구 내에서 Min-Max 스케일링을 통해 $[0,1]$ 구간으로 정규화한다. 정규화된 값 $x^{*}$는 식~\eqref{eq:minmax}과 같이 산출한다.

\begin{equation}\label{eq:minmax}
  x^{*} = \frac{x_{i} - x_{\min}}{x_{\max} - x_{\min}}
\end{equation}

이때, 값이 클수록 나쁜 결과를 의미하는 역방향 지표는 변환하여 모든 지표의 방향성을 통일하였다. 그 다음, 각 차원 내에서 정규화된 지표들을 동일 가중치 산술 평균으로 결합하여 차원별 점수를 산출하였다. 동일 가중치로 분석을 진행한 이유는, 6개 차원 간 상대적 중요도에 대한 합의된 기준이 부재한 상태에서 연구자의 자의적 판단을 최소화하기 위함이다~\citep{ITU2024}. 산출된 종합지수는 0(최저)에서 1(최고) 범위를 가지며, 서울 25개 자치구 전체에 대해 2023년과 2024년 각각 독립적으로 산출한다. 해석에 있어 중요한 점은 이 종합 지수가 지역의 절대적 UMC가 아닌, 서울 내 다른 자치구 간의 상대적 위치를 보여준다는 것이다.

\subsubsection{지역별 UMC 지수 특성}
\label{subsubsec:3-1-3}

25개 자치구의 2024년 UMC 종합지수는 0.279(중랑구)에서 0.695(서초구)까지 분포하였으며, 최상위와 최하위 간 약 2.5배의 격차가 존재하였다. 서울이 높은 인터넷 보급률을 가진 도시라는 일반적 인식에도 불구하고, 하위 행정단위에서는 상당한 상대적 수준의 디지털 연결성 격차가 관찰되었다. 상위 5개 자치구는 서초구(0.695), 영등포구(0.686), 마포구(0.665), 서대문구(0.615), 용산구(0.603)이다. 하위 5개 자치구는 중랑구(0.279), 도봉구(0.303), 강북구(0.380), 구로구(0.389), 노원구(0.389)이다. 상위 자치구는 대체로 도심부 또는 강남권에 위치하여 밀집된 상업 활동, 높은 소득 수준, 풍부한 민간 디지털 인프라의 혜택을 받는 반면, 하위 자치구는 북부 외곽 및 남서부에 집중되어 있었다.

종합지수의 공간적 분포를 살펴보면(그림~\ref{fig:3-umc-comp}), UMC 지수가 높게 나온 자치구는 두 구역에 집중되어 있었다. 이는 각각 서초-영등포-마포로 이어지는 한강과 도심권 지역과, 용산-종로-중구 중심의 전통적 도심부이다. 반면, 북부 외곽(강북구, 도봉구, 노원구)과 남서부(구로구, 금천구)는 도시 평균 이하의 점수를 보였다. 이러한 공간 분포는 서울의 사회경제적인 자원 분포와 연관되어 있었다. 즉, 디지털 격차가 도시의 구조적 공간 불평등과 밀접히 연관되어 있을 가능성을 시사하였다.

차원별 분석에서 한편으로, 어떤 자치구도 6개 차원 모두에서 균일하게 우수하지는 않다는 점이 발견되었다(그림~\ref{fig:3-umc-heatmap}). 예컨대 서울 북부에 위치한 은평구의 경우, Connectivity 차원에서 0.781로 상위권에 위치하지만, Affordability 차원에서는 0.328로 하위권에 그쳤다. 유사하게, 관악구는 Skills 차원에서 0.795로 1위를 기록하였지만 Affordability에서는 0.152로 최하위권이다. 이러한 다차원적 지역간 불균등은, 종합 순위만으로는 정책적으로 의미 있는 차원별 취약성을 포착할 수 없음을 의미하였다.

\begin{keymsg}[title=핵심 발견 \textperiodcentered{} 디지털 사막 4개 자치구]
종합지수 하위 자치구(중랑\textperiodcentered{}도봉\textperiodcentered{}강북\textperiodcentered{}노원)는 Devices\textperiodcentered{}Safety 차원에서 동반 결핍을 보이며, 공간적으로도 서울 북부에 집중되어 있다. 이는 권역 단위 협력 거버넌스 구축의 정당화 근거를 제공한다.
\end{keymsg}

결론적으로, 그림~\ref{fig:3-umc-dim}에 제시된 바와 같이 공간적으로 밀집한 종합지수 하위 4개 자치구(중랑구, 도봉구, 강북구, 노원구)를 본 프로젝트에서는 잠정적으로 디지털 사막(Digital Desert)으로 분류한다. 이들 자치구는 공통적으로 Devices와 Safety 차원에서 뚜렷한 결핍을 보이며, Affordability 역시 서울시 평균에 미치지 못했다. 한편, Connectivity 차원에서는 일부 자치구(중랑구 0.630, 강북구 0.562)가 서울시 평균을 상회하여, 물리적 네트워크 인프라 자체의 부재보다는 이를 활용하기 위한 인적\textperiodcentered{}경제적 역량의 부족이 핵심 격차 요인일 가능성을 보여준다. 이러한 발견은 \ref{subsec:3-2}절의 다층 분석에서 개인 수준 변수와 자치구 수준 변수의 상대적 기여를 분리 및 검증함으로써 보다 체계적으로 검토된다.

또한, UMC의 일부 차원에서 지리적으로 인접한 지역이 서로 유사한 패턴을 보였다. 그림~\ref{fig:3-umc-dim}은 차원별 지역의 공간적 자기상관성을 나타낸 것이다. 공간적으로 군집해 있는 자치구가 서로 유사한 패턴을 보이는 점은 디지털 사막 해소를 위한 광역적 거버넌스 체계 구축의 가능성을 시사한다.

결론적으로, 그림~\ref{fig:3-umc-dim}에 제시된 바와 같이 종합지수 하위 5개 자치구(중랑구, 도봉구, 강북구, 구로구, 노원구)를 본 프로젝트에서는 잠정적으로 디지털 사막(Digital Desert)으로 분류한다. 이들 자치구는 공통적으로 Devices와 Safety 차원에서 뚜렷한 결핍을 보이며, Affordability 역시 서울시 평균에 미치지 못했다. 한편, Connectivity 차원에서는 일부 자치구(중랑구 0.630, 강북구 0.562)가 서울시 평균을 상회하여, 물리적 네트워크 인프라 자체의 부재보다는 이를 활용하기 위한 인적\textperiodcentered{}경제적 역량의 부족이 핵심 격차 요인일 가능성을 보여준다. 이러한 발견은 \ref{subsec:3-2}절의 다층 분석에서 개인 수준 변수와 자치구 수준 변수의 상대적 기여를 분리 및 검증함으로써 보다 체계적으로 검토된다. 또한, UMC의 일부 차원에서 지리적으로 인접한 지역이 서로 유사한 패턴을 보였다. 그림~\ref{fig:3-lisa-grid}은 차원별 지역의 공간적 자기상관성을 나타낸 것이다. 공간적으로 군집해 있는 자치구가 서로 유사한 패턴을 보이는 점은 디지털 사막 해소를 위한 광역적 거버넌스 체계 구축의 가능성을 시사한다.

\begin{figure}[htbp]
  \centering
  \begin{subfigure}[t]{0.55\linewidth}
    \centering
    \includegraphics[width=\linewidth]{image9.png}
    \caption{자치구별 UMC 종합점수}
    \label{fig:3-umc-comp}
  \end{subfigure}\hfill
  \begin{subfigure}[t]{0.43\linewidth}
    \centering
    \includegraphics[width=\linewidth]{image10.png}
    \caption{UMC 차원별 점수 히트맵}
    \label{fig:3-umc-heatmap}
  \end{subfigure}
  \caption{자치구별 UMC 종합점수 및 차원별 히트맵}
  \label{fig:3-umc-two}
\end{figure}

\figfull{image11.png}{UMC 차원별 점수}{fig:3-umc-dim}

\begin{figure*}[htbp]
  \centering
  \begin{subfigure}[t]{0.32\textwidth}
    \centering
    \includegraphics[width=\linewidth]{image12.jpeg}
    \caption{Connectivity}
  \end{subfigure}\hfill
  \begin{subfigure}[t]{0.32\textwidth}
    \centering
    \includegraphics[width=\linewidth]{image13.jpeg}
    \caption{Available for Use}
  \end{subfigure}\hfill
  \begin{subfigure}[t]{0.32\textwidth}
    \centering
    \includegraphics[width=\linewidth]{image14.jpeg}
    \caption{Affordability}
  \end{subfigure}

  \vspace{4pt}
  \begin{subfigure}[t]{0.32\textwidth}
    \centering
    \includegraphics[width=\linewidth]{image15.jpeg}
    \caption{Devices}
  \end{subfigure}\hfill
  \begin{subfigure}[t]{0.32\textwidth}
    \centering
    \includegraphics[width=\linewidth]{image16.jpeg}
    \caption{Skills}
  \end{subfigure}\hfill
  \begin{subfigure}[t]{0.32\textwidth}
    \centering
    \includegraphics[width=\linewidth]{image17.jpeg}
    \caption{Safety}
  \end{subfigure}
  \caption{UMC 수준의 공간적 자기상관 (서울 25개 자치구)}
  \label{fig:3-lisa-grid}
\end{figure*}

\subsection{개인 및 지역 특성과 디지털 연결성}
\label{subsec:3-2}

본 연구의 데이터는 개인(Level 1)이 자치구(Level 2) 내에 군집되어 있는 위계적 구조를 갖는다. 여기서 Level 2는 개인의 거주지에 해당한다. 이러한 구조에서 통상적인 OLS 회귀분석은 관측치의 비독립성 가정을 위반하여, 자치구 수준 예측변수의 표준오차를 과소추정하고 1종 오류를 팽창시킬 수 있다~\citep{RaudenbushBryk2002}. 따라서 분산을 25개 자치구를 기준으로 구 내와 구 간 성분으로 분해하고, 개인 수준과 자치구 수준 변수를 동시에 투입할 수 있는 2수준 위계적 선형 모형(HLM)을 채택하였다.

두 번째 분석에서 활용한 자료인 서울서베이는 2003년부터 서울시에서 2만 개 가구를 대상으로 년마다 실시해온 행정 자료이다. 자료의 수집은 매년 8\,$\sim$\,9월에 비대면과 방문면접을 병행하여 이뤄진다.

\subsubsection{변수 구성}
\label{subsubsec:3-2-1}

먼저 분석의 종속변수인 디지털 활용 점수는 서울서베이의 응답항목 중 디지털 서비스 이용빈도 항목의 8개 항목을 합산한 점수이다. 8개 항목은 생활정보 서비스, 전자상거래 서비스, 금융거래 서비스, 공공 서비스, 커뮤니케이션 서비스, 여가 서비스, 교육학습 서비스, 모바일 결제로 구성된다. 이들은 전화와 문자와 같은 기초적 연결 행위가 아니라, 디지털 기기를 매개로 한 실질적인 추가 활용에 해당하여, 디지털 환경이 발달한 서울에 적합한 지표로 판단하였다. 각 항목은 원래 4점 서열 척도를 0--10 연속 척도로 변환하여 총 0--80점 범위의 연속변수를 구성한다. 이 점수가 높을수록 개인의 디지털 서비스 활용이 다양하고 원하는 서비스를 자유롭게 이용할 수 있음을 의미한다.

본 분석의 독립변수는 개인 특성의 변수와 지역 특성의 변수로 구분된다(표~\ref{tab:2-hlm-vars}). 개인 특성의 변수는 개인의 디지털 활용에 영향을 미칠 수 있는 변수인 성별, 연령, 직업, 학력, 소득수준, 장애여부로 구성된다. 교육수준은 중졸 이하, 고졸, 대졸(기준범주), 대학원 이상의 4범주로, 직업은 관리\textperiodcentered{}전문직, 사무직(기준범주), 생산\textperiodcentered{}기능직, 기타(학생\textperiodcentered{}주부\textperiodcentered{}무직)의 4범주로 구분하였다. 주거유형은 아파트(기준범주), 단독주택, 다세대 및 다가구, 기타의 4범주로 분류하였다. 연령과 소득은 총평균 중심화를 적용하여 절편의 해석 가능성을 높였다. 추가로, 1인가구 더미, 장애 여부 더미, 연도 더미(기준: 2023년)를 포함하였다.

지역 수준(Level 2) 변수는 \ref{subsec:3-1}절에서 구축한 UMC 차원 점수 중 지역의 디지털 환경을 결정하는 외생적 요인의 세 차원 Connectivity, Available for Use, Affordability을 활용하였다. 이는, Devices, Digital Skills, Safety가 서울디지털역량실태조사의 개인 응답을 자치구 단위로 집계한 것을 고려한 것이다. 그로 인한 내생성 문제를 피하기 위해 주 분석에서 세 가지 차원을 제외하여 지역 차원의 보편적인 의미있는 연결성에 더 적합하다고 판단되는 차원을 선택하였다. 통제변수로 자치구 평균소득과 고령화율을 포함하였으며, 모든 Level-2 변수는 앞선 \ref{subsec:3-1}절에서와 같이 Min-Max 정규화를 적용하여 $[0,1]$ 범위로 표준화하였다.

분석은 4단계로 진행된다. 무조건 모형에서 급내상관계수를 산출하여 다층 분석의 정당성을 확인한 뒤, 개인 수준 변수를 투입한 M1, 자치구 수준 변수를 추가한 M2, 교차수준 상호작용을 포함한 M3를 순차적으로 추정한다. 교차수준 상호작용은 UMC 인프라 변수(Connectivity, Available for Use, Affordability)와 취약계층 지표의 곱으로 구성하여, 자치구의 디지털 인프라 수준이 취약계층의 디지털 활용에 차별적으로 작용하는지를 검증한다.

\subsubsection{분석 결과}
\label{subsubsec:3-2-2}

분석에 활용된 표본은 서울서베이 2023\,$\sim$\,2024년 응답자 총 10,000명이며, 25개 자치구에서 평균적으로 400명이 표집되었다. 주요 변수에 대한 기술통계량을 나타낸 것은 표~\ref{tab:3-descstat}과 같다. 개인의 디지털 연결성의 평균이 80점 만점에 60점으로 나타났는데, 서울시민의 높은 디지털 연결성을 잘 보여준다. 종속변수가 음의 왜도가 존재할 수 있으나 기관 표준 척도와 10000명에 근접한 표본을 고려해 척도 선형화한 원점수대로 분석을 진행하였다.

표~\ref{tab:4-hlm-results}은 순차적으로 다층 모형의 추정 결과를 나타낸 것이다. 무조건 모형의 급내상관계수는 0.49\%로 자치구 간 분산 비율 자체는 크지 않았으나, 자치구당 평균 표본 크기 400명을 감안한 설계효과가 2.94로 임계값 2.0을 초과하여 다수준 모형의 사용에 대한 타당성을 확인하였다~\citep{HoxEtAl2018}. 그러나, 25개의 자치구의 소규모 표본 / 자치구가 설명하는 분산이 매우 적은 점을 종합하면 개인의 디지털 연결성에 지역의 연결 관련 특성이 미치는 영향이 극히 적을 것으로 추정된다. 이러한 결과는 서울의 우수한 디지털 연결성과 높은 수준의 디지털 역량으로 지역 수준의 연결성 인프라가 영향을 미치지 못한 것일 수 있다. 또한, 서울서베이는 개인의 거주 지역을 조사하기 때문에, 실제로 연결성에 직접적으로 영향을 받는 지역과 일치하지 않을 수 있다.

개인 특성 중 교육 수준의 관계가 가장 두드러졌다. 중졸 이하 학력은 대졸 대비 디지털 활용 점수를 약 22.6점 낮게 나타났다. 이는 80점 만점 기준 28\%에 해당하는 상당한 격차다. 고졸 역시 대졸보다 7.3점 낮게 나타난 반면, 대학원 이상은 오히려 2.9점 높아 교육 수준과 디지털 활용 능력 간의 뚜렷한 정적 관계를 확인할 수 있다. 연령 또한 유의한 영향을 미쳤다. 평균 연령으로부터 1세 증가할수록 개인의 디지털 연결성에 관한 점수가 0.35점 감소하였으며, 65세 이상 고령층의 경우 연령 효과와는 별개로 추가적인 6.0점 하락이 확인되었다. 65세 이상의 연령층에서 디지털 연결성의 두드러진 저하가 나타났다. 직업 측면에서는 관리직\textperiodcentered{}전문직 종사자가 사무직 대비 2.3점 높게 나타나, 고소득 직종일수록 디지털 활용 수준도 높은 경향을 보였다. 주거환경의 경우 단독주택 거주자는 아파트 대비 2.5점, 다세대주택 거주자는 1.8점 낮았는데, 이는

\begin{table*}[htbp]
  \centering\small
  \caption{분석 변수 설명}
  \label{tab:2-hlm-vars}
  \begin{tabularx}{\textwidth}{l Y c l l}
    \toprule
    \textbf{Variable} & \textbf{Description} & \textbf{Level} & \textbf{Scale} & \textbf{Source} \\
    \midrule
    digital\_use\_score & Digital service utilization score (8 items summed) & L1 & 0--80 (continuous) & Seoul Survey 2023--24 \\
    \addlinespace[2pt]
    age\_c & Age (grand-mean centered, M = 49.3) & L1 & years & Seoul Survey \\
    \addlinespace[2pt]
    female & Female (ref: male) & L1 & 0/1 & Seoul Survey \\
    \addlinespace[2pt]
    Edu & Education: $\leq$middle / high / grad+ (ref: college) & L1 & 0/1 each & Seoul Survey \\
    \addlinespace[2pt]
    Income & Monthly income (grand-mean centered) & L1 & 10K KRW & Seoul Survey \\
    \addlinespace[2pt]
    job & Job: professional / blue-collar / other (ref: white-collar) & L1 & 0/1 each & Seoul Survey \\
    \addlinespace[2pt]
    disabled\_bin & Disability status & L1 & 0/1 & Seoul Survey \\
    \addlinespace[2pt]
    hh\_size\_1 & Single-person household & L1 & 0/1 & Seoul Survey \\
    \addlinespace[2pt]
    housing\_detached / \_multi / \_other & Housing: detached / multi-family / other (ref: apt) & L1 & 0/1 each & Seoul Survey \\
    \addlinespace[2pt]
    year & Year 2024 dummy (ref: 2023) & L1 & 0/1 & Seoul Survey \\
    \addlinespace[2pt]
    elderly & Elderly ($\geq$65) & L1 & 0/1 & Seoul Survey \\
    \addlinespace[2pt]
    low\_edu & Low education ($\leq$middle school) & L1 & 0/1 & Seoul Survey \\
    \addlinespace[2pt]
    connectivity & Connectivity dimension score (normalized) & L2 & 0--1 & UMC Part 1 \\
    \addlinespace[2pt]
    wifi\_total\_n & Public WiFi density per 1K pop (normalized) & L2 & 0--1 & Local Admin License \\
    \bottomrule
  \end{tabularx}
\end{table*}

주택 유형에 따른 소득 수준의 차이뿐 아니라 건물 유형별 인터넷 연결 여건의 차이도 반영된 결과로 해석된다. 1인 가구 역시 다인 가구 대비 1.8점 낮은 점수를 보였다. 그 외 개인 특성으로는, 여성이 남성 대비 2.6점 높게 나타났다. 여성이 남성에 비해 더 높은 디지털 연결성을 갖는 것은, 국내 다양한 디지털 관련 서비스에 여성이 더 민감하게 반응한 결과일 수 있다. \textbf{[인용 필요]} 이러한 결과는 국가에 따른 사회문화적인 맥락을 고려할 필요가 있다. 장애인의 경우 3.1점 높게 나타났는데, 표본의 특성에 기인한 것으로 보인다.

자치구 수준의 UMC 인프라 차원은 주효과만으로는 대체로 유의하지 않았다. Model 2에서 세 차원 모두 통계적 유의성에 도달하지 못하였고, 우도비 검정 결과도 Model 1 대비 유의한 개선이 없었다. 그러나 교차수준 상호작용을 투입한 Model 3에서는 Infrastructure 차원이 유의한 부정적 효과를 보였다. 이는 인프라의 평균적 효과가 아닌, 집단별 차별적 효과가 존재함을 의미한다.

본 분석의 핵심 발견은 교차수준 상호작용에서 나타났다. 공공 WiFi 밀도와 고령 더미의 상호작용은 3.03점의 유의한 정적 효과를 보였다. 이는 공공 WiFi 인프라가 비고령자보다 고령자에게 상대적으로 더 큰 혜택을 제공함을 뜻한다. 디지털 인프라 수준이 높은 자치구에 거주하는 저학력 주민은 인프라 수준이 낮은 곳의 저학력 주민 대비 현저히 높은 디지털 활용 점수를 기록하였다. 이 결과는 인프라 투자가 교육 취약계층의 디지털 격차를 완화하는 데 불균형적으로 큰 효과를 가질 수 있음을 시사한다.

개인 수준 상호작용에서는 이중 취약성이 확인되었다. 독거 고령자는 비독거 고령자 대비 추가적으로 4.0점이 낮았으며, 고령일수록 저학력의 부정적인 효과가 더 커지는 것으로 나타났다. 이러한 결과는 단일 취약 요인보다 복합 취약 요인이 디지털 격차를 심화시킴을 보여준다.

\begin{table*}[htbp]
  \centering\small
  \caption{주요 변수의 기술통계량}
  \label{tab:3-descstat}
  \begin{tabularx}{\textwidth}{Y r r r r r}
    \toprule
    \textbf{변수(Variable)} & \textbf{N} & \textbf{평균(Mean)} & \textbf{표준편차(SD)} & \textbf{최소(Min)} & \textbf{최대(Max)} \\
    \midrule
    디지털 활용 점수(Digital use score) & 10,000 & 60.88 & 16.84 & 0 & 80 \\
    연령(Age) & 10,000 & 49.32 & 16.23 & 15 & 99 \\
    여성 비율(Female, proportion) & 10,000 & 0.51 & --- & 0 & 1 \\
    학력: 중졸 이하(Education: $\leq$Middle school) & 10,000 & 0.07 & --- & 0 & 1 \\
    학력: 고졸(Education: High school) & 10,000 & 0.25 & --- & 0 & 1 \\
    학력: 대학원 이상(Education: Graduate+) & 10,000 & 0.13 & --- & 0 & 1 \\
    장애 여부(Disability) & 10,000 & 0.03 & --- & 0 & 1 \\
    월평균 소득(만원 단위, Monthly income [10K KRW]) & 10,000 & 237.55 & 209.73 & 0 & 1,050 \\
    1인 가구(Single-person household) & 10,000 & 0.17 & --- & 0 & 1 \\
    고령(65세 이상, Elderly 65+) & 10,000 & 0.22 & --- & 0 & 1 \\
    (군집) 연결성 지수(Connectivity, norm.) & 50 & 0.5 & 0.18 & 0.18 & 0.98 \\
    이용 가능성(Available for Use, norm.) & 50 & 0.5 & 0.19 & 0.12 & 0.96 \\
    비용 접근성(Affordability, norm.) & 50 & 0.5 & 0.22 & 0.07 & 1 \\
    \bottomrule
    \multicolumn{6}{p{0.97\textwidth}}{\footnotesize\textit{Note.} N = 10,000명(2023\,$\sim$\,2024년, 25개 자치구, 표본 합산). 소득은 만원 단위. 자치구(군집) 단위 변수는 0\,$\sim$\,1(min-max) 정규화 지수임. N = 10,000 individuals nested within $j$ = 25 districts (2023--2024 pooled). Income measured in 10,000 KRW units. District-level variables are min-max normalized (0--1 scale).}
  \end{tabularx}
\end{table*}

\begin{table*}[htbp]
  \centering\small
  \caption{HLM 추정 결과}
  \label{tab:4-hlm-results}
  \begin{tabularx}{\textwidth}{l Y Y Y}
    \toprule
    \textbf{변수} & \textbf{Model 1} & \textbf{Model 2} & \textbf{Model 3} \\
    \midrule
    \multicolumn{4}{l}{\textit{Level 1: 개인 특성}} \\
    \addlinespace[2pt]
    연령(중심화) & $-0.347^{***}$ (0.012) & $-0.347^{***}$ (0.012) & $-0.278^{***}$ (0.014) \\
    여성 & $2.636^{***}$ (0.305) & $2.634^{***}$ (0.305) & $2.601^{***}$ (0.305) \\
    학력: 중학교 이하 & $-22.618^{***}$ (0.770) & $-22.610^{***}$ (0.770) & $-19.019^{***}$ (2.379) \\
    학력: 고등학교 & $-7.289^{***}$ (0.380) & $-7.282^{***}$ (0.380) & $-7.258^{***}$ (0.380) \\
    학력: 대학원 이상 & $2.906^{***}$ (0.483) & $2.909^{***}$ (0.483) & $2.912^{***}$ (0.481) \\
    소득(중심화) & $0.003^{**}$ (0.001) & $0.003^{**}$ (0.001) & $0.002^{**}$ (0.001) \\
    직업: 전문직 & $2.329^{***}$ (0.451) & $2.326^{***}$ (0.451) & $2.28^{***}$ (0.449) \\
    직업: 생산직 & $-0.373$ (0.424) & $-0.379$ (0.424) & $-0.385$ (0.423) \\
    직업: 기타 & $-1.233^{**}$ (0.456) & $-1.226^{**}$ (0.456) & $-1.105^{*}$ (0.454) \\
    장애 & $3.084^{***}$ (0.888) & $3.081^{***}$ (0.888) & $3.006^{***}$ (0.884) \\
    1인 가구 & $-1.775^{***}$ (0.338) & $-1.776^{***}$ (0.338) & $-0.877^{*}$ (0.367) \\
    단독주택 & $-2.478^{***}$ (0.377) & $-2.480^{***}$ (0.377) & $-2.574^{***}$ (0.379) \\
    다세대주택 & $-1.798^{***}$ (0.408) & $-1.795^{***}$ (0.408) & $-1.933^{***}$ (0.408) \\
    연도(2024) & $0.294$ (0.283) & $0.292$ (0.283) & $0.292$ (0.283) \\
    고령자(65세 이상) & $-6.046^{***}$ (1.696) & & \\
    \addlinespace[4pt]
    \multicolumn{4}{l}{\textit{Level 2: 자치구 특성}} \\
    \addlinespace[2pt]
    연결성 & $-1.127$ (0.747) & $-1.688^{*}$ (0.804) & \\
    이용 가능성(WiFi) & $0.085$ (0.780) & $-0.200$ (0.829) & \\
    비용 접근성 & $-1.578$ (1.135) & & \\
    자치구 소득 & $-2.035^{\dagger}$ (1.181) & & \\
    고령화율 & $-0.298$ (1.095) & $-0.293$ (1.111) & \\
    \addlinespace[4pt]
    \multicolumn{4}{l}{\textit{상호작용}} \\
    \addlinespace[2pt]
    고령자 $\times$ 1인 가구 & $-3.958^{***}$ (0.936) & & \\
    연령 $\times$ 저학력 & $-0.277^{***}$ (0.080) & & \\
    \addlinespace[4pt]
    \multicolumn{4}{l}{\textit{교차수준 상호작용}} \\
    \addlinespace[2pt]
    연결성 $\times$ 고령자 & $-1.920$ (1.901) & & \\
    WiFi $\times$ 고령자 & $3.026^{*}$ (1.515) & & \\
    자치구 소득 $\times$ 고령자 & $2.358$ (1.532) & & \\
    연결성 $\times$ 저학력 & $10.936^{***}$ (3.370) & & \\
    \bottomrule
    \multicolumn{4}{p{0.97\textwidth}}{\footnotesize\textit{Note.} Unstandardized coefficients with standard errors in parentheses. Model 0 = null model (intercept only). Reference categories: male, college education, white-collar job, apartment, multi-person household. $^{\dagger}p<.10$, $^{*}p<.05$, $^{**}p<.01$, $^{***}p<.001$.}
  \end{tabularx}
\end{table*}

\subsection{플랫폼 텍스트 분석: LLM 기반 디지털 연결성 사례분석}
\label{subsec:3-3}

앞선 \ref{subsec:3-1}절과 \ref{subsec:3-2}절에서는 행정 데이터와 설문 조사를 이용해 서울의 디지털 격차를 정량적으로 측정하고, 개인 특성과 지역 인프라의 역할을 분리하였다. 그러나 정량적 지표만으로는 주민들이 일상에서 언제, 어떤 맥락에서 디지털 연결의 결여를 경험하는지까지 포착하기 어렵다. 디지털 연결성에 대한 생활 세계의 체감과 구조적 지표 사이의 거리를 좁히기 위해 정성적 증거를 통합할 필요가 있다.

본 절의 목적은 지역 기반 플랫폼의 비정형 텍스트를 이용해 UMC의 6개 차원(Connectivity, Availability for Use, Affordability, Devices, Skills, Safety)에 잠재된 디지털 결여 사례를 체계적으로 추출하고, 이러한 사례를 \ref{subsec:3-1}절과 \ref{subsec:3-2}절의 정량 분석과 결합하는 방법론적 틀을 제시하는 데 있다. 플랫폼 자체를 분석의 대상으로 삼기보다는, 플랫폼이 제공하는 지역 단위의 생활 경험 데이터를 통해 구조와 경험 사이의 연결을 복원하고자 한다.

\subsubsection{당근 플랫폼 특성}
\label{subsubsec:3-3-1}

비정형 텍스트의 출처는 한국의 지역 기반 커뮤니티 플랫폼 '당근'의 동네생활 게시글이다. 당근을 선정한 이유는 세 가지다. 첫째, 각 게시글에 행정동과 자치구 수준의 지리적 메타데이터가 자동으로 부여되어 \ref{subsec:3-1}절에서 구축한 UMC 지수와 직접적으로 연동할 수 있다. 둘째, 플랫폼 이용자가 서울의 25개 자치구 전역에 골고루 분포해 있어 특정 지역이나 연령층에 치우치지 않는다. 셋째, 게시글은 일상생활 속 디지털 문제를 자발적으로 기록하기 때문에 조사자가 설문에서 미처 포착하지 못한 생활상의 결여를 드러낼 수 있다. 연구의 초점은 플랫폼 자체에 대한 평가가 아니라, 이와 같은 구조적 장점을 활용해 생활 세계의 경험을 수집하는 데 있다.

\subsubsection{왜 LLM인가: 기존 방법론의 한계와 LLM의 구조적 이점}
\label{subsubsec:3-3-2}

UMC의 여섯 차원에 따라 비정형 텍스트를 분류하려면 두 가지 요구를 충족해야 한다. 첫째, 게시글이 디지털 연결성 문제와 관련된 지 여부와 관련된 경우 해당 차원을 정확히 판별해야 한다. 둘째, 일상 언어로 쓰인 구어체 한국어는 조사나 어미 변화가 복잡한 교착어적 특성으로 인해 단순한 형태소 분석기만으로는 정확히 처리하기 어렵다. 기존의 토픽 모델링이나 키워드 기반 분류는 잠재 주제를 파악하는 데 유용하지만 사전 정의된 차원과 일치하지 않거나, 대화체 텍스트를 제대로 분석하지 못하는 한계가 있다. 이러한 한계는 한국어 형태소 분석기가 구어체 데이터에서 정확도가 낮고, 데이터 훈련이 필요한 머신러닝 모델이 충분한 라벨을 요구한다는 점에서 더욱 두드러진다.

대규모 언어 모델(LLM)은 이런 제약을 상당 부분 극복한다. 다국어로 사전 학습되어 한국어의 형태소 분석이 없어도 맥락을 파악할 수 있으며, 104개 이상의 언어를 동시에 다루는 모델은 언어간 의미론적 표현을 공유해 적은 양의 예시만으로도 새로운 태스크에 적응하는 능력을 갖추고 있다. LLM을 도구로 사용할 때는 출력의 확률성과 편차를 고려해야 한다. 동일한 입력에 대해서도 결과가 달라질 수 있으므로, 본 연구는 분석 절차를 고정하고 LLM은 분류와 추론을 수행하는 계산 모듈로만 활용했다. 이를 통해 연구자가 정의한 분류 규칙과 추론 프로토콜을 일관되게 적용하면서도 LLM의 문맥 이해력을 활용할 수 있도록 설계하였다.

\subsubsection{분석 구조}
\label{subsubsec:3-3-3}

텍스트 분석은 네 단계로 구성된다. 먼저, 크롤링한 원시 CSV 데이터에서 삭제되어 접근할 수 없는 게시글과 광고성\textperiodcentered{}중복 게시글을 제거한 뒤, 여섯 차원별 키워드 사전을 이용해 UMC 후보 게시글을 추출한다. 키워드 사전은 디지털 연결성과 직간접적으로 관련된 단어를 정리한 것으로 해당 단어를 포함하는 게시글만 1차로 추출하였다. 이후, 게시글이 디지털 연결성 문제와 관련이 있는지 여부를 판정하는 1차 분류 (umc\_related = Y/N/?)를 수행하며, 판정의 범위를 넓히고자 LLM에게 예시, 제한을 상세하게 밝히지 않고 분석하였다. 그 다음, 관련성이 확인된 게시글만을 대상으로 사용자가 제공한 상세한 예시를 포함한 few-shot 프롬프트를 LLM에게 제시해 게시글을 6개 개념적 차원의 사례로 분류한다. 이 단계에서는 모델의 자율성을 제한하여 프롬프트에 예시를 포함한 구체적인 분류 기준과 명세만을 제공했고, 모델이 판단에 참조하는 정보 공간인 컨텍스트를 엄격하게 관리하였다.

표~\ref{tab:5-text-sample}는 분석 구조와 표본 추출 과정을 요약한 것이다. 원시 데이터 1,287,761건 중 약 51.5\%는 게시자 탈퇴나 차단으로 본문을 볼 수 없어 분석에서 제외되었다. 중복\textperiodcentered{}광고성 게시글 제거 후에는 약 48만 건이 남았으며, 여섯 차원별 키워드 사전을 적용한 결과 26,742건이 UMC 관련 후보로 추출되었다. few-shot 프롬프트를 이용한 1차 판정에서는 약 57.7\%(15,438건)가 관련(Y)으로, 29.6\%(7,921건)가 비관련(N)으로, 12.6\%(3,383건)가 판단 유보(?)로 분류되었다. 이 가운데 Y로 판정된 게시글만을 대상으로 차원 분류를 진행했다.

차원 분류 판정에서 LLM에게 각 차원에서 대표적인 사례를 담은 3\,$\sim$\,4개의 예시를 프롬프트로 제공했다. 이러한 few-shot 프롬프트는 모델이 사람의 판단 기준을 학습하는 데 도움이 되며, 복잡한 작업에서 zero-shot보다 나은 성능을 보인다~\textbf{[인용 필요]}. 예시에는 어떤 유형의 게시글이 디지털 연결성 문제와 무관한지를 보여주는 음성 예시도 포함해, 관련성을 최대한 좁혀 높은 정밀도를 확보하도록 했다. 2차 차원 분류에서는 구체적인 분류 기준과 정의만을 프롬프트로 제공해 모델이 보다 자율적으로 사례를 분류하도록 설계했다. 이처럼 앞 단계는 예시 기반 few-shot 방식이고, 뒤 단계는 규칙 기반 프롬프트를 사용한다는 점에서 추론과 few-shot 사이의 균형을 보여준다(그림~\ref{fig:3-pipeline-roles} 참조).

\begin{table}[htbp]
  \centering\small
  \caption{텍스트 분석 표본 추출 및 판정 절차}
  \label{tab:5-text-sample}
  \begin{tabularx}{\linewidth}{l Y r}
    \toprule
    \textbf{단계} & \textbf{설명} & \textbf{남은 게시글 수} \\
    \midrule
    원시 데이터 & 2024년 5월\,$\sim$\,2026년 3월에 작성된 동네생활 게시글 전수 & 1,287,761 \\
    접근가능 게시글 & 탈퇴\textperiodcentered{}차단으로 본문 확인이 불가능한 게시글 제거 & 624,564 \\
    중복\textperiodcentered{}광고 제거 & 중복된 내용, 광고\textperiodcentered{}홍보 게시글 제거 & 480,972 \\
    UMC 키워드 필터 & 6개 차원별 키워드 사전을 적용해 후보 추출 & 26,742 \\
    1차 판정 & few-shot 프롬프트로 UMC 관련성 분류 & --- \\
    \bottomrule
  \end{tabularx}
\end{table}

\fig{image18.png}{분류와 추론의 역할 구분}{fig:3-pipeline-roles}

이후, 분류 결과를 자치구 단위로 집계하여 표출 프로파일을 만들고, \ref{subsec:3-1}절의 정량 지표와 결합하기 위해 베이지안 업데이트를 적용한다. 정량 지표와 텍스트 분류 결과를 결합하기 위해 베이지안 업데이트를 적용하였다. 사전 분포는 UMC의 6가지 개념의 잠재적 공간이다. 차원 $d$에서 자치구 $j$의 UMC 실측값을 $s_{jd}$로 할 때, 사전 분포는 $\mathrm{Beta}(\alpha_{jd}, \beta_{jd})$이며, 평균이 $s_{jd}$가 되도록 사전 분포를 가정한다. 강도 모수 $\kappa$를 사용해 $\alpha_{jd} = s_{jd}\cdot\kappa$, $\beta_{jd} = (1-s_{jd})\cdot\kappa$로 표현한다. 텍스트 분류에서 해당 차원으로 판정된 게시글 수를 $y_{jd}$, 총 분류된 게시글 수를 $n_{jd}$라 할 때, 우도는 $\mathrm{Binomial}(n_{jd}, \theta_{jd})$로 표현된다. 사후 분포는 $\theta_{jd}\mid y_{jd} \sim \mathrm{Beta}(\alpha_{jd}+y_{jd}, \beta_{jd}+n_{jd}-y_{jd})$로 계산된다. 사전 강도 $\kappa$는 기본값 20을 사용하고 10과 30에 대한 민감도 분석도 수행하였다. 게시글 수가 적은 자치구에서 사후 분포가 사전 분포에 의해 지나치게 지배되는 것을 막기 위해, 게시글 수 $n_{jd}$를 생활인구로 보정한 $n_{jd}^{*} = n_{jd}\times(\text{생활인구}_{j}/\overline{\text{생활인구}})$로 사용하였다.

\begin{align}
\alpha_{jd} &= \kappa \times s_{jd} \label{eq:alpha} \\
\beta_{jd}  &= \kappa \times (1 - s_{jd}) \label{eq:beta} \\
\theta_{jd}\mid y_{jd} &\sim \mathrm{Beta}\bigl(\alpha_{jd}+y_{jd},\, \beta_{jd}+n_{jd}-y_{jd}\bigr) \label{eq:posterior} \\
n_{jd}^{*} &= n_{jd}\times\bigl(\text{생활인구}_{j}/\overline{\text{생활인구}}\bigr) \label{eq:popadj}
\end{align}

베이지안 집계에서 게시글로 추정된 사후 분포가 사전 분포로 유의미하게 이동한 자치구\textperiodcentered{}차원 조합을 대상으로 인과기제에 관한 정성적 추론을 수행하였다. 구체적으로는 (1) 사후 이동이 통계적으로 유의하며, (2) 해당 자치구\textperiodcentered{}차원에서 Y 판정 게시글이 최소 5건 이상이고, (3) \ref{subsec:3-2}절의 다층모형 결과가 허용하는 범위 내에서만 해석을 실시하였다. 사전 분포에 의해 도출된 조합에 대해 대표 게시글 3\,$\sim$\,5건을 선정하고, 1\,$\sim$\,2문장으로 구조적 지표와 생활 경험 사이의 연결 경로를 서술하였다. 분류와 추론 작업은 세 명의 독립된 에이전트가 서로 다른 추론 경로를 통해 수행했으며, 최종 판정은 네 번째 에이전트가 세 결과를 비교\textperiodcentered{}검증하는 방식으로 이루어졌다. 이는 한 가지 프로토콜에 의존하지 않고 여러 관점의 상호 검증을 통해 편향을 통제하는 삼각검증(triangulation) 접근을 따른 것이다.

모든 분류와 추론 작업은 세 명의 독립된 분석 에이전트가 서로 다른 추론 방식 프롬프트와 컨텍스트를 사용해 수행하고, 네 번째 분석자가 세 결과를 비교\textperiodcentered{}검증하는 방식으로 이루어졌다. 이는 질적 연구에서 여러 자료원과 분석가를 동원해 편향을 줄이는 삼각검증(triangulation) 접근에 기반한다. 세 관점이 일치할 때만 해당 게시글을 최종 분석에 포함했으며, 해석상의 불확실성은 결과 표기에 반영하였다. 이러한 과정을 정리한 결과는 표~\ref{tab:6-units}, 그림~\ref{fig:3-pipeline-full}과 같다.

\begin{table}[htbp]
  \centering\small
  \caption{\ref{subsec:3-3}절의 분석 단위와 해석 단위}
  \label{tab:6-units}
  \begin{tabularx}{\linewidth}{l l Y Y}
    \toprule
    \textbf{구분} & \textbf{단위} & \textbf{무엇을 보여주는가} & \textbf{해석상 주의점} \\
    \midrule
    관찰 단위 & 개별 게시글 & 디지털 연결성 경험의 단서 & 단일 게시글만으로 지역 구조를 주장할 수 없음 \\
    패턴 단위 & 행정동 게시글 & 차원별 표출 비율과 분포 & 집계 기준과 표본 편향을 함께 검토해야 함 \\
    구조 단위 & UMC 지역 지표 & 지역 수준의 연결성 조건 & 텍스트 집계 결과와의 정합을 통해 제한적으로 해석 \\
    인과기제 단위 & 추론 메커니즘 & 사례와 지역 패턴의 연결 설명 & \ref{subsec:3-2}의 정량 결과가 허용하는 범위 안에서만 해석 \\
    \bottomrule
  \end{tabularx}
\end{table}

\subsubsection{데이터와 적용 범위}
\label{subsubsec:3-3-4}

원시 데이터는 서울 25개 자치구를 대상으로 2024년 5월부터 2026년 3월까지 작성된 동네생활 게시글 전수를 크롤링한 것으로, 중복 제거 후 총 1,287,761건이다. 이 중 삭제(DELETED) 663,163건(51.5\%)과 차단(BLOCKED) 103,728건(8.1\%)을 합한 766,891건(59.6\%)은 본문 확인이 불가능하여 분석에서 제외하였다. 분석 대상은 최종 약 52만 건이다.

표본 탈락의 체계적 편향 가능성을 검토한 결과, 자치구별 제거율의 표준편차는 4.6\%로 자치구 간 편차가 크지 않았으며, 디지털 결여 추론의 주요 대상인 생활/편의(47.0\%), 동네사건사고(48.4\%), 병원/약국(35.5\%) 카테고리는 사교\textperiodcentered{}홍보 목적 카테고리(동네친구 86.6\%, 교육 69.0\%)에 비해 제거율이 낮았다. 삭제의 상당 부분은 사용자 탈퇴에 수반된 일괄 삭제에 기인하는 것으로 판단된다. 다만, 삭제 게시글 중 일부가 디지털 결여 관련 사례를 포함할 가능성을 완전히 배제할 수는 없고, 삭제 행위 자체가 디지털 결여와 체계적으로 연관되었을 수 있다. 이는 본 분석의 한계에 해당한다.

당근 데이터의 구조적 한계도 인식할 필요가 있다. 플랫폼 이용자 자체가 디지털 접근이 가능한 집단으로 한정되므로, 디지털 소외 계층의 경험은 구조적으로 과소 대표된다. 게시글의 시간적\textperiodcentered{}공간적 분포가 균일하지 않아 자치구 간 비교에서 표본 편향의 가능성이 존재한다. 이러한 한계는 데이터 소스의 제약이지 방법론의 제약이 아니다. 보다 직접적인 질적 자료가 존재하는 맥락에서는 동일한 파이프라인을 적용하여 보다 정밀한 분석이 가능할 것이다.

\figfull{image19.png}{LLM 기반 텍스트 분석 파이프라인}{fig:3-pipeline-full}

\subsubsection{분석 결과}
\label{subsubsec:3-3-5}

텍스트 분류 결과는 표~\ref{tab:6-units}와 같다. 연결 품질(Connection Quality) 차원은 6개 차원 가운데 전이의 폭이 가장 작다($-0.006 \sim +0.005$). 실측과 텍스트 증거가 대체로 수렴하는 차원으로 볼 수 있다. 양천구는 실측 UMC 0.278로 낮은 편이나 관련 게시글 비율이 상대적으로 높아 $+0.005$의 전이를 보였고, 용산구는 실측 0.884로 높은 반면 관련 게시글이 드물어 $-0.006$의 전이를 보였다.

이용 가능성(Availability for Use) 차원에서는 도봉구가 실측 최하위(0.213)에 해당하면서도 관련 게시글 비율이 높아 $+0.006$의 전이를 보였다. 광진구는 실측 0.832로 높고 관련 게시글이 드물어 $-0.010$의 전이를 기록하였다.

경제적 접근성(Affordability) 차원은 서울특별시 자치구의 소득 수준과 일관된 방향의 경향을 보였다. 중랑구와 강서구는 각각 $+0.005$의 양의 전이를 보인 반면, 서초구는 실측이 0.993으로 최고 수준이면서 관련 게시글이 드물어 $-0.007$의 전이를 나타냈다. 이 차원은 통신비와 기기 구입 비용에 관한 고민이 소득 수준에 따라 차등적으로 표출되는 경향과 정합적이다.

기기 접근성(Devices) 차원은 자치구 간 편차가 가장 큰 차원이다(gap = 0.090). 중랑구는 실측 0.052에 해당하면서도 기기 관련 게시글 비율이 4.95\%로 전 자치구 가운데 가장 높아 $+0.048$의 전이를 보였다. 송파구 또한 $+0.026$의 유사한 패턴을 나타냈다. 반면 용산구($-0.042$)와 동작구($-0.033$)는 실측이 높으면서 관련 게시글이 드물게 나타났다. 다만 이 차원에서 실측 지표는 기기 보유율을 포착하는 반면, 텍스트 증거는 중고 거래, 수리 요청 등 기기와 관련한 일상적 거래와 고충을 포착하므로, 두 지표가 지시하는 대상이 부분적으로 상이할 가능성을 배제하기 어렵다.

디지털 역량(Digital Skills) 차원은 기기 차원에 이어 두 번째로 편차가 큰 차원이다(gap = 0.092). 중랑구는 실측 0.242에 관련 게시글 비율 4.8\%로 $+0.047$의 전이를 보였다. 강남구 또한 실측이 0.276으로 낮은 편이면서 관련 게시글 비율이 상대적으로 높아 $+0.013$의 전이를 나타냈다. 반면 관악구는 실측 0.795로 최고 수준이면서 관련 게시글이 드물어 $-0.045$의 전이를 보였다.

디지털 역량(Digital Skills) 차원에서는 기기 차원에 이어 두 번째로 큰 편차가 나타났다. 중랑구는 실측 0.242에 비해 관련 게시글 비율 4.8\%로 $+0.047$의 이동을 보였으며, 강서구는 실측값이 높음에도 관련 게시글이 적어 $-0.031$의 이동을 보였다. 안전\textperiodcentered{}보안(Safety) 차원은 피싱 및 보이스피싱 사례가 게시글의 주요 내용을 이루었고, 중랑구($+0.018$)와 도봉구($+0.015$)에서 양의 이동이 확인되었다. 종로구는 실측값이 매우 높고 관련 게시글이 적어 사후 분포가 낮아졌다.

안전\textperiodcentered{}보안(Safety \& Security) 차원에서는 중랑구($+0.018$)와 도봉구($+0.015$)에서 피싱 사기 및 보이스피싱 관련 게시글의 비율이 높았다. 종로구는 실측 0.984로 최고 수준이면서 관련 게시글이 드물어 $-0.019$의 전이를 나타냈는데, 관광\textperiodcentered{}상업 중심지의 주거 구성과 당근 이용자 집단의 특성이 상이할 가능성을 함께 고려할 필요가 있다.

자치구 단위로 집계했을 때, 지역의 실측값과 게시글 분포 간 차이는 크게 세 가지 유형. 첫째, 평균 전이가 양의 방향으로 나타나는 자치구이다. 중랑구는 6개 차원 중 4개 차원에서 양의 전이가 가장 크게 나타나 평균 $+0.020$으로 가장 높은 값을 보였다. 도봉구와 강남구 또한 평균적으로 양의 방향을 보였으나 그 크기는 0.001 수준으로 작았다. 특히 강남구의 경우 평균 소득이 높은 자치구임에도 이용 가능성과 디지털 역량 차원에서 양의 전이가 관찰되었는데, 자치구 내부의 이질적 집단이 관련 이슈를 표출하고 있을 가능성을 시사한다. 둘째, 전이가 0에 가까운 자치구이다. 강북구($-0.000$), 강동구($-0.001$), 동대문구($-0.001$)가 이에 해당한다. 이들 자치구에서는 차원별로 양\textperiodcentered{}음의 전이가 섞여 있어 평균적으로 상쇄되는 양상을 보인다. 셋째, 평균 전이가 음의 방향으로 나타나는 자치구이다. 마포구($-0.014$), 영등포구($-0.013$), 관악구($-0.013$), 서초구($-0.009$)가 이에 해당한다. 관악구는 대학가 중심의 인구 구성과 디지털 역량 실측값이 맞물려 도움 요청 성격의 게시글이 드물게 나타난 것으로 볼 수 있으며, 서초구는 경제적 접근성과 기기 차원에서 특히 큰 음의 전이를 보였다. 사전 강도 $\kappa$를 10, 20, 30으로 달리한 민감도 분석에서도 사후 평균의 차이는 $0.00001\sim0.00010$ 수준으로 실질적 변화가 없었다. 이는 생활인구 보정 후 유효 관측 수가 충분해 사전의 가중치 변화가 결과에 큰 영향을 미치지 않음을 시사한다.

해석에서 주의할 점은 양의 전이가 곧바로 실측 지표가 과소평가되었음을 의미하지 않는다는 것이다. 양의 이동은 실측 지표가 포착하지 못한 생활상의 불편을 반영할 수 있지만, 특정 자치구에서 특정 문제에 대한 논의가 활발한 문화적\textperiodcentered{}인구학적 요인을 동시에 포함할 수 있다~\textbf{[인용 필요]}. 또한 플랫폼 이용자의 표집 편향과 게시글의 시공간적 불균형이 결과 해석에 영향을 미칠 수 있다. 따라서 사후 분포의 방향과 크기를 정책적 함의로 연결할 때는 정량\textperiodcentered{}정성적 지표를 종합적으로 고려해야 한다. 본 절의 결과에서 유의미한 것은 전이의 부호와 크기가 차원과 자치구에 따라 체계적으로 달라진다는 것이다. 단순히 당근마켓 게시글의 빈도만을 집계할 경우 플랫폼 활성도 자체가 UMC 수준에 대한 추정과 혼동될 우려가 있다. 사전 분포로 실측값을 투입하고 텍스트 분류 결과를 우도로 결합하는 절차는, 관측된 텍스트 빈도에서 플랫폼 활성도에 기인하는 부분과 UMC 수준에 기인하는 부분을 부분적으로 분리하는 효과를 갖는다. 이 점에서 사전-사후 분포의 전이는 텍스트 증거가 실측을 얼마나 보정하는지를 보여주는 지표로 기능한다.

다만 shift의 방향을 해석할 때 두 가지 제약을 함께 고려해야 한다. 첫째, 양의 shift가 관찰된다고 해서 해당 자치구의 UMC 수준이 실측보다 낮다고 곧바로 결론지을 수 없다. 양의 shift는 실측 지표가 포착하지 못하는 생활 수준의 체감이 존재할 가능성을 시사하지만, 관련 주제가 당근마켓이라는 특정 플랫폼에서 표출되는 방식 자체가 자치구별 이용자 구성에 의존하기 때문이다. 둘째, 음의 shift 또한 해당 자치구의 UMC 수준이 실측보다 높다는 의미로 해석하기보다, 관련 이슈가 플랫폼 상에서 표출될 유인이 상대적으로 적은 상태로 읽는 것이

\begin{table*}[htbp]
  \centering\small
  \caption{자치구별 게시글 판정 결과}
  \label{tab:7-district-cls}
  \begin{tabularx}{\textwidth}{l r r r r r}
    \toprule
    \textbf{자치구} & \textbf{총 게시글(n)} & \textbf{Y 판정} & \textbf{? 판정} & \textbf{N 판정} & \textbf{Y 비율(\%)} \\
    \midrule
    \textbf{중랑구}   & 5,476  & 794 & 28  & 4,654  & 14.5 \\
    \textbf{송파구}   & 9,367  & 860 & 99  & 8,408  & 9.18 \\
    \textbf{서초구}   & 6,170  & 434 & 216 & 5,520  & 7.03 \\
    \textbf{은평구}   & 6,140  & 429 & 50  & 5,661  & 6.99 \\
    \textbf{도봉구}   & 3,491  & 233 & 5   & 3,253  & 6.67 \\
    \textbf{서대문구} & 4,500  & 289 & 5   & 4,206  & 6.42 \\
    \textbf{강북구}   & 3,885  & 247 & 31  & 3,607  & 6.36 \\
    \textbf{용산구}   & 3,012  & 181 & 12  & 2,819  & 6.01 \\
    \textbf{양천구}   & 4,975  & 296 & 7   & 4,672  & 5.95 \\
    \textbf{강동구}   & 6,254  & 342 & 231 & 5,681  & 5.47 \\
    \textbf{영등포구} & 5,964  & 306 & 5   & 5,653  & 5.13 \\
    \textbf{동대문구} & 4,754  & 238 & 29  & 4,487  & 5.01 \\
    \textbf{성북구}   & 5,437  & 266 & 141 & 5,030  & 4.89 \\
    \textbf{중구}     & 1,351  & 63  & 16  & 1,272  & 4.66 \\
    \textbf{강서구}   & 8,174  & 371 & 236 & 7,567  & 4.54 \\
    \textbf{금천구}   & 2,298  & 101 & 7   & 2,190  & 4.40 \\
    \textbf{종로구}   & 2,267  & 97  & 27  & 2,143  & 4.28 \\
    \textbf{동작구}   & 3,675  & 150 & 22  & 3,503  & 4.08 \\
    \textbf{강남구}   & 10,047 & 404 & 459 & 9,184  & 4.02 \\
    \textbf{성동구}   & 404    & 16  & 9   & 379    & 3.96 \\
    \textbf{관악구}   & 11,665 & 407 & 75  & 11,183 & 3.49 \\
    \textbf{마포구}   & 7,974  & 242 & 254 & 7,478  & 3.03 \\
    \textbf{구로구}   & 4,082  & 121 & 96  & 3,865  & 2.96 \\
    \textbf{광진구}   & 5,440  & 147 & 42  & 5,251  & 2.70 \\
    \textbf{노원구}   & 5,040  & 103 & 46  & 4,891  & 2.04 \\
    \bottomrule
  \end{tabularx}
\end{table*}

\figfull{image20.png}{지역별 사전-사후 분포 차이}{fig:3-prior-posterior}

안전하다. 중랑구와 마포구의 대조에서 드러나듯, shift의 해석은 자치구의 인구 구성과 플랫폼 이용 특성을 비롯한 다양한 맥락 정보와 결합될 때 의미를 띈다. 둘째, 음의 전이 또한 해당 자치구의 UMC 수준이 실측보다 높다는 의미로 해석하기보다, 관련 이슈가 플랫폼 상에서 표출될 유인이 상대적으로 적은 상태로 읽는 것이 타당할 수 있다. 이런 측면에서 지역 자원이 실제로 결여되어 있음에도 플랫폼 상에서 논의되지 않는 이유를 고려하는 것 또한 중요하다.

\begin{table*}[htbp]
  \centering\small
  \caption{결여 유형별 가설 생성 결과 분포}
  \label{tab:8-deficit-types}
  \begin{tabularx}{\textwidth}{l Y r r}
    \toprule
    \textbf{결여 유형} & \textbf{관련 UMC 차원} & \textbf{관련 게시글 수 (n=100)} & \textbf{비율 (\%)} \\
    \midrule
    정보부재         & Digital Skills      & 63 & 63.0 \\
    제도적 경로 결여 & Available for Use   & 61 & 61.0 \\
    기기 접근 결여   & Devices             & 50 & 50.0 \\
    디지털 역량 부재 & Digital Skills      & 45 & 45.0 \\
    연결 인프라 결여 & Connectivity        & 40 & 40.0 \\
    공공 접근점 결여 & Available for Use   & 32 & 32.0 \\
    경제적 접근 결여 & Affordability       & 27 & 27.0 \\
    안전\textperiodcentered{}보안 결여 & Safety & 16 & 16.0 \\
    복합 결여 (2개 이상) & ---             & 88 & 88.0 \\
    \bottomrule
    \multicolumn{4}{p{0.97\textwidth}}{\footnotesize\textit{Note.} 하나의 게시글에서 복수의 결여 유형이 추론될 수 있으므로 비율의 합은 100\%를 초과한다. Phase 1은 가설 생성 단계이며, Phase 2(판정)에서 삼각검증을 거쳐 최종 분류가 확정된다.}
  \end{tabularx}
\end{table*}

\subsection{분석 결과 종합}
\label{subsec:3-4}

본 절에서는 \ref{subsec:3-1}절(UMC 지수 구축), \ref{subsec:3-2}절(HLM 다수준 분석), \ref{subsec:3-3}절(LLM 기반 텍스트 분석)의 결과를 종합하여, 서울의 디지털 격차가 어떤 구조로 작동하는지를 통합적으로 해석한다. 세 분석은 각각 '어디에 격차가 있는가'(측정), '왜 격차가 발생하는가'(인과 규명), '격차가 어떻게 경험되는가'(질적 맥락)에 답하며, 이들의 교차 지점에서 정책적으로 유의미한 함의가 도출된다.

세 분석이 공통으로 지지하는 첫 번째 발견이자 본 프로젝트의 핵심 실증 결과는, 서울의 디지털 격차가 자치구 간 차이가 아닌 개인 간 차이에 의해 구조화되어 있다는 것이다. \ref{subsec:3-2}절의 HLM 분석에서 ICC는 0.49\%에 불과하여, 디지털 활용 점수 분산의 99.5\%가 동일 자치구 내 개인 간 차이에 기인하였다. 이 결과는 \ref{subsec:3-1}절에서 구축한 UMC 종합지수의 자치구 간 편차가 실재함에도 불구하고, 해당 편차가 개인의 디지털 활용 수준에 미치는 직접적 영향은 제한적임을 의미한다. 그러나 이러한 낮은 ICC는 자치구 수준 변수의 무의미함을 뜻하는 것이 아니라, 그 효과가 집단별로 차별적으로 발현됨을 시사한다. Model 3의 교차수준 상호작용 결과가 이를 뒷받침한다.

\ref{subsec:3-3}절의 텍스트 분석은 이 양적 발견이 일상에서 어떻게 발현되는지를 구체화한다. '정보부재'가 가장 빈번한 결여 유형(63\%)으로 나타난 것은, 교육수준이 낮을수록 디지털 서비스의 존재 자체를 인지하지 못하는 정보 격차가 더 심화될 수 있음을 시사한다. 프린터 이용을 위해 편의점 무인 인쇄를 모르는 사례, ISP 고객센터 번호를 모르는 사례, 정부24의 공문서 출력 기능을 모르는 사례 등은 모두 '서비스가 존재하지만 이용자에게 보이지 않는' 구조적 문제를 가리킨다. 이는 HLM에서 교육 기울기가 가파른 이유를 경험적 수준에서 설명하는 보완적 증거로 해석할 수 있다.

두 번째 발견은, 자치구 인프라의 효과가 주효과 수준에서는 유의하지 않으나 특정 취약집단에 대해서는 조절 효과로서 유의하게 작용한다는 것이다. Model 2에서 UMC 차원 중 Connectivity와 관련된 인프라 변수의 주효과는 통계적 유의성에 도달하지 못하였는데($p = 0.237$), 이는 ICC가 0.49\%인 상황에서 자치구 수준 변수의 주효과가 제한적인 것과 일관된 결과이다. 그러나 Model 3에서 교차수준 상호작용이 유의하게 나타남으로써, 인프라의 효과가 평균적 효과가 아닌 집단별 조절 효과로서 작동함이 확인되었다.

\ref{subsec:3-3}절의 텍스트 분석은 이 차별적 효과의 질적 맥락을 제공한다. 공공 접근점(프린트 카페, 도서관 복합기, 무인 인쇄)에 대한 정보부재가 32\%의 게시글에서 추론된 것은, 인프라가 물리적으로 존재하더라도 이용자의 인지와 접근 역량이 수반되지 않으면 실질적 혜택으로 전환되지 않음을 보여준다. 은평구 신사동의 사례('정녕 없는 건지요')는 온라인 검색을 이미 시도했으나 결과가 불충분했을 가능성을 시사하며, 이는 인프라의 존재와 인프라의 '발견가능성(discoverability)' 사이의 간극을 조명한다. HLM에서 인프라 주효과가 유의하지 않았던 결과와 결합하면, 인프라 투자의 효과는 취약계층의 정보 접근과 역량 지원이 병행될 때 비로소 발현된다는 해석이 가능하다.

세 번째 발견은, 디지털 격차가 단일 축이 아닌 복합 취약성의 교차로 심화된다는 것이다. HLM에서 독거 고령자는 비독거 고령자 대비 추가적으로 4.0점이 낮았고, 연령 증가에 따른 저학력의 부정적 효과가 가속화되는 패턴이 확인되었다. \ref{subsec:3-3}절에서도 결여의 복합성이 일관되게 나타났다. 게시글당 평균 2.5개 이상의 결여 유형이 동시에 식별되었으며, 정보부재\textperiodcentered{}역량부재\textperiodcentered{}제도적 경로 차단이 중첩될 때 이용자가 공식 경로에서 이탈하여 커뮤니티에 의존하는 구조가 확인되었다. 디지털 지원망의 부재('자녀분들이 없으시거나 떨어져 살아서 묻기 힘들거나')는 가구 구성이 디지털 역량 획득의 비공식적 경로로 기능함을 보여주며, 독거 고령자의 추가적 격차를 설명하는 기제로 해석된다.

네 번째 발견은, 디지털 사막의 공간적 군집과 차원별 이질성이다. \ref{subsec:3-1}절에서 하위 5개 지역 중 특히 서울 북부(도봉구, 강북구, 노원구, 중랑구)가 UMC 종합지수의 Low-Low 클러스터를 형성하였으며, 이들은 Connectivity 차원에서는 도시 평균 이상을 기록하면서도 Devices와 Safety 차원에서 극단적 결핍을 보였다. 이러한 차원별 프로파일의 이질성은 일률적 인프라 투자보다 차원별 맞춤 개입의 필요성을 시사한다.

방법론적 관점에서, 세 분석의 결합은 각각의 한계를 상호 보완한다. \ref{subsec:3-1}절의 종합지수는 자치구 간 상대적 위치를 파악하되 개인 내 변이를 포착하지 못하며, \ref{subsec:3-2}절의 HLM은 개인과 지역의 기여를 분리하되 디지털 격차의 질적 발현 경로를 설명하지 못한다. \ref{subsec:3-3}절의 텍스트 분석은 결여의 경험적 맥락을 가시화하되 대표성에 한계가 있다. 세 분석을 병렬이 아닌 순차적 삼각검증(측정 $\rightarrow$ 인과 규명 $\rightarrow$ 질적 맥락)으로 결합함으로써, '어디에'(자치구), '왜'(개인 특성과 인프라의 상호작용), '어떻게'(결여의 발현 경로)라는 세 질문에 대한 통합적 답을 구성할 수 있었다.

이상의 종합 결과는 4장의 정책 제언으로 직결된다. 개인 특성이 지배적인 설명변수라는 발견은 people-based 정책의 우선성을, 인프라의 차별적 효과는 취약계층 타겟형 place-based 정책의 필요성을, 복합 취약성의 확인은 단일 프로그램이 아닌 통합 지원 체계의 필요성을, 공간적 군집의 확인은 권역 단위 거버넌스의 근거를 각각 제공한다.
)
\section{분석}
\label{sec:analysis}

\noindent{\sffamily\small\color{muted}\textbf{CHAPTER 3}}

\begin{quoteblock}
측정 $\rightarrow$ HLM $\rightarrow$ 베이지안 통합의 3단계 분석으로 서울 디지털 격차의 공간 분포, 개인-지역 기여 분리, 그리고 텍스트 기반 잠재 격차 신호까지 통합 추정한다.
\end{quoteblock}

\figfull{image6.png}{핵심 발견 요약}{fig:3-keyfindings}
\figfull{image7.png}{디지털 사막 --- 종합지수 하위 4개 자치구}{fig:3-desert}

본 장의 세 분석은 순차적으로 연결된다. 먼저 서울 자치구별 UMC 지수를 구축하여 지역 격차를 측정하고(\ref{subsec:3-1}), 이 결과를 지역 수준 변수로 활용하여 개인 및 지역 특성이 개인의 디지털 연결성에 미치는 영향을 분석한다(\ref{subsec:3-2}). 마지막으로 디지털 플랫폼 텍스트를 통해 양적 분석이 포착하지 못하는 구체적 사례를 도출한다. 이 흐름은 측정 $\rightarrow$ 인과 규명 $\rightarrow$ 현상 보완으로 구성되며, \ref{subsec:3-4}절에서 세 분석의 결과를 종합한다.

\subsection{서울의 UMC 지역 격차}
\label{subsec:3-1}

\subsubsection{측정 지표 및 데이터}
\label{subsubsec:3-1-1}

본 프로젝트에서 구성한 측정 지표는 UMC의 6개 차원에 대한 개념적 정의를 바탕으로 한다. 각 차원은 Connectivity, Available for Use, Affordability, Devices, Skills, Safety로 구분되며, 개별 측정 지표는 의미성과 보편성을 동시에 고려한 두 가지 핵심 구성 요소에 따라 설정하였다~\citep{ITU2023}. 본 분석에 활용된 서울의 UMC를 구성하는 6개 차원별 측정 지표는 표~\ref{tab:1-umc-dims}과 같다.

측정 지표 선정 과정에서는 개념적 정의와의 부합성, 서울의 디지털 기술 환경, 데이터 가용성을 종합적으로 검토하였다. 분석에는 데이터를 확보할 수 있는 가장 최근 시점인 2023년과 2024년, 총 2개 시점의 자료를 활용하였다. 2개 시점의 국가 행정 자료 등을 활용하였다. 단, Devices, Skills, Security 차원은 데이터 제공 시점의 한계로 인해 2023년 조사 자료를 기준으로 두 시점 모두 동일한 값을 적용하였다.

Connectivity의 경우, 한국전파누리에서 제공하는 이동통신 무선국 데이터와 서울시 행정데이터인 생활인구 데이터를 활용하여 측정 지표를 산출하였다. \citet{ITU2023}에 따르면, Connectivity는 네트워크 커버리지 내에 거주하는 인구를 기준으로 측정될 수 있다. 그러나 한국은 직주 분리가 뚜렷한 경우가 많아, 거주지역과 실질적으로 디지털 환경에 접속하는 지역이 상이할 수 있다. 이에 본 연구에서는 단순 거주인구 대신 주간 유동인구 데이터를 활용하여 4G 및 5G 네트워크 커버리지를 산출함으로써, 한국의 맥락에 보다 적합한 지표를 고안하였다. 구체적으로, 4G 기지국 수와 5G 기지국 수를 각각 주간 생활인구 천 명당 기지국 수로 환산하여 집계구별로 산출한 뒤, 이를 자치구별로 평균한 값을 최종 지표로 사용하였다.

Available for Use는 디지털 서비스의 실질적 이용 가능성을 측정하는 차원이다. SKT 통신정보에서 제공하는 월별 지역 내 가구당 모바일 데이터 사용량과 금융, 쇼핑, 동영상, 배달 4개 범주의 온라인 서비스 이용일수를 인구가중 평균으로 자치구별 월별로 집계하고, 12개월 평균값을 사용하였다. 여기에 자치구 지방정부의 공공WiFi 등록 현황 데이터를 활용하여 전체 인구 1,000명당 공공WiFi AP 수를 산출하여 포함하였다. 공공WiFi 밀도는 가정 내 인터넷 접속 환경이 제한적인 고령층에게 특히 중요한 디지털 접근 경로로 기능한다~\citep{vanDijk2020}.

Affordability는 ITU ICT Development Index 2023의 두 가지 축 구조를 따라 구성하였다~\citep{ITU2024}. 구매력(capacity) 측면에서는 서울시 상권분석서비스에서 제공하는 자치구별 월평균 가구소득을 활용하고, 가격부담(burden) 측면에서는 SKT 통신정보의 최근 3개월 이내 요금 연체 비율을 활용하였다. 연체율은 행정동, 성별, 연령대별 인구가중 평균으로 산출한 뒤 자치구 단위로 집계하였으며, 값이 높을수록 경제적 접근성이 낮음을 의미하는 역방향 지표이다.

Devices는 디지털 기기 보유 및 이용의 다양성을 측정한다. 서울시민 디지털역량실태조사의 가구 내 핵심 디지털 기기 보유 유형 수와 개인이 실제로 사용하는 기기 유형 수를 자치구별 가중평균으로 산출하였다. \citet{ITU2023}의 보고서에 따르면, 휴대폰과 컴퓨터 등 기기의 특성에 따라서 기기의 보유와 이용은 서로 다르게 해석된다. 이를 한국 국내 맥락에선, 컴퓨터의 보유 여부도 디지털 연결성을 드러내는 지표가 될 수 있다. 스마트폰 보급률이 사실상 포화 상태인 서울에서는 데스크톱, 노트북, 태블릿 등 보조 기기의 보유 여부가 디지털 활용 역량의 차별화 요인으로 작용한다~\citep{vanDeursenHelsper2015}.

Digital Skills는 기초, 응용, 고급, 사이버활동의 4개 하위 영역으로 구성된다. 기초 역량은 스마트 기기 기본 조작 능력을, 응용 역량은 디지털 서비스 활용 능력을, 고급 역량은 스마트 기기 활동 숙련도를, 사이버활동 역량은 사이버 공간에서의 활동 능력을 각각 측정한다. 각 하위 영역의 원점수를 만점으로 나누어 0\,$\sim$\,1 비율로 변환한 뒤 자치구별 가중평균을 산출하였다.

Safety는 디지털 보안 행동과 보안 인식의 2개 하위 지표로 구성된다. 보안 행동은 비밀번호 관리, 소프트웨어 업데이트, 2단계 인증 사용 등 실제 보안 실천 행위를 측정하고, 보안 인식은 개인정보 보호의 중요성에 대한 인지 수준을 측정한다. 각 하위 지표를 만점 기준 비율로 변환하여 자치구별 가중평균을 산출하였다.

\begin{table*}[htbp]
  \centering\small
  \caption{UMC 차원별 측정 지표}
  \label{tab:1-umc-dims}
  \begin{tabularx}{\textwidth}{l Y Y Y}
    \toprule
    \textbf{Dimension} & \textbf{Indicators} & \textbf{Data Source} & \textbf{Note} \\
    \midrule
    Connectivity & station\_density\_4g, station\_density\_5g, download\_speed & Spectrum Resource Mgmt System; NIA Communication Quality & Per 1K daytime living pop; CQ uses 2022--2024 avg \\
    \addlinespace[2pt]
    Available for Use & mobile\_data\_usage, online\_service\_days, wifi\_density & SKT Telecom Data; Local Admin WiFi License Data & Pop-weighted district avg; WiFi per 1K total pop \\
    \addlinespace[2pt]
    Affordability & avg\_income, delinq\_rate (reverse) & Seoul Commercial Analysis Service; SKT Telecom Data & ITU IDI capacity + burden; delinq = reverse-coded \\
    \addlinespace[2pt]
    Devices & device\_home\_core, device\_use\_core & Seoul Digital Competency Survey 2023 & Home: excl. smartphone; Use: incl. smartphone \\
    \addlinespace[2pt]
    Digital Skills & Basic / Applied / Advanced / Cyber activity & Seoul Digital Competency Survey 2023 & --- \\
    \addlinespace[2pt]
    Safety & safety\_behavior, safety\_awareness & Seoul Digital Competency Survey 2023 & Security practice (7 items) + Awareness (3 items) \\
    \bottomrule
  \end{tabularx}
\end{table*}

\subsubsection{측정 지표 지수화}
\label{subsubsec:3-1-2}

모든 원시 지표는 25개 자치구 내에서 Min-Max 스케일링을 통해 $[0,1]$ 구간으로 정규화한다. 정규화된 값 $x^{*}$는 식~\eqref{eq:minmax}과 같이 산출한다.

\begin{equation}\label{eq:minmax}
  x^{*} = \frac{x_{i} - x_{\min}}{x_{\max} - x_{\min}}
\end{equation}

이때, 값이 클수록 나쁜 결과를 의미하는 역방향 지표는 변환하여 모든 지표의 방향성을 통일하였다. 그 다음, 각 차원 내에서 정규화된 지표들을 동일 가중치 산술 평균으로 결합하여 차원별 점수를 산출하였다. 동일 가중치로 분석을 진행한 이유는, 6개 차원 간 상대적 중요도에 대한 합의된 기준이 부재한 상태에서 연구자의 자의적 판단을 최소화하기 위함이다~\citep{ITU2024}. 산출된 종합지수는 0(최저)에서 1(최고) 범위를 가지며, 서울 25개 자치구 전체에 대해 2023년과 2024년 각각 독립적으로 산출한다. 해석에 있어 중요한 점은 이 종합 지수가 지역의 절대적 UMC가 아닌, 서울 내 다른 자치구 간의 상대적 위치를 보여준다는 것이다.

\subsubsection{지역별 UMC 지수 특성}
\label{subsubsec:3-1-3}

25개 자치구의 2024년 UMC 종합지수는 0.279(중랑구)에서 0.695(서초구)까지 분포하였으며, 최상위와 최하위 간 약 2.5배의 격차가 존재하였다. 서울이 높은 인터넷 보급률을 가진 도시라는 일반적 인식에도 불구하고, 하위 행정단위에서는 상당한 상대적 수준의 디지털 연결성 격차가 관찰되었다. 상위 5개 자치구는 서초구(0.695), 영등포구(0.686), 마포구(0.665), 서대문구(0.615), 용산구(0.603)이다. 하위 5개 자치구는 중랑구(0.279), 도봉구(0.303), 강북구(0.380), 구로구(0.389), 노원구(0.389)이다. 상위 자치구는 대체로 도심부 또는 강남권에 위치하여 밀집된 상업 활동, 높은 소득 수준, 풍부한 민간 디지털 인프라의 혜택을 받는 반면, 하위 자치구는 북부 외곽 및 남서부에 집중되어 있었다.

종합지수의 공간적 분포를 살펴보면(그림~\ref{fig:3-umc-comp}), UMC 지수가 높게 나온 자치구는 두 구역에 집중되어 있었다. 이는 각각 서초-영등포-마포로 이어지는 한강과 도심권 지역과, 용산-종로-중구 중심의 전통적 도심부이다. 반면, 북부 외곽(강북구, 도봉구, 노원구)과 남서부(구로구, 금천구)는 도시 평균 이하의 점수를 보였다. 이러한 공간 분포는 서울의 사회경제적인 자원 분포와 연관되어 있었다. 즉, 디지털 격차가 도시의 구조적 공간 불평등과 밀접히 연관되어 있을 가능성을 시사하였다.

차원별 분석에서 한편으로, 어떤 자치구도 6개 차원 모두에서 균일하게 우수하지는 않다는 점이 발견되었다(그림~\ref{fig:3-umc-heatmap}). 예컨대 서울 북부에 위치한 은평구의 경우, Connectivity 차원에서 0.781로 상위권에 위치하지만, Affordability 차원에서는 0.328로 하위권에 그쳤다. 유사하게, 관악구는 Skills 차원에서 0.795로 1위를 기록하였지만 Affordability에서는 0.152로 최하위권이다. 이러한 다차원적 지역간 불균등은, 종합 순위만으로는 정책적으로 의미 있는 차원별 취약성을 포착할 수 없음을 의미하였다.

\begin{keymsg}[title=핵심 발견 \textperiodcentered{} 디지털 사막 4개 자치구]
종합지수 하위 자치구(중랑\textperiodcentered{}도봉\textperiodcentered{}강북\textperiodcentered{}노원)는 Devices\textperiodcentered{}Safety 차원에서 동반 결핍을 보이며, 공간적으로도 서울 북부에 집중되어 있다. 이는 권역 단위 협력 거버넌스 구축의 정당화 근거를 제공한다.
\end{keymsg}

결론적으로, 그림~\ref{fig:3-umc-dim}에 제시된 바와 같이 공간적으로 밀집한 종합지수 하위 4개 자치구(중랑구, 도봉구, 강북구, 노원구)를 본 프로젝트에서는 잠정적으로 디지털 사막(Digital Desert)으로 분류한다. 이들 자치구는 공통적으로 Devices와 Safety 차원에서 뚜렷한 결핍을 보이며, Affordability 역시 서울시 평균에 미치지 못했다. 한편, Connectivity 차원에서는 일부 자치구(중랑구 0.630, 강북구 0.562)가 서울시 평균을 상회하여, 물리적 네트워크 인프라 자체의 부재보다는 이를 활용하기 위한 인적\textperiodcentered{}경제적 역량의 부족이 핵심 격차 요인일 가능성을 보여준다. 이러한 발견은 \ref{subsec:3-2}절의 다층 분석에서 개인 수준 변수와 자치구 수준 변수의 상대적 기여를 분리 및 검증함으로써 보다 체계적으로 검토된다.

또한, UMC의 일부 차원에서 지리적으로 인접한 지역이 서로 유사한 패턴을 보였다. 그림~\ref{fig:3-umc-dim}은 차원별 지역의 공간적 자기상관성을 나타낸 것이다. 공간적으로 군집해 있는 자치구가 서로 유사한 패턴을 보이는 점은 디지털 사막 해소를 위한 광역적 거버넌스 체계 구축의 가능성을 시사한다.

결론적으로, 그림~\ref{fig:3-umc-dim}에 제시된 바와 같이 종합지수 하위 5개 자치구(중랑구, 도봉구, 강북구, 구로구, 노원구)를 본 프로젝트에서는 잠정적으로 디지털 사막(Digital Desert)으로 분류한다. 이들 자치구는 공통적으로 Devices와 Safety 차원에서 뚜렷한 결핍을 보이며, Affordability 역시 서울시 평균에 미치지 못했다. 한편, Connectivity 차원에서는 일부 자치구(중랑구 0.630, 강북구 0.562)가 서울시 평균을 상회하여, 물리적 네트워크 인프라 자체의 부재보다는 이를 활용하기 위한 인적\textperiodcentered{}경제적 역량의 부족이 핵심 격차 요인일 가능성을 보여준다. 이러한 발견은 \ref{subsec:3-2}절의 다층 분석에서 개인 수준 변수와 자치구 수준 변수의 상대적 기여를 분리 및 검증함으로써 보다 체계적으로 검토된다. 또한, UMC의 일부 차원에서 지리적으로 인접한 지역이 서로 유사한 패턴을 보였다. 그림~\ref{fig:3-lisa-grid}은 차원별 지역의 공간적 자기상관성을 나타낸 것이다. 공간적으로 군집해 있는 자치구가 서로 유사한 패턴을 보이는 점은 디지털 사막 해소를 위한 광역적 거버넌스 체계 구축의 가능성을 시사한다.

\begin{figure}[htbp]
  \centering
  \begin{subfigure}[t]{0.55\linewidth}
    \centering
    \includegraphics[width=\linewidth]{image9.png}
    \caption{자치구별 UMC 종합점수}
    \label{fig:3-umc-comp}
  \end{subfigure}\hfill
  \begin{subfigure}[t]{0.43\linewidth}
    \centering
    \includegraphics[width=\linewidth]{image10.png}
    \caption{UMC 차원별 점수 히트맵}
    \label{fig:3-umc-heatmap}
  \end{subfigure}
  \caption{자치구별 UMC 종합점수 및 차원별 히트맵}
  \label{fig:3-umc-two}
\end{figure}

\figfull{image11.png}{UMC 차원별 점수}{fig:3-umc-dim}

\begin{figure*}[htbp]
  \centering
  \begin{subfigure}[t]{0.32\textwidth}
    \centering
    \includegraphics[width=\linewidth]{image12.jpeg}
    \caption{Connectivity}
  \end{subfigure}\hfill
  \begin{subfigure}[t]{0.32\textwidth}
    \centering
    \includegraphics[width=\linewidth]{image13.jpeg}
    \caption{Available for Use}
  \end{subfigure}\hfill
  \begin{subfigure}[t]{0.32\textwidth}
    \centering
    \includegraphics[width=\linewidth]{image14.jpeg}
    \caption{Affordability}
  \end{subfigure}

  \vspace{4pt}
  \begin{subfigure}[t]{0.32\textwidth}
    \centering
    \includegraphics[width=\linewidth]{image15.jpeg}
    \caption{Devices}
  \end{subfigure}\hfill
  \begin{subfigure}[t]{0.32\textwidth}
    \centering
    \includegraphics[width=\linewidth]{image16.jpeg}
    \caption{Skills}
  \end{subfigure}\hfill
  \begin{subfigure}[t]{0.32\textwidth}
    \centering
    \includegraphics[width=\linewidth]{image17.jpeg}
    \caption{Safety}
  \end{subfigure}
  \caption{UMC 수준의 공간적 자기상관 (서울 25개 자치구)}
  \label{fig:3-lisa-grid}
\end{figure*}

\subsection{개인 및 지역 특성과 디지털 연결성}
\label{subsec:3-2}

본 연구의 데이터는 개인(Level 1)이 자치구(Level 2) 내에 군집되어 있는 위계적 구조를 갖는다. 여기서 Level 2는 개인의 거주지에 해당한다. 이러한 구조에서 통상적인 OLS 회귀분석은 관측치의 비독립성 가정을 위반하여, 자치구 수준 예측변수의 표준오차를 과소추정하고 1종 오류를 팽창시킬 수 있다~\citep{RaudenbushBryk2002}. 따라서 분산을 25개 자치구를 기준으로 구 내와 구 간 성분으로 분해하고, 개인 수준과 자치구 수준 변수를 동시에 투입할 수 있는 2수준 위계적 선형 모형(HLM)을 채택하였다.

두 번째 분석에서 활용한 자료인 서울서베이는 2003년부터 서울시에서 2만 개 가구를 대상으로 년마다 실시해온 행정 자료이다. 자료의 수집은 매년 8\,$\sim$\,9월에 비대면과 방문면접을 병행하여 이뤄진다.

\subsubsection{변수 구성}
\label{subsubsec:3-2-1}

먼저 분석의 종속변수인 디지털 활용 점수는 서울서베이의 응답항목 중 디지털 서비스 이용빈도 항목의 8개 항목을 합산한 점수이다. 8개 항목은 생활정보 서비스, 전자상거래 서비스, 금융거래 서비스, 공공 서비스, 커뮤니케이션 서비스, 여가 서비스, 교육학습 서비스, 모바일 결제로 구성된다. 이들은 전화와 문자와 같은 기초적 연결 행위가 아니라, 디지털 기기를 매개로 한 실질적인 추가 활용에 해당하여, 디지털 환경이 발달한 서울에 적합한 지표로 판단하였다. 각 항목은 원래 4점 서열 척도를 0--10 연속 척도로 변환하여 총 0--80점 범위의 연속변수를 구성한다. 이 점수가 높을수록 개인의 디지털 서비스 활용이 다양하고 원하는 서비스를 자유롭게 이용할 수 있음을 의미한다.

본 분석의 독립변수는 개인 특성의 변수와 지역 특성의 변수로 구분된다(표~\ref{tab:2-hlm-vars}). 개인 특성의 변수는 개인의 디지털 활용에 영향을 미칠 수 있는 변수인 성별, 연령, 직업, 학력, 소득수준, 장애여부로 구성된다. 교육수준은 중졸 이하, 고졸, 대졸(기준범주), 대학원 이상의 4범주로, 직업은 관리\textperiodcentered{}전문직, 사무직(기준범주), 생산\textperiodcentered{}기능직, 기타(학생\textperiodcentered{}주부\textperiodcentered{}무직)의 4범주로 구분하였다. 주거유형은 아파트(기준범주), 단독주택, 다세대 및 다가구, 기타의 4범주로 분류하였다. 연령과 소득은 총평균 중심화를 적용하여 절편의 해석 가능성을 높였다. 추가로, 1인가구 더미, 장애 여부 더미, 연도 더미(기준: 2023년)를 포함하였다.

지역 수준(Level 2) 변수는 \ref{subsec:3-1}절에서 구축한 UMC 차원 점수 중 지역의 디지털 환경을 결정하는 외생적 요인의 세 차원 Connectivity, Available for Use, Affordability을 활용하였다. 이는, Devices, Digital Skills, Safety가 서울디지털역량실태조사의 개인 응답을 자치구 단위로 집계한 것을 고려한 것이다. 그로 인한 내생성 문제를 피하기 위해 주 분석에서 세 가지 차원을 제외하여 지역 차원의 보편적인 의미있는 연결성에 더 적합하다고 판단되는 차원을 선택하였다. 통제변수로 자치구 평균소득과 고령화율을 포함하였으며, 모든 Level-2 변수는 앞선 \ref{subsec:3-1}절에서와 같이 Min-Max 정규화를 적용하여 $[0,1]$ 범위로 표준화하였다.

분석은 4단계로 진행된다. 무조건 모형에서 급내상관계수를 산출하여 다층 분석의 정당성을 확인한 뒤, 개인 수준 변수를 투입한 M1, 자치구 수준 변수를 추가한 M2, 교차수준 상호작용을 포함한 M3를 순차적으로 추정한다. 교차수준 상호작용은 UMC 인프라 변수(Connectivity, Available for Use, Affordability)와 취약계층 지표의 곱으로 구성하여, 자치구의 디지털 인프라 수준이 취약계층의 디지털 활용에 차별적으로 작용하는지를 검증한다.

\subsubsection{분석 결과}
\label{subsubsec:3-2-2}

분석에 활용된 표본은 서울서베이 2023\,$\sim$\,2024년 응답자 총 10,000명이며, 25개 자치구에서 평균적으로 400명이 표집되었다. 주요 변수에 대한 기술통계량을 나타낸 것은 표~\ref{tab:3-descstat}과 같다. 개인의 디지털 연결성의 평균이 80점 만점에 60점으로 나타났는데, 서울시민의 높은 디지털 연결성을 잘 보여준다. 종속변수가 음의 왜도가 존재할 수 있으나 기관 표준 척도와 10000명에 근접한 표본을 고려해 척도 선형화한 원점수대로 분석을 진행하였다.

표~\ref{tab:4-hlm-results}은 순차적으로 다층 모형의 추정 결과를 나타낸 것이다. 무조건 모형의 급내상관계수는 0.49\%로 자치구 간 분산 비율 자체는 크지 않았으나, 자치구당 평균 표본 크기 400명을 감안한 설계효과가 2.94로 임계값 2.0을 초과하여 다수준 모형의 사용에 대한 타당성을 확인하였다~\citep{HoxEtAl2018}. 그러나, 25개의 자치구의 소규모 표본 / 자치구가 설명하는 분산이 매우 적은 점을 종합하면 개인의 디지털 연결성에 지역의 연결 관련 특성이 미치는 영향이 극히 적을 것으로 추정된다. 이러한 결과는 서울의 우수한 디지털 연결성과 높은 수준의 디지털 역량으로 지역 수준의 연결성 인프라가 영향을 미치지 못한 것일 수 있다. 또한, 서울서베이는 개인의 거주 지역을 조사하기 때문에, 실제로 연결성에 직접적으로 영향을 받는 지역과 일치하지 않을 수 있다.

개인 특성 중 교육 수준의 관계가 가장 두드러졌다. 중졸 이하 학력은 대졸 대비 디지털 활용 점수를 약 22.6점 낮게 나타났다. 이는 80점 만점 기준 28\%에 해당하는 상당한 격차다. 고졸 역시 대졸보다 7.3점 낮게 나타난 반면, 대학원 이상은 오히려 2.9점 높아 교육 수준과 디지털 활용 능력 간의 뚜렷한 정적 관계를 확인할 수 있다. 연령 또한 유의한 영향을 미쳤다. 평균 연령으로부터 1세 증가할수록 개인의 디지털 연결성에 관한 점수가 0.35점 감소하였으며, 65세 이상 고령층의 경우 연령 효과와는 별개로 추가적인 6.0점 하락이 확인되었다. 65세 이상의 연령층에서 디지털 연결성의 두드러진 저하가 나타났다. 직업 측면에서는 관리직\textperiodcentered{}전문직 종사자가 사무직 대비 2.3점 높게 나타나, 고소득 직종일수록 디지털 활용 수준도 높은 경향을 보였다. 주거환경의 경우 단독주택 거주자는 아파트 대비 2.5점, 다세대주택 거주자는 1.8점 낮았는데, 이는

\begin{table*}[htbp]
  \centering\small
  \caption{분석 변수 설명}
  \label{tab:2-hlm-vars}
  \begin{tabularx}{\textwidth}{l Y c l l}
    \toprule
    \textbf{Variable} & \textbf{Description} & \textbf{Level} & \textbf{Scale} & \textbf{Source} \\
    \midrule
    digital\_use\_score & Digital service utilization score (8 items summed) & L1 & 0--80 (continuous) & Seoul Survey 2023--24 \\
    \addlinespace[2pt]
    age\_c & Age (grand-mean centered, M = 49.3) & L1 & years & Seoul Survey \\
    \addlinespace[2pt]
    female & Female (ref: male) & L1 & 0/1 & Seoul Survey \\
    \addlinespace[2pt]
    Edu & Education: $\leq$middle / high / grad+ (ref: college) & L1 & 0/1 each & Seoul Survey \\
    \addlinespace[2pt]
    Income & Monthly income (grand-mean centered) & L1 & 10K KRW & Seoul Survey \\
    \addlinespace[2pt]
    job & Job: professional / blue-collar / other (ref: white-collar) & L1 & 0/1 each & Seoul Survey \\
    \addlinespace[2pt]
    disabled\_bin & Disability status & L1 & 0/1 & Seoul Survey \\
    \addlinespace[2pt]
    hh\_size\_1 & Single-person household & L1 & 0/1 & Seoul Survey \\
    \addlinespace[2pt]
    housing\_detached / \_multi / \_other & Housing: detached / multi-family / other (ref: apt) & L1 & 0/1 each & Seoul Survey \\
    \addlinespace[2pt]
    year & Year 2024 dummy (ref: 2023) & L1 & 0/1 & Seoul Survey \\
    \addlinespace[2pt]
    elderly & Elderly ($\geq$65) & L1 & 0/1 & Seoul Survey \\
    \addlinespace[2pt]
    low\_edu & Low education ($\leq$middle school) & L1 & 0/1 & Seoul Survey \\
    \addlinespace[2pt]
    connectivity & Connectivity dimension score (normalized) & L2 & 0--1 & UMC Part 1 \\
    \addlinespace[2pt]
    wifi\_total\_n & Public WiFi density per 1K pop (normalized) & L2 & 0--1 & Local Admin License \\
    \bottomrule
  \end{tabularx}
\end{table*}

주택 유형에 따른 소득 수준의 차이뿐 아니라 건물 유형별 인터넷 연결 여건의 차이도 반영된 결과로 해석된다. 1인 가구 역시 다인 가구 대비 1.8점 낮은 점수를 보였다. 그 외 개인 특성으로는, 여성이 남성 대비 2.6점 높게 나타났다. 여성이 남성에 비해 더 높은 디지털 연결성을 갖는 것은, 국내 다양한 디지털 관련 서비스에 여성이 더 민감하게 반응한 결과일 수 있다. \textbf{[인용 필요]} 이러한 결과는 국가에 따른 사회문화적인 맥락을 고려할 필요가 있다. 장애인의 경우 3.1점 높게 나타났는데, 표본의 특성에 기인한 것으로 보인다.

자치구 수준의 UMC 인프라 차원은 주효과만으로는 대체로 유의하지 않았다. Model 2에서 세 차원 모두 통계적 유의성에 도달하지 못하였고, 우도비 검정 결과도 Model 1 대비 유의한 개선이 없었다. 그러나 교차수준 상호작용을 투입한 Model 3에서는 Infrastructure 차원이 유의한 부정적 효과를 보였다. 이는 인프라의 평균적 효과가 아닌, 집단별 차별적 효과가 존재함을 의미한다.

본 분석의 핵심 발견은 교차수준 상호작용에서 나타났다. 공공 WiFi 밀도와 고령 더미의 상호작용은 3.03점의 유의한 정적 효과를 보였다. 이는 공공 WiFi 인프라가 비고령자보다 고령자에게 상대적으로 더 큰 혜택을 제공함을 뜻한다. 디지털 인프라 수준이 높은 자치구에 거주하는 저학력 주민은 인프라 수준이 낮은 곳의 저학력 주민 대비 현저히 높은 디지털 활용 점수를 기록하였다. 이 결과는 인프라 투자가 교육 취약계층의 디지털 격차를 완화하는 데 불균형적으로 큰 효과를 가질 수 있음을 시사한다.

개인 수준 상호작용에서는 이중 취약성이 확인되었다. 독거 고령자는 비독거 고령자 대비 추가적으로 4.0점이 낮았으며, 고령일수록 저학력의 부정적인 효과가 더 커지는 것으로 나타났다. 이러한 결과는 단일 취약 요인보다 복합 취약 요인이 디지털 격차를 심화시킴을 보여준다.

\begin{table*}[htbp]
  \centering\small
  \caption{주요 변수의 기술통계량}
  \label{tab:3-descstat}
  \begin{tabularx}{\textwidth}{Y r r r r r}
    \toprule
    \textbf{변수(Variable)} & \textbf{N} & \textbf{평균(Mean)} & \textbf{표준편차(SD)} & \textbf{최소(Min)} & \textbf{최대(Max)} \\
    \midrule
    디지털 활용 점수(Digital use score) & 10,000 & 60.88 & 16.84 & 0 & 80 \\
    연령(Age) & 10,000 & 49.32 & 16.23 & 15 & 99 \\
    여성 비율(Female, proportion) & 10,000 & 0.51 & --- & 0 & 1 \\
    학력: 중졸 이하(Education: $\leq$Middle school) & 10,000 & 0.07 & --- & 0 & 1 \\
    학력: 고졸(Education: High school) & 10,000 & 0.25 & --- & 0 & 1 \\
    학력: 대학원 이상(Education: Graduate+) & 10,000 & 0.13 & --- & 0 & 1 \\
    장애 여부(Disability) & 10,000 & 0.03 & --- & 0 & 1 \\
    월평균 소득(만원 단위, Monthly income [10K KRW]) & 10,000 & 237.55 & 209.73 & 0 & 1,050 \\
    1인 가구(Single-person household) & 10,000 & 0.17 & --- & 0 & 1 \\
    고령(65세 이상, Elderly 65+) & 10,000 & 0.22 & --- & 0 & 1 \\
    (군집) 연결성 지수(Connectivity, norm.) & 50 & 0.5 & 0.18 & 0.18 & 0.98 \\
    이용 가능성(Available for Use, norm.) & 50 & 0.5 & 0.19 & 0.12 & 0.96 \\
    비용 접근성(Affordability, norm.) & 50 & 0.5 & 0.22 & 0.07 & 1 \\
    \bottomrule
    \multicolumn{6}{p{0.97\textwidth}}{\footnotesize\textit{Note.} N = 10,000명(2023\,$\sim$\,2024년, 25개 자치구, 표본 합산). 소득은 만원 단위. 자치구(군집) 단위 변수는 0\,$\sim$\,1(min-max) 정규화 지수임. N = 10,000 individuals nested within $j$ = 25 districts (2023--2024 pooled). Income measured in 10,000 KRW units. District-level variables are min-max normalized (0--1 scale).}
  \end{tabularx}
\end{table*}

\begin{table*}[htbp]
  \centering\small
  \caption{HLM 추정 결과}
  \label{tab:4-hlm-results}
  \begin{tabularx}{\textwidth}{l Y Y Y}
    \toprule
    \textbf{변수} & \textbf{Model 1} & \textbf{Model 2} & \textbf{Model 3} \\
    \midrule
    \multicolumn{4}{l}{\textit{Level 1: 개인 특성}} \\
    \addlinespace[2pt]
    연령(중심화) & $-0.347^{***}$ (0.012) & $-0.347^{***}$ (0.012) & $-0.278^{***}$ (0.014) \\
    여성 & $2.636^{***}$ (0.305) & $2.634^{***}$ (0.305) & $2.601^{***}$ (0.305) \\
    학력: 중학교 이하 & $-22.618^{***}$ (0.770) & $-22.610^{***}$ (0.770) & $-19.019^{***}$ (2.379) \\
    학력: 고등학교 & $-7.289^{***}$ (0.380) & $-7.282^{***}$ (0.380) & $-7.258^{***}$ (0.380) \\
    학력: 대학원 이상 & $2.906^{***}$ (0.483) & $2.909^{***}$ (0.483) & $2.912^{***}$ (0.481) \\
    소득(중심화) & $0.003^{**}$ (0.001) & $0.003^{**}$ (0.001) & $0.002^{**}$ (0.001) \\
    직업: 전문직 & $2.329^{***}$ (0.451) & $2.326^{***}$ (0.451) & $2.28^{***}$ (0.449) \\
    직업: 생산직 & $-0.373$ (0.424) & $-0.379$ (0.424) & $-0.385$ (0.423) \\
    직업: 기타 & $-1.233^{**}$ (0.456) & $-1.226^{**}$ (0.456) & $-1.105^{*}$ (0.454) \\
    장애 & $3.084^{***}$ (0.888) & $3.081^{***}$ (0.888) & $3.006^{***}$ (0.884) \\
    1인 가구 & $-1.775^{***}$ (0.338) & $-1.776^{***}$ (0.338) & $-0.877^{*}$ (0.367) \\
    단독주택 & $-2.478^{***}$ (0.377) & $-2.480^{***}$ (0.377) & $-2.574^{***}$ (0.379) \\
    다세대주택 & $-1.798^{***}$ (0.408) & $-1.795^{***}$ (0.408) & $-1.933^{***}$ (0.408) \\
    연도(2024) & $0.294$ (0.283) & $0.292$ (0.283) & $0.292$ (0.283) \\
    고령자(65세 이상) & $-6.046^{***}$ (1.696) & & \\
    \addlinespace[4pt]
    \multicolumn{4}{l}{\textit{Level 2: 자치구 특성}} \\
    \addlinespace[2pt]
    연결성 & $-1.127$ (0.747) & $-1.688^{*}$ (0.804) & \\
    이용 가능성(WiFi) & $0.085$ (0.780) & $-0.200$ (0.829) & \\
    비용 접근성 & $-1.578$ (1.135) & & \\
    자치구 소득 & $-2.035^{\dagger}$ (1.181) & & \\
    고령화율 & $-0.298$ (1.095) & $-0.293$ (1.111) & \\
    \addlinespace[4pt]
    \multicolumn{4}{l}{\textit{상호작용}} \\
    \addlinespace[2pt]
    고령자 $\times$ 1인 가구 & $-3.958^{***}$ (0.936) & & \\
    연령 $\times$ 저학력 & $-0.277^{***}$ (0.080) & & \\
    \addlinespace[4pt]
    \multicolumn{4}{l}{\textit{교차수준 상호작용}} \\
    \addlinespace[2pt]
    연결성 $\times$ 고령자 & $-1.920$ (1.901) & & \\
    WiFi $\times$ 고령자 & $3.026^{*}$ (1.515) & & \\
    자치구 소득 $\times$ 고령자 & $2.358$ (1.532) & & \\
    연결성 $\times$ 저학력 & $10.936^{***}$ (3.370) & & \\
    \bottomrule
    \multicolumn{4}{p{0.97\textwidth}}{\footnotesize\textit{Note.} Unstandardized coefficients with standard errors in parentheses. Model 0 = null model (intercept only). Reference categories: male, college education, white-collar job, apartment, multi-person household. $^{\dagger}p<.10$, $^{*}p<.05$, $^{**}p<.01$, $^{***}p<.001$.}
  \end{tabularx}
\end{table*}

\subsection{플랫폼 텍스트 분석: LLM 기반 디지털 연결성 사례분석}
\label{subsec:3-3}

앞선 \ref{subsec:3-1}절과 \ref{subsec:3-2}절에서는 행정 데이터와 설문 조사를 이용해 서울의 디지털 격차를 정량적으로 측정하고, 개인 특성과 지역 인프라의 역할을 분리하였다. 그러나 정량적 지표만으로는 주민들이 일상에서 언제, 어떤 맥락에서 디지털 연결의 결여를 경험하는지까지 포착하기 어렵다. 디지털 연결성에 대한 생활 세계의 체감과 구조적 지표 사이의 거리를 좁히기 위해 정성적 증거를 통합할 필요가 있다.

본 절의 목적은 지역 기반 플랫폼의 비정형 텍스트를 이용해 UMC의 6개 차원(Connectivity, Availability for Use, Affordability, Devices, Skills, Safety)에 잠재된 디지털 결여 사례를 체계적으로 추출하고, 이러한 사례를 \ref{subsec:3-1}절과 \ref{subsec:3-2}절의 정량 분석과 결합하는 방법론적 틀을 제시하는 데 있다. 플랫폼 자체를 분석의 대상으로 삼기보다는, 플랫폼이 제공하는 지역 단위의 생활 경험 데이터를 통해 구조와 경험 사이의 연결을 복원하고자 한다.

\subsubsection{당근 플랫폼 특성}
\label{subsubsec:3-3-1}

비정형 텍스트의 출처는 한국의 지역 기반 커뮤니티 플랫폼 '당근'의 동네생활 게시글이다. 당근을 선정한 이유는 세 가지다. 첫째, 각 게시글에 행정동과 자치구 수준의 지리적 메타데이터가 자동으로 부여되어 \ref{subsec:3-1}절에서 구축한 UMC 지수와 직접적으로 연동할 수 있다. 둘째, 플랫폼 이용자가 서울의 25개 자치구 전역에 골고루 분포해 있어 특정 지역이나 연령층에 치우치지 않는다. 셋째, 게시글은 일상생활 속 디지털 문제를 자발적으로 기록하기 때문에 조사자가 설문에서 미처 포착하지 못한 생활상의 결여를 드러낼 수 있다. 연구의 초점은 플랫폼 자체에 대한 평가가 아니라, 이와 같은 구조적 장점을 활용해 생활 세계의 경험을 수집하는 데 있다.

\subsubsection{왜 LLM인가: 기존 방법론의 한계와 LLM의 구조적 이점}
\label{subsubsec:3-3-2}

UMC의 여섯 차원에 따라 비정형 텍스트를 분류하려면 두 가지 요구를 충족해야 한다. 첫째, 게시글이 디지털 연결성 문제와 관련된 지 여부와 관련된 경우 해당 차원을 정확히 판별해야 한다. 둘째, 일상 언어로 쓰인 구어체 한국어는 조사나 어미 변화가 복잡한 교착어적 특성으로 인해 단순한 형태소 분석기만으로는 정확히 처리하기 어렵다. 기존의 토픽 모델링이나 키워드 기반 분류는 잠재 주제를 파악하는 데 유용하지만 사전 정의된 차원과 일치하지 않거나, 대화체 텍스트를 제대로 분석하지 못하는 한계가 있다. 이러한 한계는 한국어 형태소 분석기가 구어체 데이터에서 정확도가 낮고, 데이터 훈련이 필요한 머신러닝 모델이 충분한 라벨을 요구한다는 점에서 더욱 두드러진다.

대규모 언어 모델(LLM)은 이런 제약을 상당 부분 극복한다. 다국어로 사전 학습되어 한국어의 형태소 분석이 없어도 맥락을 파악할 수 있으며, 104개 이상의 언어를 동시에 다루는 모델은 언어간 의미론적 표현을 공유해 적은 양의 예시만으로도 새로운 태스크에 적응하는 능력을 갖추고 있다. LLM을 도구로 사용할 때는 출력의 확률성과 편차를 고려해야 한다. 동일한 입력에 대해서도 결과가 달라질 수 있으므로, 본 연구는 분석 절차를 고정하고 LLM은 분류와 추론을 수행하는 계산 모듈로만 활용했다. 이를 통해 연구자가 정의한 분류 규칙과 추론 프로토콜을 일관되게 적용하면서도 LLM의 문맥 이해력을 활용할 수 있도록 설계하였다.

\subsubsection{분석 구조}
\label{subsubsec:3-3-3}

텍스트 분석은 네 단계로 구성된다. 먼저, 크롤링한 원시 CSV 데이터에서 삭제되어 접근할 수 없는 게시글과 광고성\textperiodcentered{}중복 게시글을 제거한 뒤, 여섯 차원별 키워드 사전을 이용해 UMC 후보 게시글을 추출한다. 키워드 사전은 디지털 연결성과 직간접적으로 관련된 단어를 정리한 것으로 해당 단어를 포함하는 게시글만 1차로 추출하였다. 이후, 게시글이 디지털 연결성 문제와 관련이 있는지 여부를 판정하는 1차 분류 (umc\_related = Y/N/?)를 수행하며, 판정의 범위를 넓히고자 LLM에게 예시, 제한을 상세하게 밝히지 않고 분석하였다. 그 다음, 관련성이 확인된 게시글만을 대상으로 사용자가 제공한 상세한 예시를 포함한 few-shot 프롬프트를 LLM에게 제시해 게시글을 6개 개념적 차원의 사례로 분류한다. 이 단계에서는 모델의 자율성을 제한하여 프롬프트에 예시를 포함한 구체적인 분류 기준과 명세만을 제공했고, 모델이 판단에 참조하는 정보 공간인 컨텍스트를 엄격하게 관리하였다.

표~\ref{tab:5-text-sample}는 분석 구조와 표본 추출 과정을 요약한 것이다. 원시 데이터 1,287,761건 중 약 51.5\%는 게시자 탈퇴나 차단으로 본문을 볼 수 없어 분석에서 제외되었다. 중복\textperiodcentered{}광고성 게시글 제거 후에는 약 48만 건이 남았으며, 여섯 차원별 키워드 사전을 적용한 결과 26,742건이 UMC 관련 후보로 추출되었다. few-shot 프롬프트를 이용한 1차 판정에서는 약 57.7\%(15,438건)가 관련(Y)으로, 29.6\%(7,921건)가 비관련(N)으로, 12.6\%(3,383건)가 판단 유보(?)로 분류되었다. 이 가운데 Y로 판정된 게시글만을 대상으로 차원 분류를 진행했다.

차원 분류 판정에서 LLM에게 각 차원에서 대표적인 사례를 담은 3\,$\sim$\,4개의 예시를 프롬프트로 제공했다. 이러한 few-shot 프롬프트는 모델이 사람의 판단 기준을 학습하는 데 도움이 되며, 복잡한 작업에서 zero-shot보다 나은 성능을 보인다~\textbf{[인용 필요]}. 예시에는 어떤 유형의 게시글이 디지털 연결성 문제와 무관한지를 보여주는 음성 예시도 포함해, 관련성을 최대한 좁혀 높은 정밀도를 확보하도록 했다. 2차 차원 분류에서는 구체적인 분류 기준과 정의만을 프롬프트로 제공해 모델이 보다 자율적으로 사례를 분류하도록 설계했다. 이처럼 앞 단계는 예시 기반 few-shot 방식이고, 뒤 단계는 규칙 기반 프롬프트를 사용한다는 점에서 추론과 few-shot 사이의 균형을 보여준다(그림~\ref{fig:3-pipeline-roles} 참조).

\begin{table}[htbp]
  \centering\small
  \caption{텍스트 분석 표본 추출 및 판정 절차}
  \label{tab:5-text-sample}
  \begin{tabularx}{\linewidth}{l Y r}
    \toprule
    \textbf{단계} & \textbf{설명} & \textbf{남은 게시글 수} \\
    \midrule
    원시 데이터 & 2024년 5월\,$\sim$\,2026년 3월에 작성된 동네생활 게시글 전수 & 1,287,761 \\
    접근가능 게시글 & 탈퇴\textperiodcentered{}차단으로 본문 확인이 불가능한 게시글 제거 & 624,564 \\
    중복\textperiodcentered{}광고 제거 & 중복된 내용, 광고\textperiodcentered{}홍보 게시글 제거 & 480,972 \\
    UMC 키워드 필터 & 6개 차원별 키워드 사전을 적용해 후보 추출 & 26,742 \\
    1차 판정 & few-shot 프롬프트로 UMC 관련성 분류 & --- \\
    \bottomrule
  \end{tabularx}
\end{table}

\fig{image18.png}{분류와 추론의 역할 구분}{fig:3-pipeline-roles}

이후, 분류 결과를 자치구 단위로 집계하여 표출 프로파일을 만들고, \ref{subsec:3-1}절의 정량 지표와 결합하기 위해 베이지안 업데이트를 적용한다. 정량 지표와 텍스트 분류 결과를 결합하기 위해 베이지안 업데이트를 적용하였다. 사전 분포는 UMC의 6가지 개념의 잠재적 공간이다. 차원 $d$에서 자치구 $j$의 UMC 실측값을 $s_{jd}$로 할 때, 사전 분포는 $\mathrm{Beta}(\alpha_{jd}, \beta_{jd})$이며, 평균이 $s_{jd}$가 되도록 사전 분포를 가정한다. 강도 모수 $\kappa$를 사용해 $\alpha_{jd} = s_{jd}\cdot\kappa$, $\beta_{jd} = (1-s_{jd})\cdot\kappa$로 표현한다. 텍스트 분류에서 해당 차원으로 판정된 게시글 수를 $y_{jd}$, 총 분류된 게시글 수를 $n_{jd}$라 할 때, 우도는 $\mathrm{Binomial}(n_{jd}, \theta_{jd})$로 표현된다. 사후 분포는 $\theta_{jd}\mid y_{jd} \sim \mathrm{Beta}(\alpha_{jd}+y_{jd}, \beta_{jd}+n_{jd}-y_{jd})$로 계산된다. 사전 강도 $\kappa$는 기본값 20을 사용하고 10과 30에 대한 민감도 분석도 수행하였다. 게시글 수가 적은 자치구에서 사후 분포가 사전 분포에 의해 지나치게 지배되는 것을 막기 위해, 게시글 수 $n_{jd}$를 생활인구로 보정한 $n_{jd}^{*} = n_{jd}\times(\text{생활인구}_{j}/\overline{\text{생활인구}})$로 사용하였다.

\begin{align}
\alpha_{jd} &= \kappa \times s_{jd} \label{eq:alpha} \\
\beta_{jd}  &= \kappa \times (1 - s_{jd}) \label{eq:beta} \\
\theta_{jd}\mid y_{jd} &\sim \mathrm{Beta}\bigl(\alpha_{jd}+y_{jd},\, \beta_{jd}+n_{jd}-y_{jd}\bigr) \label{eq:posterior} \\
n_{jd}^{*} &= n_{jd}\times\bigl(\text{생활인구}_{j}/\overline{\text{생활인구}}\bigr) \label{eq:popadj}
\end{align}

베이지안 집계에서 게시글로 추정된 사후 분포가 사전 분포로 유의미하게 이동한 자치구\textperiodcentered{}차원 조합을 대상으로 인과기제에 관한 정성적 추론을 수행하였다. 구체적으로는 (1) 사후 이동이 통계적으로 유의하며, (2) 해당 자치구\textperiodcentered{}차원에서 Y 판정 게시글이 최소 5건 이상이고, (3) \ref{subsec:3-2}절의 다층모형 결과가 허용하는 범위 내에서만 해석을 실시하였다. 사전 분포에 의해 도출된 조합에 대해 대표 게시글 3\,$\sim$\,5건을 선정하고, 1\,$\sim$\,2문장으로 구조적 지표와 생활 경험 사이의 연결 경로를 서술하였다. 분류와 추론 작업은 세 명의 독립된 에이전트가 서로 다른 추론 경로를 통해 수행했으며, 최종 판정은 네 번째 에이전트가 세 결과를 비교\textperiodcentered{}검증하는 방식으로 이루어졌다. 이는 한 가지 프로토콜에 의존하지 않고 여러 관점의 상호 검증을 통해 편향을 통제하는 삼각검증(triangulation) 접근을 따른 것이다.

모든 분류와 추론 작업은 세 명의 독립된 분석 에이전트가 서로 다른 추론 방식 프롬프트와 컨텍스트를 사용해 수행하고, 네 번째 분석자가 세 결과를 비교\textperiodcentered{}검증하는 방식으로 이루어졌다. 이는 질적 연구에서 여러 자료원과 분석가를 동원해 편향을 줄이는 삼각검증(triangulation) 접근에 기반한다. 세 관점이 일치할 때만 해당 게시글을 최종 분석에 포함했으며, 해석상의 불확실성은 결과 표기에 반영하였다. 이러한 과정을 정리한 결과는 표~\ref{tab:6-units}, 그림~\ref{fig:3-pipeline-full}과 같다.

\begin{table}[htbp]
  \centering\small
  \caption{\ref{subsec:3-3}절의 분석 단위와 해석 단위}
  \label{tab:6-units}
  \begin{tabularx}{\linewidth}{l l Y Y}
    \toprule
    \textbf{구분} & \textbf{단위} & \textbf{무엇을 보여주는가} & \textbf{해석상 주의점} \\
    \midrule
    관찰 단위 & 개별 게시글 & 디지털 연결성 경험의 단서 & 단일 게시글만으로 지역 구조를 주장할 수 없음 \\
    패턴 단위 & 행정동 게시글 & 차원별 표출 비율과 분포 & 집계 기준과 표본 편향을 함께 검토해야 함 \\
    구조 단위 & UMC 지역 지표 & 지역 수준의 연결성 조건 & 텍스트 집계 결과와의 정합을 통해 제한적으로 해석 \\
    인과기제 단위 & 추론 메커니즘 & 사례와 지역 패턴의 연결 설명 & \ref{subsec:3-2}의 정량 결과가 허용하는 범위 안에서만 해석 \\
    \bottomrule
  \end{tabularx}
\end{table}

\subsubsection{데이터와 적용 범위}
\label{subsubsec:3-3-4}

원시 데이터는 서울 25개 자치구를 대상으로 2024년 5월부터 2026년 3월까지 작성된 동네생활 게시글 전수를 크롤링한 것으로, 중복 제거 후 총 1,287,761건이다. 이 중 삭제(DELETED) 663,163건(51.5\%)과 차단(BLOCKED) 103,728건(8.1\%)을 합한 766,891건(59.6\%)은 본문 확인이 불가능하여 분석에서 제외하였다. 분석 대상은 최종 약 52만 건이다.

표본 탈락의 체계적 편향 가능성을 검토한 결과, 자치구별 제거율의 표준편차는 4.6\%로 자치구 간 편차가 크지 않았으며, 디지털 결여 추론의 주요 대상인 생활/편의(47.0\%), 동네사건사고(48.4\%), 병원/약국(35.5\%) 카테고리는 사교\textperiodcentered{}홍보 목적 카테고리(동네친구 86.6\%, 교육 69.0\%)에 비해 제거율이 낮았다. 삭제의 상당 부분은 사용자 탈퇴에 수반된 일괄 삭제에 기인하는 것으로 판단된다. 다만, 삭제 게시글 중 일부가 디지털 결여 관련 사례를 포함할 가능성을 완전히 배제할 수는 없고, 삭제 행위 자체가 디지털 결여와 체계적으로 연관되었을 수 있다. 이는 본 분석의 한계에 해당한다.

당근 데이터의 구조적 한계도 인식할 필요가 있다. 플랫폼 이용자 자체가 디지털 접근이 가능한 집단으로 한정되므로, 디지털 소외 계층의 경험은 구조적으로 과소 대표된다. 게시글의 시간적\textperiodcentered{}공간적 분포가 균일하지 않아 자치구 간 비교에서 표본 편향의 가능성이 존재한다. 이러한 한계는 데이터 소스의 제약이지 방법론의 제약이 아니다. 보다 직접적인 질적 자료가 존재하는 맥락에서는 동일한 파이프라인을 적용하여 보다 정밀한 분석이 가능할 것이다.

\figfull{image19.png}{LLM 기반 텍스트 분석 파이프라인}{fig:3-pipeline-full}

\subsubsection{분석 결과}
\label{subsubsec:3-3-5}

텍스트 분류 결과는 표~\ref{tab:6-units}와 같다. 연결 품질(Connection Quality) 차원은 6개 차원 가운데 전이의 폭이 가장 작다($-0.006 \sim +0.005$). 실측과 텍스트 증거가 대체로 수렴하는 차원으로 볼 수 있다. 양천구는 실측 UMC 0.278로 낮은 편이나 관련 게시글 비율이 상대적으로 높아 $+0.005$의 전이를 보였고, 용산구는 실측 0.884로 높은 반면 관련 게시글이 드물어 $-0.006$의 전이를 보였다.

이용 가능성(Availability for Use) 차원에서는 도봉구가 실측 최하위(0.213)에 해당하면서도 관련 게시글 비율이 높아 $+0.006$의 전이를 보였다. 광진구는 실측 0.832로 높고 관련 게시글이 드물어 $-0.010$의 전이를 기록하였다.

경제적 접근성(Affordability) 차원은 서울특별시 자치구의 소득 수준과 일관된 방향의 경향을 보였다. 중랑구와 강서구는 각각 $+0.005$의 양의 전이를 보인 반면, 서초구는 실측이 0.993으로 최고 수준이면서 관련 게시글이 드물어 $-0.007$의 전이를 나타냈다. 이 차원은 통신비와 기기 구입 비용에 관한 고민이 소득 수준에 따라 차등적으로 표출되는 경향과 정합적이다.

기기 접근성(Devices) 차원은 자치구 간 편차가 가장 큰 차원이다(gap = 0.090). 중랑구는 실측 0.052에 해당하면서도 기기 관련 게시글 비율이 4.95\%로 전 자치구 가운데 가장 높아 $+0.048$의 전이를 보였다. 송파구 또한 $+0.026$의 유사한 패턴을 나타냈다. 반면 용산구($-0.042$)와 동작구($-0.033$)는 실측이 높으면서 관련 게시글이 드물게 나타났다. 다만 이 차원에서 실측 지표는 기기 보유율을 포착하는 반면, 텍스트 증거는 중고 거래, 수리 요청 등 기기와 관련한 일상적 거래와 고충을 포착하므로, 두 지표가 지시하는 대상이 부분적으로 상이할 가능성을 배제하기 어렵다.

디지털 역량(Digital Skills) 차원은 기기 차원에 이어 두 번째로 편차가 큰 차원이다(gap = 0.092). 중랑구는 실측 0.242에 관련 게시글 비율 4.8\%로 $+0.047$의 전이를 보였다. 강남구 또한 실측이 0.276으로 낮은 편이면서 관련 게시글 비율이 상대적으로 높아 $+0.013$의 전이를 나타냈다. 반면 관악구는 실측 0.795로 최고 수준이면서 관련 게시글이 드물어 $-0.045$의 전이를 보였다.

디지털 역량(Digital Skills) 차원에서는 기기 차원에 이어 두 번째로 큰 편차가 나타났다. 중랑구는 실측 0.242에 비해 관련 게시글 비율 4.8\%로 $+0.047$의 이동을 보였으며, 강서구는 실측값이 높음에도 관련 게시글이 적어 $-0.031$의 이동을 보였다. 안전\textperiodcentered{}보안(Safety) 차원은 피싱 및 보이스피싱 사례가 게시글의 주요 내용을 이루었고, 중랑구($+0.018$)와 도봉구($+0.015$)에서 양의 이동이 확인되었다. 종로구는 실측값이 매우 높고 관련 게시글이 적어 사후 분포가 낮아졌다.

안전\textperiodcentered{}보안(Safety \& Security) 차원에서는 중랑구($+0.018$)와 도봉구($+0.015$)에서 피싱 사기 및 보이스피싱 관련 게시글의 비율이 높았다. 종로구는 실측 0.984로 최고 수준이면서 관련 게시글이 드물어 $-0.019$의 전이를 나타냈는데, 관광\textperiodcentered{}상업 중심지의 주거 구성과 당근 이용자 집단의 특성이 상이할 가능성을 함께 고려할 필요가 있다.

자치구 단위로 집계했을 때, 지역의 실측값과 게시글 분포 간 차이는 크게 세 가지 유형. 첫째, 평균 전이가 양의 방향으로 나타나는 자치구이다. 중랑구는 6개 차원 중 4개 차원에서 양의 전이가 가장 크게 나타나 평균 $+0.020$으로 가장 높은 값을 보였다. 도봉구와 강남구 또한 평균적으로 양의 방향을 보였으나 그 크기는 0.001 수준으로 작았다. 특히 강남구의 경우 평균 소득이 높은 자치구임에도 이용 가능성과 디지털 역량 차원에서 양의 전이가 관찰되었는데, 자치구 내부의 이질적 집단이 관련 이슈를 표출하고 있을 가능성을 시사한다. 둘째, 전이가 0에 가까운 자치구이다. 강북구($-0.000$), 강동구($-0.001$), 동대문구($-0.001$)가 이에 해당한다. 이들 자치구에서는 차원별로 양\textperiodcentered{}음의 전이가 섞여 있어 평균적으로 상쇄되는 양상을 보인다. 셋째, 평균 전이가 음의 방향으로 나타나는 자치구이다. 마포구($-0.014$), 영등포구($-0.013$), 관악구($-0.013$), 서초구($-0.009$)가 이에 해당한다. 관악구는 대학가 중심의 인구 구성과 디지털 역량 실측값이 맞물려 도움 요청 성격의 게시글이 드물게 나타난 것으로 볼 수 있으며, 서초구는 경제적 접근성과 기기 차원에서 특히 큰 음의 전이를 보였다. 사전 강도 $\kappa$를 10, 20, 30으로 달리한 민감도 분석에서도 사후 평균의 차이는 $0.00001\sim0.00010$ 수준으로 실질적 변화가 없었다. 이는 생활인구 보정 후 유효 관측 수가 충분해 사전의 가중치 변화가 결과에 큰 영향을 미치지 않음을 시사한다.

해석에서 주의할 점은 양의 전이가 곧바로 실측 지표가 과소평가되었음을 의미하지 않는다는 것이다. 양의 이동은 실측 지표가 포착하지 못한 생활상의 불편을 반영할 수 있지만, 특정 자치구에서 특정 문제에 대한 논의가 활발한 문화적\textperiodcentered{}인구학적 요인을 동시에 포함할 수 있다~\textbf{[인용 필요]}. 또한 플랫폼 이용자의 표집 편향과 게시글의 시공간적 불균형이 결과 해석에 영향을 미칠 수 있다. 따라서 사후 분포의 방향과 크기를 정책적 함의로 연결할 때는 정량\textperiodcentered{}정성적 지표를 종합적으로 고려해야 한다. 본 절의 결과에서 유의미한 것은 전이의 부호와 크기가 차원과 자치구에 따라 체계적으로 달라진다는 것이다. 단순히 당근마켓 게시글의 빈도만을 집계할 경우 플랫폼 활성도 자체가 UMC 수준에 대한 추정과 혼동될 우려가 있다. 사전 분포로 실측값을 투입하고 텍스트 분류 결과를 우도로 결합하는 절차는, 관측된 텍스트 빈도에서 플랫폼 활성도에 기인하는 부분과 UMC 수준에 기인하는 부분을 부분적으로 분리하는 효과를 갖는다. 이 점에서 사전-사후 분포의 전이는 텍스트 증거가 실측을 얼마나 보정하는지를 보여주는 지표로 기능한다.

다만 shift의 방향을 해석할 때 두 가지 제약을 함께 고려해야 한다. 첫째, 양의 shift가 관찰된다고 해서 해당 자치구의 UMC 수준이 실측보다 낮다고 곧바로 결론지을 수 없다. 양의 shift는 실측 지표가 포착하지 못하는 생활 수준의 체감이 존재할 가능성을 시사하지만, 관련 주제가 당근마켓이라는 특정 플랫폼에서 표출되는 방식 자체가 자치구별 이용자 구성에 의존하기 때문이다. 둘째, 음의 shift 또한 해당 자치구의 UMC 수준이 실측보다 높다는 의미로 해석하기보다, 관련 이슈가 플랫폼 상에서 표출될 유인이 상대적으로 적은 상태로 읽는 것이

\begin{table*}[htbp]
  \centering\small
  \caption{자치구별 게시글 판정 결과}
  \label{tab:7-district-cls}
  \begin{tabularx}{\textwidth}{l r r r r r}
    \toprule
    \textbf{자치구} & \textbf{총 게시글(n)} & \textbf{Y 판정} & \textbf{? 판정} & \textbf{N 판정} & \textbf{Y 비율(\%)} \\
    \midrule
    \textbf{중랑구}   & 5,476  & 794 & 28  & 4,654  & 14.5 \\
    \textbf{송파구}   & 9,367  & 860 & 99  & 8,408  & 9.18 \\
    \textbf{서초구}   & 6,170  & 434 & 216 & 5,520  & 7.03 \\
    \textbf{은평구}   & 6,140  & 429 & 50  & 5,661  & 6.99 \\
    \textbf{도봉구}   & 3,491  & 233 & 5   & 3,253  & 6.67 \\
    \textbf{서대문구} & 4,500  & 289 & 5   & 4,206  & 6.42 \\
    \textbf{강북구}   & 3,885  & 247 & 31  & 3,607  & 6.36 \\
    \textbf{용산구}   & 3,012  & 181 & 12  & 2,819  & 6.01 \\
    \textbf{양천구}   & 4,975  & 296 & 7   & 4,672  & 5.95 \\
    \textbf{강동구}   & 6,254  & 342 & 231 & 5,681  & 5.47 \\
    \textbf{영등포구} & 5,964  & 306 & 5   & 5,653  & 5.13 \\
    \textbf{동대문구} & 4,754  & 238 & 29  & 4,487  & 5.01 \\
    \textbf{성북구}   & 5,437  & 266 & 141 & 5,030  & 4.89 \\
    \textbf{중구}     & 1,351  & 63  & 16  & 1,272  & 4.66 \\
    \textbf{강서구}   & 8,174  & 371 & 236 & 7,567  & 4.54 \\
    \textbf{금천구}   & 2,298  & 101 & 7   & 2,190  & 4.40 \\
    \textbf{종로구}   & 2,267  & 97  & 27  & 2,143  & 4.28 \\
    \textbf{동작구}   & 3,675  & 150 & 22  & 3,503  & 4.08 \\
    \textbf{강남구}   & 10,047 & 404 & 459 & 9,184  & 4.02 \\
    \textbf{성동구}   & 404    & 16  & 9   & 379    & 3.96 \\
    \textbf{관악구}   & 11,665 & 407 & 75  & 11,183 & 3.49 \\
    \textbf{마포구}   & 7,974  & 242 & 254 & 7,478  & 3.03 \\
    \textbf{구로구}   & 4,082  & 121 & 96  & 3,865  & 2.96 \\
    \textbf{광진구}   & 5,440  & 147 & 42  & 5,251  & 2.70 \\
    \textbf{노원구}   & 5,040  & 103 & 46  & 4,891  & 2.04 \\
    \bottomrule
  \end{tabularx}
\end{table*}

\figfull{image20.png}{지역별 사전-사후 분포 차이}{fig:3-prior-posterior}

안전하다. 중랑구와 마포구의 대조에서 드러나듯, shift의 해석은 자치구의 인구 구성과 플랫폼 이용 특성을 비롯한 다양한 맥락 정보와 결합될 때 의미를 띈다. 둘째, 음의 전이 또한 해당 자치구의 UMC 수준이 실측보다 높다는 의미로 해석하기보다, 관련 이슈가 플랫폼 상에서 표출될 유인이 상대적으로 적은 상태로 읽는 것이 타당할 수 있다. 이런 측면에서 지역 자원이 실제로 결여되어 있음에도 플랫폼 상에서 논의되지 않는 이유를 고려하는 것 또한 중요하다.

\begin{table*}[htbp]
  \centering\small
  \caption{결여 유형별 가설 생성 결과 분포}
  \label{tab:8-deficit-types}
  \begin{tabularx}{\textwidth}{l Y r r}
    \toprule
    \textbf{결여 유형} & \textbf{관련 UMC 차원} & \textbf{관련 게시글 수 (n=100)} & \textbf{비율 (\%)} \\
    \midrule
    정보부재         & Digital Skills      & 63 & 63.0 \\
    제도적 경로 결여 & Available for Use   & 61 & 61.0 \\
    기기 접근 결여   & Devices             & 50 & 50.0 \\
    디지털 역량 부재 & Digital Skills      & 45 & 45.0 \\
    연결 인프라 결여 & Connectivity        & 40 & 40.0 \\
    공공 접근점 결여 & Available for Use   & 32 & 32.0 \\
    경제적 접근 결여 & Affordability       & 27 & 27.0 \\
    안전\textperiodcentered{}보안 결여 & Safety & 16 & 16.0 \\
    복합 결여 (2개 이상) & ---             & 88 & 88.0 \\
    \bottomrule
    \multicolumn{4}{p{0.97\textwidth}}{\footnotesize\textit{Note.} 하나의 게시글에서 복수의 결여 유형이 추론될 수 있으므로 비율의 합은 100\%를 초과한다. Phase 1은 가설 생성 단계이며, Phase 2(판정)에서 삼각검증을 거쳐 최종 분류가 확정된다.}
  \end{tabularx}
\end{table*}

\subsection{분석 결과 종합}
\label{subsec:3-4}

본 절에서는 \ref{subsec:3-1}절(UMC 지수 구축), \ref{subsec:3-2}절(HLM 다수준 분석), \ref{subsec:3-3}절(LLM 기반 텍스트 분석)의 결과를 종합하여, 서울의 디지털 격차가 어떤 구조로 작동하는지를 통합적으로 해석한다. 세 분석은 각각 '어디에 격차가 있는가'(측정), '왜 격차가 발생하는가'(인과 규명), '격차가 어떻게 경험되는가'(질적 맥락)에 답하며, 이들의 교차 지점에서 정책적으로 유의미한 함의가 도출된다.

세 분석이 공통으로 지지하는 첫 번째 발견이자 본 프로젝트의 핵심 실증 결과는, 서울의 디지털 격차가 자치구 간 차이가 아닌 개인 간 차이에 의해 구조화되어 있다는 것이다. \ref{subsec:3-2}절의 HLM 분석에서 ICC는 0.49\%에 불과하여, 디지털 활용 점수 분산의 99.5\%가 동일 자치구 내 개인 간 차이에 기인하였다. 이 결과는 \ref{subsec:3-1}절에서 구축한 UMC 종합지수의 자치구 간 편차가 실재함에도 불구하고, 해당 편차가 개인의 디지털 활용 수준에 미치는 직접적 영향은 제한적임을 의미한다. 그러나 이러한 낮은 ICC는 자치구 수준 변수의 무의미함을 뜻하는 것이 아니라, 그 효과가 집단별로 차별적으로 발현됨을 시사한다. Model 3의 교차수준 상호작용 결과가 이를 뒷받침한다.

\ref{subsec:3-3}절의 텍스트 분석은 이 양적 발견이 일상에서 어떻게 발현되는지를 구체화한다. '정보부재'가 가장 빈번한 결여 유형(63\%)으로 나타난 것은, 교육수준이 낮을수록 디지털 서비스의 존재 자체를 인지하지 못하는 정보 격차가 더 심화될 수 있음을 시사한다. 프린터 이용을 위해 편의점 무인 인쇄를 모르는 사례, ISP 고객센터 번호를 모르는 사례, 정부24의 공문서 출력 기능을 모르는 사례 등은 모두 '서비스가 존재하지만 이용자에게 보이지 않는' 구조적 문제를 가리킨다. 이는 HLM에서 교육 기울기가 가파른 이유를 경험적 수준에서 설명하는 보완적 증거로 해석할 수 있다.

두 번째 발견은, 자치구 인프라의 효과가 주효과 수준에서는 유의하지 않으나 특정 취약집단에 대해서는 조절 효과로서 유의하게 작용한다는 것이다. Model 2에서 UMC 차원 중 Connectivity와 관련된 인프라 변수의 주효과는 통계적 유의성에 도달하지 못하였는데($p = 0.237$), 이는 ICC가 0.49\%인 상황에서 자치구 수준 변수의 주효과가 제한적인 것과 일관된 결과이다. 그러나 Model 3에서 교차수준 상호작용이 유의하게 나타남으로써, 인프라의 효과가 평균적 효과가 아닌 집단별 조절 효과로서 작동함이 확인되었다.

\ref{subsec:3-3}절의 텍스트 분석은 이 차별적 효과의 질적 맥락을 제공한다. 공공 접근점(프린트 카페, 도서관 복합기, 무인 인쇄)에 대한 정보부재가 32\%의 게시글에서 추론된 것은, 인프라가 물리적으로 존재하더라도 이용자의 인지와 접근 역량이 수반되지 않으면 실질적 혜택으로 전환되지 않음을 보여준다. 은평구 신사동의 사례('정녕 없는 건지요')는 온라인 검색을 이미 시도했으나 결과가 불충분했을 가능성을 시사하며, 이는 인프라의 존재와 인프라의 '발견가능성(discoverability)' 사이의 간극을 조명한다. HLM에서 인프라 주효과가 유의하지 않았던 결과와 결합하면, 인프라 투자의 효과는 취약계층의 정보 접근과 역량 지원이 병행될 때 비로소 발현된다는 해석이 가능하다.

세 번째 발견은, 디지털 격차가 단일 축이 아닌 복합 취약성의 교차로 심화된다는 것이다. HLM에서 독거 고령자는 비독거 고령자 대비 추가적으로 4.0점이 낮았고, 연령 증가에 따른 저학력의 부정적 효과가 가속화되는 패턴이 확인되었다. \ref{subsec:3-3}절에서도 결여의 복합성이 일관되게 나타났다. 게시글당 평균 2.5개 이상의 결여 유형이 동시에 식별되었으며, 정보부재\textperiodcentered{}역량부재\textperiodcentered{}제도적 경로 차단이 중첩될 때 이용자가 공식 경로에서 이탈하여 커뮤니티에 의존하는 구조가 확인되었다. 디지털 지원망의 부재('자녀분들이 없으시거나 떨어져 살아서 묻기 힘들거나')는 가구 구성이 디지털 역량 획득의 비공식적 경로로 기능함을 보여주며, 독거 고령자의 추가적 격차를 설명하는 기제로 해석된다.

네 번째 발견은, 디지털 사막의 공간적 군집과 차원별 이질성이다. \ref{subsec:3-1}절에서 하위 5개 지역 중 특히 서울 북부(도봉구, 강북구, 노원구, 중랑구)가 UMC 종합지수의 Low-Low 클러스터를 형성하였으며, 이들은 Connectivity 차원에서는 도시 평균 이상을 기록하면서도 Devices와 Safety 차원에서 극단적 결핍을 보였다. 이러한 차원별 프로파일의 이질성은 일률적 인프라 투자보다 차원별 맞춤 개입의 필요성을 시사한다.

방법론적 관점에서, 세 분석의 결합은 각각의 한계를 상호 보완한다. \ref{subsec:3-1}절의 종합지수는 자치구 간 상대적 위치를 파악하되 개인 내 변이를 포착하지 못하며, \ref{subsec:3-2}절의 HLM은 개인과 지역의 기여를 분리하되 디지털 격차의 질적 발현 경로를 설명하지 못한다. \ref{subsec:3-3}절의 텍스트 분석은 결여의 경험적 맥락을 가시화하되 대표성에 한계가 있다. 세 분석을 병렬이 아닌 순차적 삼각검증(측정 $\rightarrow$ 인과 규명 $\rightarrow$ 질적 맥락)으로 결합함으로써, '어디에'(자치구), '왜'(개인 특성과 인프라의 상호작용), '어떻게'(결여의 발현 경로)라는 세 질문에 대한 통합적 답을 구성할 수 있었다.

이상의 종합 결과는 4장의 정책 제언으로 직결된다. 개인 특성이 지배적인 설명변수라는 발견은 people-based 정책의 우선성을, 인프라의 차별적 효과는 취약계층 타겟형 place-based 정책의 필요성을, 복합 취약성의 확인은 단일 프로그램이 아닌 통합 지원 체계의 필요성을, 공간적 군집의 확인은 권역 단위 거버넌스의 근거를 각각 제공한다.
